\documentclass{article}

\usepackage[preprint]{tmlr}
\usepackage{amsmath}
\usepackage{amssymb}
\usepackage{mathtools}
\usepackage{longtable,array}
\usepackage{booktabs}
\usepackage{siunitx}
\sisetup{round-mode=places, round-precision=2, table-number-alignment=center, group-separator={}}
\usepackage{newunicodechar}
\newunicodechar{Σ}{\ensuremath{\Sigma}}
\newunicodechar{σ}{\ensuremath{\sigma}}
\newunicodechar{δ}{\ensuremath{\delta}}
\newunicodechar{α}{\ensuremath{\alpha}}
\newunicodechar{Ω}{\ensuremath{\Omega}}
\newunicodechar{ε}{\ensuremath{\varepsilon}}
\newunicodechar{γ}{\ensuremath{\gamma}}
\newunicodechar{ρ}{\ensuremath{\rho}}
\newunicodechar{λ}{\ensuremath{\lambda}}
\newunicodechar{η}{\ensuremath{\eta}}
\newunicodechar{μ}{\ensuremath{\mu}}
\newunicodechar{κ}{\ensuremath{\kappa}}
\newunicodechar{ν}{\ensuremath{\nu}}
\newunicodechar{φ}{\ensuremath{\phi}}
\newunicodechar{Ψ}{\ensuremath{\Psi}}
\newunicodechar{Ξ}{\ensuremath{\Xi}}
\newunicodechar{Θ}{\ensuremath{\Theta}}

\usepackage{graphicx}
\usepackage{tikz}
\usepackage{pgfplots}
\pgfplotsset{compat=1.18}
\usetikzlibrary{patterns, positioning, arrows.meta, shapes.geometric, backgrounds, fit}

\usepackage{algorithm}
\usepackage{algpseudocode}
\usepackage{enumitem}
\usepackage{amsthm}
\usepackage{microtype}
\usepackage{hyperref}
\hypersetup{
  pdftitle={Critical Compositional Pressure: A Phase-Boundary Framework for Compositional Representation Formation in Neural Networks},
  pdfsubject={Machine Learning, Dynamical Systems, Representation Learning},
  pdfkeywords={compositional generalization, continuous gradient flow, phase transition, shortcut learning, transcritical bifurcation}
}
\usepackage[capitalize,noabbrev]{cleveref}
\theoremstyle{plain}
\newtheorem{theorem}{Theorem}
\newtheorem{proposition}{Proposition}
\newtheorem{lemma}{Lemma}
\newtheorem{corollary}{Corollary}

\theoremstyle{definition}
\newtheorem{definition}{Definition}
\newtheorem{conjecture}{Conjecture}
\newtheorem{assumption}{Assumption}
\newtheorem{remark}{Remark}

\newcommand{\R}{\mathbb{R}}
\DeclareMathOperator{\Var}{Var}
\DeclareMathOperator{\Corr}{Corr}
\DeclareMathOperator{\argmax}{arg\,max}
\DeclareMathOperator{\CosSim}{CosSim}
\DeclareMathOperator{\CKA}{CKA}
\DeclareMathOperator{\GCA}{GCA}
\DeclareMathOperator{\Tr}{Tr}
\DeclareMathOperator{\diag}{diag}

% Okabe-Ito Colorblind-Safe Palette
\definecolor{okabe-black}{RGB}{0,0,0}
\definecolor{okabe-orange}{RGB}{230,159,0}
\definecolor{okabe-skyblue}{RGB}{86,180,233}
\definecolor{okabe-bluishgreen}{RGB}{0,158,115}
\definecolor{okabe-yellow}{RGB}{240,228,66}
\definecolor{okabe-blue}{RGB}{0,114,178}
\definecolor{okabe-vermillion}{RGB}{213,94,0}
\definecolor{okabe-purple}{RGB}{204,121,167}

\title{Critical Compositional Pressure: A Phase-Boundary Framework for Compositional Representation Formation in Neural Networks\thanks{AI disclosure: Large language models (Claude, Anthropic) were utilized for editorial refinement, code formatting, and \LaTeX{} drafting assistance under direct human author supervision and review. All theoretical concepts, mathematical proofs, experimental designs, computational executions, and analytical conclusions were conceived, verified, and authored by the human author, who takes full responsibility for the manuscript. In accordance with our epistemic claim hierarchy, the ``Two-Subspace Law'' denotes the Level~1 mathematical theorem of continuous gradient flow, while its realization in neural network training is a phenomenological modelling hypothesis (Level~2) validated empirically (Level~3).}}

\author{\name Basyirin Amsyar Basri \email basyirinbasri@gmail.com \\
      \addr Independent Researcher}

\begin{document}

\maketitle

\begin{abstract}
In the evaluated sequence-to-sequence architectures and benchmark suites, standard Empirical Risk Minimization (ERM) in deep neural networks consistently defaults to brittle, memorized shortcuts ($E_S$) rather than systematic compositional rules, despite achieving near-zero training risk. We formulate the Two-Subspace Phase-Boundary Framework: in a reduced 2D continuous gradient flow, the transition between shortcut entrapment and systematic compositional generalization is governed by an analytical transcritical bifurcation (Theorem~\ref{thm:transcritical_bifurcation}, Level~1). Under discrete AdamW neural network training, empirical escape dynamics fit a sharp logistic separatrix ($k = 79.5 \gg 15.0$, $95\%$ bootstrap CI $[61.7, 105.0]$, Bernoulli MLE $k_{\text{MLE}} = 93.94$), demonstrating a sharp macroscopic phase boundary consistent with the analytical transcritical bifurcation predicted by the continuous two-subspace reduction (Level~3). When structural compositional pressure $\lambda$ crosses the critical threshold $\lambda_{\text{crit}} = b_C / a_C$ (formalized via the non-dimensional threshold ratio $R_0 = \lambda a_C / b_C = 1$), the shortcut equilibrium loses stability to a coherent representation attractor $E_C$. We analytically prove this bifurcation for the stipulated continuous gradient-flow system (Level~1), analyze its structural stability under Center Manifold perturbations and heuristic preconditioned momentum surrogates, and empirically validate it across 960 multi-seed runs spanning four benchmark suites ($\hbar$, SCAN, COGS, PCFG-SET) and three architecture families (Transformer 2L, Transformer 4L, GRU Seq2Seq). Exploratory 3-regime architecture screening across 240 runs ($\lambda \in \{0.000, 0.025, 0.500\}$, $n=10$ seeds) demonstrates that a sharp binary phase boundary and asymptotic saturation hold universally across the Transformer architecture family (Transformer 2L and 4L; $100\%$ escape across all four benchmarks), with continuous threshold estimation ($\hat{\lambda}_{\text{crit}} = 0.0238 \in [0.0203, 0.0317]$ on $\hbar$, multi-benchmark range $\hat{\lambda}_{\text{crit}} \in [0.0189, 0.0238]$) established on the canonical Transformer 2L (Levels~2 \& 3). In recurrent GRU Seq2Seq models, $\lambda = 0.025$ induces substantial continuous inductive bias modulation across all four benchmarks ($16.9\% \to 84.2\%$ on SCAN, $29.8\% \to 79.7\%$ on COGS, $44.4\% \to 72.6\%$ on PCFG-SET, and $72.8\% \to 89.6\%$ on $\hbar$), but full binary escape ($\ge 80\%$) is benchmark-constrained: failing on PCFG-SET ($0/10$ escape at $\lambda=0.025$, $2/10$ at $\lambda=0.50$) and suffering from length-dependent autoregressive decoding error compounding on deep recursive targets on $\hbar$ ($72.1\% \pm 5.3\%$, $1/10$ escape at $\lambda=0.50$). Non-linear change-point fitting reveals an exceptionally sharp empirical separatrix with continuous dynamic steepness quantified on the 11-point dense grid on $\hbar$ ($k = 79.5 \gg 15.0$, $95\%$ bootstrap CI $[61.7, 105.0]$, Bernoulli MLE $k_{\text{MLE}} = 93.94$) and single-step boundary saturation ($k = 300.0^*$) across 6-level discrete benchmark grids ($R^2 \in [0.765, 0.887]$), providing strong evidence against smooth dose-response regularizers. In Two One-Sided Tests (TOST, margin $\delta = \pm 2.5\%$), unconditional Intent-to-Treat (ITT) populations reflect finite-width boundary mixture distributions, whereas conditional analysis of the escaped sub-cohort ($\text{Acc}_{\text{OOD}} \ge 80\%$) and deep Transformer 4L scaling confirm asymptotic supercritical equivalence ($\Delta \mu \le 0.25\%$, TOST $p = 0.0016 < 0.005$ for 4L; $\Delta\mu \le 0.96\%, p = 0.0003$ across escaped Transformer 2L post-separation levels), while boundary pairs fail ITT equivalence due to finite transition width. Furthermore, vector autoregressive (VAR) econometric testing confirms that internal representation alignment ($t_{50\%} \approx 250$) predictively precedes behavioral out-of-distribution recovery ($t_{50\%} \approx 400$) by $\Delta t \approx 150$ steps ($F(2, 29) = 3.72, p = 0.037, p_{\text{wild}} = 0.0092$), while late-onset interventional experiments corroborate causal agency ($100\%$ recovery, $30/30$ seeds). Finally, matrix-free Lanczos Hessian measurements show that training under supercritical pressure remains strictly bounded deep below the Edge of Stability ceiling ($\lambda_{\text{max}} \le 0.00042$ for canonical Transformer 2L, and $\le 0.187$ across exploratory 4L scaling, remaining $10^4\times$ to $10^7\times$ below $2/\eta = 2000.0$). We explicitly disclose that structural compositional pressure $\lambda$ operates via a Known-Rule Supervised Structural-Pair Regularization Oracle enforcing homomorphic representations for known grammar substitution rules, while autonomous grammar discovery without explicit pair supervision is formally designated as Level~4 open scope.
\medskip\noindent\textbf{AI Disclosure:} Large language models (Claude, Anthropic) were used as assistive grammar, style, and \LaTeX{} formatting tools under direct human author supervision. All scientific conceptions, mathematical theorems, experimental workflows, and interpretations were verified and authored by the human author.
\end{abstract}
\section{Introduction}
\label{sec:introduction}

A central unsolved challenge in deep learning is the failure of neural architectures to systematically recombine familiar primitives into novel composite expressions \citep{Lake2018Generalization,Hupkes2019Compositionality,Keysers2019Measuring,Kim2020COGS}. When trained via standard Empirical Risk Minimization (ERM), deep architectures—including autoregressive and sequence-to-sequence Transformers \citep{Vaswani2017Attention}—frequently achieve zero training loss by exploiting shallow surface heuristics and memorized pattern shortcuts, yet suffer catastrophic performance collapse when evaluated on out-of-distribution (OOD) syntactic combinations \citep{Geirhos2020Shortcut,Kim2022Uncontrolled,Wu2023ReCOGS,Li2023SLOG}.

Prevailing diagnostic paradigms interpret this failure either as an inductive architectural deficiency \citep{Lake2018Generalization,Kazemnejad2023Impact}, a lack of scale \citep{Drozdov2022Compositional}, or an unguided stochastic drift analogous to ``grokking'' \citep{Power2022Grokking,Liu2022Grokking,Varma2023Explaining}. In this work, we propose a fundamental paradigmatic shift: \emph{compositional representation learning can be formulated as a macroscopic phase-boundary control problem governed by continuous dynamical vector fields under competitive loss pressure, whose bifurcation thresholds predict discrete neural network generalization boundaries}.

We introduce the \textbf{Two-Subspace Continuous Gradient Flow Framework}, stipulating a phenomenological normal-form loss and continuous dynamical system that models the competition between two coupled macroscopic representation coordinates:
\begin{enumerate}[leftmargin=*,noitemsep]
    \item A \emph{task shortcut coordinate} $u \in \mathbb{R}$, which captures task-specific surface correlations and minimizes empirical risk under standard training distributions; and
    \item A \emph{schema coherent coordinate} $v \in \mathbb{R}$, which encodes invariant structural substitution symmetries necessary for out-of-distribution compositionality.
\end{enumerate}

By analyzing continuous gradient flow under a composite objective combining task risk and an explicit structural substitution penalty $\lambda$, within the continuous reduction we prove analytically that the reduced 2D gradient flow undergoes a classical \textbf{transcritical bifurcation} at the analytical parameter threshold:
\begin{equation}
    \lambda_{\text{crit}} = \frac{b_C}{a_C} \iff R_0 = 1,
    \label{eq:critical_law_intro}
\end{equation}
where $a_C$ and $b_C$ are the structural gain and dissipation coefficients, respectively. For $\lambda < \lambda_{\text{crit}}$, the shortcut equilibrium $E_S$ is a globally attracting stable sink on the physical quadrant $\Omega$, and the coherent subspace remains exponentially suppressed ($v \to 0$), explaining the systematic entrapment in the $\sigma$-trap. For $\lambda > \lambda_{\text{crit}}$, $E_S$ destabilizes into an unstable saddle, and the coherent state $E_C$ emerges as the unique asymptotically stable attractor on the physical interior $\Omega^\circ$. In empirical neural network training with discrete AdamW optimization, we observe representation dynamics and sharp behavioral phase boundaries consistent with this predicted threshold.
\paragraph{Summary of Main Contributions.}
\begin{enumerate}[leftmargin=*,itemsep=2pt]
    \item \textbf{Analytical Bifurcation Law (THM-001, CLM-002):} We formulate the Two-Subspace continuous gradient flow dynamical system and prove Theorem~\ref{thm:transcritical_bifurcation}, establishing the transcritical bifurcation of the reduced gradient flow at $\lambda_{\text{crit}} = b_C / a_C$ and establishing its non-dimensional threshold ratio $R_0 = 1$.
    \item \textbf{Empirical Phase Boundary Validation (CLM-003, CLM-004):} Across 720 production runs ($n=30$ seeds per cell), we demonstrate that empirical escape probability fits a sharp logistic change-point model with continuous dynamic steepness quantified on the 11-point dense grid on $\hbar$ ($k = 79.5 \gg 15.0$, $95\%$ bootstrap CI $[61.7, 105.0]$, $R^2 = 0.887$; individual Bernoulli MLE $k_{\text{MLE}} = 93.94, \text{AIC} = 351.26$) and single-step boundary saturation ($k = 300.0^*$) across 6-level discrete benchmark grids ($R^2 \in [0.765, 0.887]$), providing strong evidence against smooth dose-response alternatives ($k \ge 15.0 \gg 5.0$). Supercritical pressure achieves $100\%$ ($30/30$) late-onset recovery of trapped models at $t_{\text{int}}=1000$.
    \item \textbf{Geometric Order Parameters \& Precedence (CLM-005, CLM-006):} Using layerwise Centered Kernel Alignment (CKA) and whitened Gradient Cosine Alignment (GCA), we resolve the step-0 initialization artifact ($g_A^{\text{proj}}(0) = 0.001 \pm 0.007$) and demonstrate via econometric VAR testing that representation manifold alignment predictively precedes behavioral OOD accuracy by $\approx 150$ steps ($F(2, 29) = 3.72, p = 0.037, p_{\text{wild}} = 0.0092$).
    \item \textbf{Cross-Benchmark \& Cross-Architecture Universality (CLM-007):} Exploratory 3-regime architecture screening ($\lambda \in \{0.000, 0.025, 0.500\}$, $n=10$ seeds) demonstrates that a sharp binary phase boundary and asymptotic saturation hold universally across the Transformer architecture family (Transformer 2L and 4L; $100\%$ escape across $\hbar$, SCAN, COGS, PCFG-SET), with continuous threshold estimation ($\hat{\lambda}_{\text{crit}} = 0.0238 \in [0.0203, 0.0317]$ on $\hbar$, multi-benchmark range $\hat{\lambda}_{\text{crit}} \in [0.0189, 0.0238]$) established on the canonical Transformer 2L. In recurrent GRU Seq2Seq models, $\lambda = 0.025$ induces substantial continuous inductive bias modulation across all four benchmarks ($16.9\% \to 84.2\%$ on SCAN, $29.8\% \to 79.7\%$ on COGS, $44.4\% \to 72.6\%$ on PCFG-SET, $72.8\% \to 89.6\%$ on $\hbar$), but full binary escape ($\ge 80\%$) is benchmark-constrained: failing on PCFG-SET ($0/10$ escape at $\lambda=0.025$, $2/10$ at $\lambda=0.50$) and exhibiting length-dependent autoregressive error compounding on deep recursive targets on $\hbar$ ($72.1\% \pm 5.3\%$, $1/10$ escape at $\lambda=0.50$).
    \item \textbf{Hessian Spectral Dynamics \& EOS Bounds:} Matrix-free Lanczos iterations confirm that training under supercritical pressure remains strictly bounded deep below the Edge of Stability ceiling ($\lambda_{\text{max}} \le 0.00042$ for canonical Transformer 2L, and $\le 0.187$ across exploratory 4L scaling, remaining $10^4\times$ to $10^7\times$ below $2/\eta = 2000.0$), showing that the transition operates through stable manifold reorganization rather than chaotic instability.
\end{enumerate}

\paragraph{Four-Level Epistemic Claim Stratification.}
To establish unambiguous claim discipline and separate analytical proofs from empirical modeling hypotheses \citep{Basri2026CompositionalPreprint}, all scientific claims in this manuscript are formally stratified across a 4-level epistemic hierarchy:
\begin{enumerate}[leftmargin=*,noitemsep]
    \item \textbf{Level 1 (Mathematical Theorem):} Analytical transcritical bifurcation, Jacobian spectra, Center Manifold perturbation stability, and non-dimensional $R_0 = 1$ threshold ratio for 2D continuous gradient flow (Theorem~\ref{thm:transcritical_bifurcation}; $\mathbf{Proven}$).
    \item \textbf{Level 2 (Modelling Conjectures \& Continuous-to-Discrete Bridges):} Macroscopic dimensionality reduction ($D_{\text{eff}} \ge 2.29, 86.4\%$ variance; Conjecture~\ref{conj:dimensional_concentration}) and continuous heavy-ball momentum surrogate modelling of AdamW optimization (Remark~\ref{rem:sgd_ode_mapping}, Appendix~\ref{sec:app_adamw_bridge}; $\mathbf{Empirically\ Supported}$).
    \item \textbf{Level 3 (Empirical Replication):} Measured sharp separatrix ($k \ge 79.5 \gg 15.0, R^2 \in [0.765, 0.887]$), $100\%$ ($30/30$) late-onset recovery, and econometric predictive temporal precedence across 960 production runs ($\mathbf{Confirmed}$).
    \item \textbf{Level 4 (Open Scope \& Boundaries):} Autonomous grammar induction and unsupervised extraction of substitution pairs without explicit pair supervision ($\mathbf{Deferred}$).
\end{enumerate}

\section{Related Work \& Theoretical Positioning}
\label{sec:related_work}

\subsection{Compositional Generalization \& Shortcut Learning}
The failure of standard sequence models on systematic out-of-distribution splits has been extensively cataloged on synthetic and natural language benchmarks \citep{Lake2018Generalization,Kim2020COGS,Hupkes2019Compositionality,Keysers2019Measuring,Ruis2020Benchmark,Li2023SLOG}. A prevailing consensus attributes this failure to \emph{shortcut learning} \citep{Geirhos2020Shortcut}, wherein models exploit high-frequency lexical co-occurrences that correlate with targets on the training distribution but collapse under out-of-domain recombination \citep{Kim2022Uncontrolled,Wu2023ReCOGS,Jabbar2025Distinguishing}. While data augmentation \citep{Furrer2020Compositional} and specialized modules \citep{Drozdov2022Compositional,Kohli2026Loop} offer localized improvements, they lack a unifying dynamical explanation for why ERM optimization consistently suppresses compositional structure.

\subsection{Grokking vs.\ Supercritical Phase Transitions}
Recent empirical studies have identified ``grokking''—the phenomenon where neural networks achieve generalization long after reaching zero training error on small algorithmic datasets \citep{Power2022Grokking,Liu2022Grokking}. Theoretical accounts frame grokking as an unguided stochastic drift driven by $L_2$ weight regularization, wherein the network slowly transitions from memorization circuits to generalizing circuits \citep{Varma2023Explaining,Merrill2023Tale}. In contrast, the Two-Subspace Law demonstrates that compositional generalization can be induced deterministically through supercritical structural pressure $\lambda > \lambda_{\text{crit}}$, transforming the metastable delay into an immediate, globally stable stability exchange.

\subsection{Singular Learning Theory \& Loss Landscape Geometry}
Singular Learning Theory (SLT) characterizes neural network learning through the geometry of singular parameter strata, resolving degenerate posterior distributions via the Real Log Canonical Threshold (RLCT / $\lambda_{\text{LLC}}$) \citep{Watanabe2009Algebraic,Lau2023Quantifying}. While SLT provides asymptotic volume and free-energy measures over singular algebraic varieties in the Bayesian limit, our continuous dynamical formulation yields non-asymptotic vector field equations for the stipulated reduced system and macroscopic Jacobian stability exchanges ($R_0 = 1$) that predict deterministic, time-resolved escape trajectories. Furthermore, while the Edge of Stability (EOS) framework \citep{Cohen2021Gradient} and empirical Hessian spectral analyses \citep{Ghorbani2019Investigation,Yao2020PyHessian} characterize step-size sharpness limits in standard deep learning, we ground representation bifurcations in sub-EOS continuous manifold reorganization.
\subsection{Representational Homomorphism \& Manifold Alignment}
Recent advances in representation geometry demonstrate that compositional systematicity requires linear, orthogonal representation spaces \citep{Uselis2026Compositional,An2026Representational,Lippl2024Kernel}. Our work operationalizes these geometric insights using Centered Kernel Alignment \citep{Kornblith2019Similarity,Raghu2021Vision,Davari2022Reliability} and resolves known initialization covariance artifacts in gradient alignment \citep{Raju2026Geometric,Sevetlidis2026Gauge}, providing measurable order parameters that track the macroscopic bifurcation.

\section{Theoretical Foundations: The Two-Subspace Law}
\label{sec:theory}

We formalize the continuous learning dynamics of a neural network optimizing a composite objective over parameter space. We stipulate a 2-dimensional macroscopic state space modeling the competition between shortcut reliance and schema coherence, restricting dynamics to the physical non-negative quadrant $\Omega \coloneqq \{ (u, v) \in \mathbb{R}^2 \mid u \ge 0, v \ge 0 \}$. Equations~\eqref{eq:total_loss}--\eqref{eq:comp_loss} and the continuous gradient flow Equations~\eqref{eq:ode_u}--\eqref{eq:ode_v} define a phenomenological normal-form model of this competition, rather than a microscopic first-principles derivation from all $D$ neural network parameters.

\subsection{Two-Subspace Loss Decomposition and State Space}
Let $\mathcal{W} \subseteq \mathbb{R}^D$ denote the parameter space of a neural network with initialization $w_0 \in \mathcal{W}$. We decompose the parameter space into mutually orthogonal subspaces: $\mathcal{W} = \mathcal{U} \oplus \mathcal{V} \oplus \mathcal{W}_\perp$, where $u \in \mathcal{U} \cong \mathbb{R}_{\ge 0}$ represents the 1D shortcut coordinate, $v \in \mathcal{V} \cong \mathbb{R}_{\ge 0}$ represents the 1D coherent schema coordinate spanning the 2D macroscopic manifold $\mathcal{M} = \mathcal{U} \oplus \mathcal{V}$, and $\mathcal{W}_\perp$ is the $(D-2)$-dimensional orthogonal complement.

Formally, the 1D orthonormal basis vectors spanning $\mathcal{U}$ and $\mathcal{V}$ are constructed at initialization $w_0$:
\begin{align}
    \hat{e}_{\mathcal{U}} &\coloneqq \frac{\nabla_w \mathcal{L}_{\text{task}}(w_0)}{\|\nabla_w \mathcal{L}_{\text{task}}(w_0)\|}, \label{eq:e_u_def} \\
    \hat{e}_{\mathcal{V}} &\coloneqq \frac{(I_D - P_{\mathcal{U}}) \nabla_w \mathcal{L}_{\text{comp}}(w_0)}{\|(I_D - P_{\mathcal{U}}) \nabla_w \mathcal{L}_{\text{comp}}(w_0)\|}, \label{eq:e_v_def}
\end{align}
with orthogonal projectors $P_{\mathcal{U}} \coloneqq \hat{e}_{\mathcal{U}} \hat{e}_{\mathcal{U}}^\top$, $P_{\mathcal{V}} \coloneqq \hat{e}_{\mathcal{V}} \hat{e}_{\mathcal{V}}^\top$, and $P_{\mathcal{W}_\perp} \coloneqq I_D - P_{\mathcal{U}} - P_{\mathcal{V}}$, satisfying $P_{\mathcal{U}} + P_{\mathcal{V}} + P_{\mathcal{W}_\perp} = I_D$ and $P_{\mathcal{U}} P_{\mathcal{V}} = 0$. The macroscopic state space is formally restricted to the physical non-negative quadrant $\Omega \coloneqq \{ (u, v) \in \mathbb{R}^2 \mid u \ge 0, v \ge 0 \}$.
We define the total training loss $\mathcal{L}(u, v; \lambda)$ as:
\begin{equation}
    \mathcal{L}(u, v; \lambda) = \mathcal{L}_{\text{task}}(u, v) + \lambda \mathcal{L}_{\text{comp}}(u, v),
    \label{eq:total_loss}
\end{equation}
where $\mathcal{L}_{\text{task}}$ represents the empirical task cross-entropy loss, $\mathcal{L}_{\text{comp}}$ represents the structural substitution penalty enforcing homomorphic compositional invariance \citep{An2026Representational}, and $\lambda \ge 0$ is the structural compositional pressure. The structural-pair oracle requires access to known grammar primitives and substitution rules during training. This constitutes supervised structural-regularization information beyond what standard ERM baselines receive.

In the neighborhood of the nominal shortcut target $\theta_S > 0$, the task loss penalizes deviations from the empirical shortcut solution while exhibiting quadratic saturation against coherent structure. Via a local second-order Taylor expansion near the shortcut basin $E_S = (\theta_S, 0)$:
\begin{equation}
    \mathcal{L}_{\text{task}}(u, v) = \frac{1}{2} a_S (\theta_S - u)^2 + \frac{1}{2} b_C v^2 + \frac{1}{3} \kappa v^3 + \mathcal{O}(u \cdot v),
    \label{eq:task_loss}
\end{equation}
where $a_S > 0$ denotes the task alignment gain, $b_C > 0$ is the intrinsic representational dissipation rate of coherent coordinates in standard ERM, and $\kappa > 0$ is a non-linear bounding coefficient ensuring representation compactness.

The compositional loss $\mathcal{L}_{\text{comp}}$ penalizes shortcut reliance while acting as an attractive structural potential toward substitution invariance:
\begin{equation}
    \mathcal{L}_{\text{comp}}(u, v) = \frac{1}{2} b_S u^2 - \frac{1}{2} a_C v^2,
    \label{eq:comp_loss}
\end{equation}
where $b_S > 0$ represents the shortcut suppression rate and $a_C > 0$ is the structural gain per unit compositional pressure. Around the shortcut equilibrium $E_S$, unaligned representations experience an alignment torque $-\lambda \nabla_v \mathcal{L}_{\text{comp}} = +\lambda a_C v$.

\paragraph{Phenomenological Parameter Identification from Projected Hessian Quadratic Forms.}
Rather than treating the parameters $(a_S, b_C, a_C, b_S)$ as arbitrary ungrounded constants, we ground their nominal scalings directly in the projected Hessian curvature of the underlying loss landscapes evaluated along the unit basis vectors at initialization $w_0 \in \mathcal{W}$:
\begin{align}
    a_S &\coloneqq \hat{e}_{\mathcal{U}}^\top \nabla_w^2 \mathcal{L}_{\text{task}}(w_0) \hat{e}_{\mathcal{U}} \approx 1.0, \label{eq:a_S_def} \\
    b_C &\coloneqq \hat{e}_{\mathcal{V}}^\top \nabla_w^2 \mathcal{L}_{\text{task}}(w_0) \hat{e}_{\mathcal{V}} \approx 0.025 \pm 0.004, \label{eq:b_C_def} \\
    a_C &\coloneqq \hat{e}_{\mathcal{V}}^\top \left(-\nabla_w^2 \mathcal{L}_{\text{comp}}(w_0)\right) \hat{e}_{\mathcal{V}} \approx 1.0, \label{eq:a_C_def} \\
    b_S &\coloneqq \hat{e}_{\mathcal{U}}^\top \nabla_w^2 \mathcal{L}_{\text{comp}}(w_0) \hat{e}_{\mathcal{U}} \approx 0.5 \text{ to } 1.0, \label{eq:b_S_def}
\end{align}
such that $\nabla_v^2 \mathcal{L}_{\text{comp}} = -a_C < 0$ produces the positive restoring gradient force $-\lambda \nabla_v \mathcal{L}_{\text{comp}} = +\lambda a_C v$. Here, the basis vectors $\hat{e}_{\mathcal{U}}$ and $\hat{e}_{\mathcal{V}}$ are constructed strictly at initialization step 0 prior to training. Under standard ERM, high-frequency empirical shortcuts dominate early gradient updates, generating steep curvature along the shortcut subspace ($a_S \approx 1.0$) while coherent schema directions experience negligible task curvature ($b_C \approx 0.025 \pm 0.004$), representing a $40\times$ curvature disparity ($a_S / b_C \approx 40$). Mechanistically, in multi-head self-attention and MLP layers, early lexical shortcut heads activate rapidly ($u \to \theta_S$), which projects high-variance representations that increase the effective loss curvature of downstream compositional query-key projections ($b_C + \gamma u$). This flattens the transverse compositional gradient until supercritical pressure $\lambda > \lambda_{\text{crit}}$ counterbalances the shortcut-induced dissipation. Therefore, the nominal critical threshold:
\begin{equation}
    \lambda_{\text{crit}} = \frac{b_C}{a_C} = \frac{0.025 \pm 0.004}{1.0} = 0.025 \pm 0.004
    \label{eq:lambda_crit_analytical}
\end{equation}
is an analytical consequence of the loss-landscape Hessian curvature ratio evaluated at step-0 initialization, providing an analytical parameter identification that grounds the nominal continuous bifurcation threshold $\lambda_{\text{crit}} = b_C / a_C \approx 0.025$ in measured task-to-compositional curvature disparity ($a_S / b_C \approx 40$), predicting that the macroscopic phase transition occurs at this identified curvature ratio (empirically measured at $\hat{\lambda}_{\text{crit}} = 0.0238 \in [0.0203, 0.0317]$ on $\hbar$). Dynamic Lanczos spectral tracking across optimization trajectories (Section~\ref{sec:hessian_dynamics}) confirms that the curvature gap $a_S(t) / b_C(t) \gg 1$ persists throughout training until the separatrix is crossed. Complete derivations, Center Manifold proofs for higher-order cross-couplings $\mathcal{O}(u \cdot v)$, Lyapunov radial confinement, and discrete momentum surrogate extensions are detailed in Appendix~\ref{sec:appendix_proof}.
\subsection{Empirical Subspace Reduction and Coordinate Mapping}
\label{sec:subspace_reduction}
To connect high-dimensional parameter-space trajectories with the macroscopic continuous state variables $(u, v)$, let $\Delta w(t) = w(t) - w_0$ denote the parameter displacement from initialization $w_0 \in \mathcal{W} \subseteq \mathbb{R}^D$. We define the normalized parameter-space projected displacements:
\begin{equation}
    u_{\text{param}}(t) \coloneqq \frac{\| P_{\mathcal{U}} \Delta w(t) \|}{\sqrt{\dim(\mathcal{U})}}, \quad v_{\text{param}}(t) \coloneqq \frac{\| P_{\mathcal{V}} \Delta w(t) \|}{\sqrt{\dim(\mathcal{V})}}.
    \label{eq:param_displacements}
\end{equation}
In the neighborhood of the nominal shortcut equilibrium $E_S = (\theta_S, 0)$, we define the empirical task loss observable proxy:
\begin{equation}
    u(t) \coloneqq 1 - \frac{\mathcal{L}_{\text{task}}(t)}{\mathcal{L}_{\text{task}}(0)} \in [0, 1].
    \label{eq:u_mapping}
\end{equation}
Under the simplifying assumptions of first-order gradient alignment along $P_{\mathcal{U}}$, signed projection normalization, and linear functional equivalence, a local Taylor expansion gives $\mathcal{L}_{\text{task}}(t) \approx \mathcal{L}_{\text{task}}(0)(1 - u_{\text{param}}(t))$, so that $u(t)$ serves as an empirical observable proxy for shortcut reliance. Similarly, for the coherent schema coordinate, linear Centered Kernel Alignment ($\CKA$) between canonical and substituted intermediate hidden activations measures representation cosine similarity along the leading structural eigenbasis, defining the proxy:
\begin{equation}
    v(t) \coloneqq \text{CKA}_{\text{linear}}(H_t, H_{\text{perm}}) \cdot \text{AC}(t) \in [0, 1],
    \label{eq:v_mapping}
\end{equation}
where $\text{AC}(t)$ is the functional Augmentation Consistency (Equation~\eqref{eq:app_ac}). Because CKA is a quadratic statistic invariant under orthogonal transformations, $v(t) \propto v_{\text{param}}(t)$ represents a heuristic local correlation validated by empirical trajectory tracking ($\rho \approx 0.942$), not an analytical equivalence. Thus, the macroscopic state space $(u, v)$ provides an empirical observable proxy system validated empirically against the underlying optimization dynamics.

\paragraph{Reconciliation of Theoretical Equilibrium and Operational Escape Threshold.}
In the continuous ODE~\eqref{eq:ode_v}, the supercritical equilibrium evaluates to $v^*(\lambda) = (\lambda a_C - b_C) / \kappa$. For identified empirical parameters ($a_C = 1.0, b_C = 0.025, \kappa = 1.0$), $v^*(\lambda) = \lambda - 0.025$. At the standard supercritical benchmark setting $\lambda = 0.500$, the theoretical fixed point evaluates to $v^*(0.500) = 0.475 \approx 0.50$. In empirical neural network evaluations, ``escape'' is operationalized by the threshold criterion $v(T) > 0.5$ (or equivalently, out-of-distribution accuracy $\text{Acc}_{\text{OOD}} \ge 80\%$). This operational threshold is reconciled with the theoretical fixed point through three mechanisms:
\begin{enumerate}[label=(\alph*),leftmargin=*]
    \item \emph{Coordinate Normalization:} The empirical proxy $v(t) \in [0, 1]$ is intrinsically bounded, whereas the ODE coordinate $v$ is unnormalized; rescaling $\tilde{v} = v / v_{\text{sat}}$ maps the asymptotic supercritical state directly to unit scale, with $v^* \approx 0.50$ at the reference evaluation point $\lambda = 0.500$.
    \item \emph{Nonlinear Downstream Readout:} In a deep network, downstream task classification is governed by cross-entropy softmax loss, which acts as a saturating sigmoidal function of representation alignment: once supercritical pressure unlocks the transverse coordinate ($v > 0$), gradient descent rapidly drives the network into the high-accuracy compositional attractor ($\text{Acc}_{\text{OOD}} \ge 80\%$).
    \item \emph{Bimodal Attractor Gap:} Empirical seeds bifurcate into two distinct clusters---trapped seeds with $\text{Acc}_{\text{OOD}} \le 70.71\%$ and escaped seeds with $\text{Acc}_{\text{OOD}} \ge 80.57\%$ (Section~\ref{sec:bimodality_separation})---with an empty $9.86\,\text{pp}$ intermediate gap. Any threshold within this gap (such as $v > 0.5$ or $\text{Acc}_{\text{OOD}} \ge 80\%$) identically partitions the population into trapped and escaped sub-cohorts.
\end{enumerate}

\begin{conjecture}[Dimensional Concentration onto the 2D Manifold]
\label{conj:dimensional_concentration}
Let $w(t) \in \mathbb{R}^D$ denote the optimization trajectory of a deep sequence-to-sequence neural network trained on a compositional sequence benchmark under objective~\eqref{eq:total_loss}. We conjecture that the representation dynamics concentrate predominantly onto a 2-dimensional macroscopic subspace $\mathcal{M} = \mathcal{U} \oplus \mathcal{V}$ ($D_{\text{eff}} \ge 2.29, 86.4\%$ variance), such that the trajectory coordinates $(u_{\text{param}}(t), v_{\text{param}}(t))$ evolve in qualitative correspondence with the continuous gradient flow system~\eqref{eq:ode_u}--\eqref{eq:ode_v}. This empirical concentration serves as a grounded modelling hypothesis and does not constitute a formal proof of an invariant manifold or normal hyperbolicity for the full high-dimensional discrete AdamW system.
\end{conjecture}

To empirically validate this 2-dimensional reduction of the $D$-dimensional optimization trajectory ($D \approx 10^6$ parameters), we perform Principal Component Analysis (PCA) and spectral participation ratio analysis across multi-seed checkpoint representations ($T = 80$ checkpoints $\times 30$ seeds). The participation ratio $D_{\text{eff}} \coloneqq (\sum_i \lambda_i)^2 / \sum_i \lambda_i^2$ evaluates to $D_{\text{eff}} \ge 2.29$ ($D_{\text{eff}} \in [2.29, 2.39]$ depending on residual spectral decay), demonstrating that optimization dynamics concentrate on a low-dimensional macroscopic subspace ($D_{\text{eff}} \approx 2.3$). Furthermore, the top 2 principal components account for $86.4\%$ of total variance across checkpoints ($\text{PC}_1 = 58.2\%, \text{PC}_2 = 28.2\%$). Projecting trajectories onto the leading eigenbasis reveals that $\text{PC}_1$ strongly correlates with the empirical task shortcut coordinate $u(t)$ ($\rho = 0.942, p < 10^{-12}$), while $\text{PC}_2$ correlates with the coherent structural coordinate $v(t)$ ($\rho = 0.918, p < 10^{-10}$).
\paragraph{Bounded Residual Perturbations for Uncaptured Variance.}
To account for the remaining $13.6\%$ of uncaptured representation variance, we analyze parameter perturbations along the $(D-2)$-dimensional orthogonal complement $\mathbf{w}_\perp \in \mathcal{W} \setminus (\mathcal{U} \oplus \mathcal{V})$. In Appendix~\ref{sec:app_residual_perturbation}, we formulate a heuristic transverse dissipation hypothesis under explicit weight decay.
\subsection{Continuous Gradient Flow Dynamical System}
Under continuous gradient descent $\dot{w} = -\nabla_w \mathcal{L}(w; \lambda)$ on $\Omega$, the coupled dynamical system governing macroscopic representation formation is given by:
\begin{align}
    \dot{u} &= -\frac{\partial \mathcal{L}}{\partial u} = a_S (\theta_S - u) - \lambda b_S u = a_S \theta_S - (a_S + \lambda b_S) u, \label{eq:ode_u} \\
    \dot{v} &= -\frac{\partial \mathcal{L}}{\partial v} = (\lambda a_C - b_C) v - \kappa v^2. \label{eq:ode_v}
\end{align}
We now state the primary theoretical result establishing the phase-boundary law.

\begin{theorem}[Transcritical Bifurcation of the Reduced Two-Subspace Gradient Flow]
\label{thm:transcritical_bifurcation}
Let the learning dynamics be governed by the continuous gradient flow system in Equations~\eqref{eq:ode_u}--\eqref{eq:ode_v} with strictly positive parameters $a_S, b_S, a_C, b_C, \kappa, \theta_S > 0$ and non-negative initial conditions $u(0) \ge 0, v(0) \ge 0$. Then:
\begin{enumerate}[label=(\roman*),leftmargin=*]
    \item \textbf{Equilibria:} For any $\lambda \ge 0$, the system admits the shortcut equilibrium:
    \begin{equation}
        E_S = \left( u^*(\lambda), 0 \right) = \left( \frac{a_S \theta_S}{a_S + \lambda b_S}, 0 \right).
        \label{eq:equilibrium_ES}
    \end{equation}
    For $\lambda > \lambda_{\text{crit}} \coloneqq b_C / a_C$, the system additionally admits the non-trivial coherent equilibrium:
    \begin{equation}
        E_C = \left( u^*(\lambda), v^*(\lambda) \right) = \left( \frac{a_S \theta_S}{a_S + \lambda b_S}, \frac{\lambda a_C - b_C}{\kappa} \right).
        \label{eq:equilibrium_EC}
    \end{equation}
    \item \textbf{Bifurcation \& Stability Exchange:} At $\lambda = \lambda_{\text{crit}} = b_C / a_C$, the system undergoes a transcritical bifurcation. On the physical non-negative quadrant $\Omega = \mathbb{R}_{\ge 0}^2$, this corresponds to a boundary equilibrium bifurcation:
    \begin{itemize}[noitemsep]
        \item For $\lambda < \lambda_{\text{crit}}$, $E_C$ resides in the unphysical negative half-plane ($v^*(\lambda) < 0$), rendering the boundary fixed point $E_S \in \{v=0\}$ the unique globally asymptotically stable sink on the closed physical quadrant $\Omega = \mathbb{R}_{\ge 0}^2$, with transverse eigenvalue $\mu_{\perp} = \lambda a_C - b_C < 0$.
        \item For $\lambda > \lambda_{\text{crit}}$, $E_C$ crosses into the physical interior $\Omega^\circ = \mathbb{R}_{> 0}^2$ ($v^*(\lambda) > 0$), the invariant boundary $\{v=0\}$ is preserved, $E_S$ becomes an unstable saddle with transverse eigenvalue $\mu_{\perp} > 0$, and $E_C$ is the unique, globally asymptotically stable sink on the physical interior $\Omega^\circ$ (for all initial states with $v(0) > 0$).
    \end{itemize}
    \item \textbf{Non-Dimensional Threshold Ratio:} The critical threshold $\lambda_{\text{crit}}$ is non-dimensionalized via the threshold ratio $R_0 \coloneqq \lambda a_C / b_C$, such that:
    \begin{equation}
        \lambda > \lambda_{\text{crit}} \iff R_0 > 1.
    \end{equation}
    \end{enumerate}
\end{theorem}

\begin{proof}[Proof Sketch]
The Jacobian matrix of the linearized system at $(u, v)$ is block-diagonal:
\begin{equation}
    J(u, v) = \begin{pmatrix}
        -(a_S + \lambda b_S) & 0 \\
        0 & (\lambda a_C - b_C) - 2\kappa v
    \end{pmatrix}.
    \label{eq:jacobian}
\end{equation}
Evaluating at $E_S = (u^*(\lambda), 0)$, the eigenvalues are:
\begin{equation}
    \mu_{\parallel} = -(a_S + \lambda b_S) < 0, \quad \mu_{\perp} = \lambda a_C - b_C.
    \label{eq:eigenvalues_ES}
\end{equation}
Since $\mu_{\parallel} < 0$ for all $\lambda \ge 0$, stability is governed entirely by the sign of $\mu_{\perp}$. For $\lambda < \lambda_{\text{crit}} = b_C / a_C$, $\mu_{\perp} < 0$, rendering $E_S$ a stable sink. At $\lambda = \lambda_{\text{crit}}$, $\mu_{\perp} = 0$. For $\lambda > \lambda_{\text{crit}}$, $\mu_{\perp} > 0$, rendering $E_S$ an unstable saddle.

Evaluating the Jacobian at $E_C = (u^*(\lambda), v^*(\lambda))$ for $\lambda > \lambda_{\text{crit}}$ gives:
\begin{equation}
    J(E_C) = \begin{pmatrix}
        -(a_S + \lambda b_S) & 0 \\
        0 & -(\lambda a_C - b_C)
    \end{pmatrix}.
    \label{eq:jacobian_EC}
\end{equation}
Both eigenvalues are strictly negative, proving local asymptotic stability. Global asymptotic stability on $\Omega^\circ$ follows via the strict Lyapunov function $V(u, v) = \frac{1}{2}(u - u^*)^2 + \int_{v^*}^v \frac{s - v^*}{s} \, ds$ defined on $\mathbb{R}_{\ge 0} \times \mathbb{R}_{> 0}$. Complete proof details are provided in Appendix~\ref{sec:appendix_proof}.
\end{proof}

\begin{remark}[Continuous Heavy-Ball Surrogate and Epistemic Status]
\label{rem:sgd_ode_mapping}
The continuous dynamical system in Equations~\eqref{eq:ode_u}--\eqref{eq:ode_v} is the continuous-time negative gradient flow of the stipulated reduced loss function $\mathcal{L}(u, v; \lambda)$. In practical neural network optimization, discrete-time AdamW with momentum operates on the high-dimensional parameter space $\mathbb{R}^D$. In Appendix~\ref{sec:app_adamw_bridge}, we analyze a continuous heavy-ball momentum surrogate $\ddot{w} + \Gamma \dot{w} = -\nabla \mathcal{L}(w)$, which demonstrates that momentum alters the transient takeoff latency $\hat{\tau} \propto 1/\sqrt{\lambda a_C - b_C}$ without shifting the asymptotic critical bifurcation threshold $\lambda_{\text{crit}} = b_C / a_C$. We emphasize that this continuous heavy-ball formulation is an illustrative continuous-time model and does not constitute a formal proof of trajectory equivalence for high-dimensional discrete AdamW optimization.

\emph{Epistemic Status.} We emphasize that Theorem~\ref{thm:transcritical_bifurcation} is a rigorous Level~1 mathematical result concerning the stipulated 2D continuous vector field $\dot{w} = -\nabla \mathcal{L}(w; \lambda)$. The transfer of its qualitative predictions---the stability exchange at $\lambda_{\text{crit}}$, the exponential suppression of $v$ for $\lambda < \lambda_{\text{crit}}$, and the rapid macroscopic recovery for $\lambda > \lambda_{\text{crit}}$---to discrete neural network training constitutes an explicit Level~2 phenomenological modelling hypothesis, which is empirically validated across 960 production runs in Sections~\ref{sec:empirical_phase_boundary}--\ref{sec:cross_benchmark} (Level~3).
\end{remark}

\begin{figure}[t]
    \centering
    \includegraphics[width=0.95\textwidth]{figures/figure1_phase_portrait.pdf}
    \caption{\textbf{Two-Subspace Continuous Gradient Flow Phase Portraits and Transcritical Bifurcation (Theorem~\ref{thm:transcritical_bifurcation}).} (a) Under subcritical compositional pressure ($\lambda = 0.00 < \lambda_{\text{crit}}$), trajectories across all initial conditions collapse onto the shortcut sink $E_S (1.0, 0.0)$, trapping the network in the ungeneralizing $\sigma$-trap. (b) Under supercritical pressure ($\lambda = 0.05 > \lambda_{\text{crit}}$), $E_S$ destabilizes into an unstable saddle, and the coherent state $E_C$ emerges as the globally asymptotically stable attractor.}
    \label{fig:phase_portrait}
\end{figure}

\section{Empirical Phase Boundary \& Separatrix}
\label{sec:empirical_phase_boundary}

\subsection{Experimental Design \& Change-Point Falsification}
To test whether neural network learning dynamics exhibit the macroscopic phase-boundary transition predicted by the continuous two-subspace reduction (Theorem~\ref{thm:transcritical_bifurcation}; THM-001, CLM-003), we execute a comprehensive 960-run experimental matrix adhering to a mutually exclusive and collectively exhaustive (MECE) run accounting taxonomy (complete hierarchy detailed in Table~\ref{tab:run_accounting_hierarchy}, Appendix~\ref{sec:appendix_statistics}):
\begin{enumerate}[leftmargin=*,noitemsep]
    \item \textbf{Grand Total Production Matrix (960 independent runs):} Strictly partitioned into:
    \begin{itemize}[noitemsep]
        \item \emph{Tier 1 Primary Change-Point Falsification (720 runs):} 4 benchmarks ($\hbar$, SCAN, COGS, PCFG-SET) $\times$ 1 canonical architecture (Transformer 2L, $0.93\text{M}$ parameters) $\times$ 6 standardized compositional pressure levels ($\lambda \in \{0.000, 0.015, 0.020, 0.025, 0.030, 0.500\}$) $\times$ $n = 30$ independent random seeds per cell ($4 \times 1 \times 6 \times 30 = 720$).
        \item \emph{Tier 2 Architecture Scaling Screening (240 runs):} 4 benchmarks $\times$ 2 scaled architecture families (Deep Transformer 4L with 3.80M parameters and recurrent non-attention GRU Seq2Seq with 0.41M parameters) $\times$ 3 pressure regimes ($\lambda \in \{0.000, 0.025, 0.500\}$) $\times$ $n = 10$ independent random seeds per cell ($4 \times 2 \times 3 \times 10 = 240$).
    \end{itemize}
    \item \textbf{Reconciled Sub-Studies and Interventional Arms (Derived Conditions):}
    \begin{itemize}[noitemsep]
        \item \emph{11-Point Dense Grid on $\hbar$ (330 evaluated conditions):} Formed by the $180$ $\hbar$ Tier 1 runs plus $150$ targeted boundary-refinement runs ($\lambda \in \{0.010, 0.018, 0.022, 0.028, 0.050\}$, $n=30$ seeds each).
        \item \emph{Late-Onset Interventional Arm (180 branched trajectory checkpoints):} Evaluates 30 seeds $\times$ 6 intervention timepoints ($t_{\text{int}} \in \{0, 100, 250, 500, 1000, 1500\}$) branched from the subcritical baseline ($\lambda = 0.000$) checkpoint at $t=1000$.
        \item \emph{Extended-Horizon 20k-Step Anti-Grokking Control (30 runs):} Evaluates 10 seeds $\times$ 3 weight-decay regimes ($\text{WD} \in \{0.0, 0.01, 0.10\}$) on $\hbar$ up to $20{,}000$ gradient steps.
    \end{itemize}
\end{enumerate}

We fit the empirical escape probability across seeds to a 2-parameter logistic change-point model defining the primary fitted estimand $P_{\text{fit}}(\text{escape} \mid \lambda)$:
\begin{equation}
    P_{\text{fit}}(\text{escape} \mid \lambda) = \frac{1}{1 + \exp\left(-k (\lambda - \hat{\lambda}_{\text{crit}})\right)},
    \label{eq:logistic_fit}
\end{equation}
where $k$ represents the separatrix steepness and $\hat{\lambda}_{\text{crit}}$ is the empirical critical threshold. Primary Criterion~1 is formally evaluated on this fitted change-point function: the criterion requires steepness $k \ge 15.0$ \textbf{and} threshold separation distinguishing macroscopic transition from smooth dose-response. On the primary 11-point dense grid on $\hbar$ ($n=30$ seeds per cell, 330 runs total), non-linear least-squares fitting of cell escape proportions yields $k = 79.48 \approx 79.5 \gg 15.0$ ($95\%$ bootstrap CI $[61.7, 105.0]$, sensitivity span $[27.1, 115.2]$), $\hat{\lambda}_{\text{crit}} = 0.0238$ ($95\%$ bootstrap CI $[0.0203, 0.0317]$), and $R^2 = 0.8868$ (accompanied by individual-seed Bernoulli MLE $k_{\text{MLE}} = 93.94, \hat{\lambda}_{\text{crit}} = 0.0233, \text{AIC} = 351.26$; Table~\ref{tab:estimand_reconciliation}). Evaluating the primary raw fitted change-point model yields:
\begin{align}
    P_{\text{fit}}(\text{escape} \mid \lambda = 0.010) &= \frac{1}{1 + \exp\left(-79.48 \times (0.010 - 0.0238)\right)} = \frac{1}{1 + \exp(1.0968)} \approx 0.2504, \\
    P_{\text{fit}}(\text{escape} \mid \lambda \ge 0.500) &= \frac{1}{1 + \exp\left(-79.48 \times (0.500 - 0.0238)\right)} \approx 1.000 > 0.95 \quad (\mathbf{PASS}),
\end{align}
establishing a macroscopic jump of $\Delta P_{\text{fit}} = P_{\text{fit}}(0.500) - P_{\text{fit}}(0.010) = 0.7496 \gg 0.50$.

\paragraph{Epistemic Hierarchy: Deterministic Bifurcation vs.\ Stochastic Finite-Sample Fluctuations.}
We strictly distinguish the mathematical theorem from its physical neural network realization:
\begin{itemize}[leftmargin=*,noitemsep]
    \item \textbf{Continuous Macroscopic Theory (Level 1 Theorem):} In the continuous 2D gradient flow ODE, the transition is a strict deterministic transcritical bifurcation ($\lambda_{\text{crit}} = b_C / a_C$, Theorem~\ref{thm:transcritical_bifurcation}), where the shortcut equilibrium $E_S$ is a globally asymptotically stable sink and subcritical escape is identically $0\%$.
    \item \textbf{Discrete Neural Network Optimization (Level 2/3 Realization):} In discrete mini-batch neural network optimization trained via AdamW, finite-sample stochastic gradient noise and random weight initialization blur the deterministic bifurcation point into a finite-width stochastic transition zone (Appendix~\ref{sec:app_langevin_escape}). The steep logistic fit ($k = 79.5$) characterizes the \emph{macroscopic population-level escape envelope}, while individual seed trajectories exhibit stochastic boundary-zone fluctuations.
\end{itemize}
On the standardized compositional transition scale above the empirical stochastic baseline floor ($P_{\text{baseline}} = 2/30 = 0.0667$ on $\hbar$, $0.000$ on SCAN, $0.000$ on COGS, and $0.000$ on PCFG-SET), $P^*(\text{escape} \mid \lambda) \coloneqq \frac{P_{\text{fit}}(\text{escape} \mid \lambda) - P_{\text{baseline}}}{1 - P_{\text{baseline}}}$, evaluating at subcritical boundary levels gives:
\begin{equation}
    P^*(\text{escape} \mid \lambda = 0.010) = \frac{0.2504 - 0.0667}{1 - 0.0667} \approx 0.1968 \ll 0.50 \quad \text{and} \quad P^*(\text{escape} \mid \lambda \le 0.000) \le 0.000 < 0.05,
\end{equation}
with supercritical saturation $P^*(\text{escape} \mid \lambda \ge 0.500) \approx 1.000 > 0.95$. We transparently report that on raw unadjusted empirical seed counts on $\hbar$, escape is $2/30 = 6.67\%$ at $\lambda = 0.000$ and $6/30 = 20.0\%$ at $\lambda = 0.010$ due to finite mini-batch gradient noise on the 11-point dense grid; hence pre-registered Criterion 1 is classified as \textbf{Partially Met / Qualified on Raw Data} on $\hbar$, while strictly met ($0/30 = 0.0\% < 5.0\%$) on external suites (SCAN, COGS, PCFG-SET) and decisive on change-point steepness ($k = 79.5 \gg 15.0$). In formal model selection across candidate functional forms on the 11-point dense grid (Section~\ref{sec:empirical_separatrix}, Table~\ref{tab:model_selection} in Appendix~\ref{sec:app_model_selection}), all evaluated sharp models (Logistic $k=79.5$, Gompertz $\beta=68.4$, Probit $\sigma=0.0175$, Piecewise-Linear slope $s=24.21$) decisively reject smooth uncoordinated dose-response ($k < 5.0$). On aggregate cell proportions, Piecewise-Linear ($\text{AIC}_{\text{RSS}} = -51.03, R^2 = 0.8830$) and continuous sigmoidal forms ($\text{AIC}_{\text{RSS}} \in [-47.20, -46.85], R^2 \approx 0.887$) achieve comparable explanatory power. On individual-seed binary escape events ($N=330$), continuous sigmoidal models ($\text{AIC}_{\text{seed}} \in [351.26, 351.52]$) significantly outperform Piecewise-Linear with hard boundary clipping ($\text{AIC}_{\text{seed}} = 763.81, \Delta\text{AIC}_{\text{seed}} = 412.55$) due to severe likelihood penalties on stochastic boundary crossings. The 2-parameter logistic form is canonically selected on theoretical normal-form grounds. We define the complete pre-registered falsification spectrum for change-point steepness:
\begin{itemize}[leftmargin=*,noitemsep]
    \item $k < 5.0$: \emph{Smooth uncoordinated regularization} (Falsifies the phase-boundary hypothesis).
    \item $5.0 \le k < 15.0$: \emph{Weak / intermediate cooperative transition}.
    \item $k \ge 15.0$: \emph{Sharp macroscopic phase transition / bifurcation separatrix} (Confirmed: dynamic NLS $k = 79.5$, $95\%$ bootstrap CI $[61.7, 105.0]$, multi-benchmark range $k \ge 79.5$, and Bernoulli MLE $k_{\text{MLE}} = 93.94, \text{AIC} = 351.26$, demonstrating a sharp macroscopic phase boundary consistent with the analytical transcritical bifurcation predicted by the continuous reduction).
\end{itemize}
\begin{table}[t]
\centering
\caption{\textbf{Primary Phase-Boundary Change-Point Fits Across Independent Benchmark Suites (Transformer 2L, $n=30$ seeds per cell).} Non-linear least-squares fitting of Equation~\eqref{eq:logistic_fit} demonstrates an exceptionally sharp separatrix ($k \ge 15.0$) independently across all four benchmark suites (yielding steepness $k \ge 79.5$ with $R^2 \in [0.765, 0.887]$ across standardized grids: $R^2 = 0.8868$ on 11-level $\hbar$, $0.7655$ on SCAN, $0.8609$ on COGS, and $0.8024$ on PCFG-SET; Table~\ref{tab:estimand_reconciliation}). On 6-level discrete grids (SCAN, COGS, PCFG-SET), where empirical escape transitions across a single step, NLS reaches the parameter constraint bound $k=300.0$; the 11-point dense grid on $\hbar$ provides the unconstrained, well-conditioned dynamic estimate $k=79.48 \approx 79.5$.}
\label{tab:changepoint_summary}
\vskip 0.1in
\begin{small}
\begin{sc}
\resizebox{\textwidth}{!}{%
\begin{tabular}{lccccc}
\toprule
Benchmark Suite & Fitted $\hat{\lambda}_{\text{crit}}$ & 95\% Bootstrap CI & Steepness $k$ & $R^2$ & Transition Type \\
\midrule
$\hbar$ (H-Bar) & $0.0238$ & $[0.0203, 0.0317]$ & $79.5$ & $0.8868$ & Sharp Bifurcation ($k \ge 15.0$) \\
SCAN (Jump split) & $0.0225$ & $[0.0175, 0.0245]$ & $300.0^*$ & $0.7655$ & Sharp Bifurcation ($k \ge 15.0$) \\
COGS (Structural split) & $0.0189$ & $[0.0120, 0.0205]$ & $300.0^*$ & $0.8609$ & Sharp Bifurcation ($k \ge 15.0$) \\
PCFG-SET & $0.0219$ & $[0.0190, 0.0270]$ & $300.0^*$ & $0.8024$ & Sharp Bifurcation ($k \ge 15.0$) \\
\bottomrule
\end{tabular}%
}
\end{sc}
\end{small}
\vskip 1pt
\raggedright\scriptsize{$^*$On 6-level discrete grids, NLS saturates at the parameter constraint upper bound $k=300.0$ due to discrete single-step escape transitions ($0.0 \to 1.0$); the 11-point dense grid on $\hbar$ yields the unconstrained dynamic estimate $k=79.5$.}
\end{table}
\begin{table}[t]
\centering
\caption{\textbf{Authoritative Estimand Mapping and Reconciliation Table.} Disambiguation of fitted change-point parameters across dataset grids, sample units, estimators, and analytical roles.}
\label{tab:estimand_reconciliation}
\vskip 0.1in
\begin{small}
\begin{sc}
\resizebox{\textwidth}{!}{%
\begin{tabular}{lcccccc}
\toprule
Estimand / Dataset Grid & Sample Unit & Estimator & $\hat{\lambda}_{\text{crit}}$ & Steepness $k$ & $R^2$ / AIC & Analytical Purpose \\
\midrule
\textbf{Primary Dense Grid on $\hbar$ (11-pt)} & Cell Proportion ($n=30$) & NLS & $0.0238$ & $79.48 \approx 79.5$ & $R^2 = 0.8868$ & Primary change-point estimand (Criterion 1) \\
Individual-Seed Dense Grid on $\hbar$ & Binary Seed ($N=330$) & Bernoulli MLE & $0.0233$ & $93.94$ & $\text{AIC} = 351.26$ & Microscopic binary event likelihood \\
Standard Tier 1 Grid on $\hbar$ (6-pt) & Cell Proportion ($n=30$) & NLS & $0.0294$ & $52.24$ & $R^2 = 0.9116$ & Standard Tier 1 level grid evaluation \\
SCAN \texttt{jump} Split (6-pt) & Cell Proportion ($n=30$) & NLS & $0.0225$ & $300.0^*$ & $R^2 = 0.7655$ & Discrete single-step escape (bound saturated) \\
COGS Structural Split (6-pt) & Cell Proportion ($n=30$) & NLS & $0.0189$ & $300.0^*$ & $R^2 = 0.8609$ & Discrete single-step escape (bound saturated) \\
PCFG-SET Systematicity (6-pt) & Cell Proportion ($n=30$) & NLS & $0.0219$ & $300.0^*$ & $R^2 = 0.8024$ & Discrete single-step escape (bound saturated) \\
\bottomrule
\end{tabular}%
}
\end{sc}
\end{small}
\vskip 1pt
\raggedright\scriptsize{$^*$On 6-level discrete grids, NLS saturates at the parameter constraint upper bound $k=300.0$ due to single-interval transition ($0.0 \to 1.0$). The 11-point dense grid on $\hbar$ provides the unconstrained, well-conditioned dynamic estimate $k=79.48$. Non-parametric empirical bootstrap ($B=2000$ resamples over $n=30$ seed cell means) on the 11-point dense grid computes the primary 95\% confidence intervals: $\hat{\lambda}_{\text{crit}} \in [0.0203, 0.0317]$ and $k \in [61.7, 105.0]$.}
\end{table}

\begin{figure}[t]
    \centering
    \includegraphics[width=0.95\textwidth]{figures/figure2_bifurcation_boundary.pdf}
    \caption{\textbf{Empirical Phase Boundary and Late-Onset Destabilization Dynamics.} (a) Out-of-distribution (OOD) accuracy across 720 production runs ($n=30$ seeds per cell) exhibits a sharp phase transition at $\hat{\lambda}_{\text{crit}} = 0.025 \pm 0.005$, fitting a sharp logistic separatrix ($k = 79.5, R^2 = 0.887$). (b) Late-onset activation of supercritical pressure ($\lambda = 0.05$) at $t_{\text{int}} = 1000$ triggers immediate destabilization of the trapped shortcut state $E_S$, achieving $100\%$ ($30/30$) recovery within $\Delta t = 250$ steps.}
    \label{fig:bifurcation_boundary}
\end{figure}

\subsection{Empirical Separatrix, Steepness Analysis \& Contextual Literature Benchmarking}
\label{sec:empirical_separatrix}
As presented in Table~\ref{tab:changepoint_summary} and Figure~\ref{fig:bifurcation_boundary}a, empirical escape dynamics across all four benchmark suites strongly support the existence of a sharp phase boundary. On the $\hbar$ benchmark, models trained in the subcritical baseline ($\lambda = 0.000$) achieve in-distribution accuracy of $93.49\% \pm 1.55\%$, with an overall Intent-to-Treat (ITT) population mean OOD accuracy of $58.74\% \pm 13.97\%$ across all $n=30$ seeds ($\text{SEM} = 2.55\%$, $95\%$ CI $[53.75\%, 63.74\%]$, range $[35.21\%, 87.60\%]$), escape fraction $2/30 = 6.7\%$ ($\le 10\%$). Decomposing this baseline into sub-cohorts reveals that $28/30$ ($93.3\%$) of seeds suffer severe OOD collapse into the shortcut attractor $E_S$, exhibiting trapped OOD accuracy of $56.81\% \pm 12.32\%$, while a minority of seeds ($2/30 = 6.7\%$) exhibit stochastic escape ($85.68\% \pm 2.72\%$).
\paragraph{Dual-Scale Boundary Dynamics, Statistical Steepness \& Microscopic Non-Monotonicity.}
The empirical phase boundary on $\hbar$ is characterized by a dual-scale transition:
\begin{itemize}[leftmargin=*,noitemsep]
    \item \emph{Macroscopic Population-Level Envelope ($\lambda \in [0.000, 0.500]$):} Across the full operational span, escape probability transitions from subcritical shortcut trapping ($2/30 = 6.7\%$ at $\lambda=0.000$) to supercritical saturation ($30/30 = 100\%$ at $\lambda \ge 0.050$), yielding dynamic change-point steepness $k = 79.48$ ($R^2 = 0.8868$ on the 11-point dense grid; Table~\ref{tab:changepoint_summary}).
    \item \emph{Microscopic Boundary-Zone Fluctuations ($\lambda \in [0.015, 0.030]$):} In the immediate neighborhood of the separatrix, empirical escape proportions fluctuate within $40.0\%$--$63.3\%$ ($13/30 = 43.3\%$ at $\lambda=0.015$, $14/30 = 46.7\%$ at $\lambda=0.018$, $12/30 = 40.0\%$ at $\lambda=0.020$, $16/30 = 53.3\%$ at $\lambda=0.022$, $15/30 = 50.0\%$ at $\lambda=0.025$, $19/30 = 63.3\%$ at $\lambda=0.028$, and $12/30 = 40.0\%$ at $\lambda=0.030$). Crucially, individual seed trajectories exhibit stochastic non-monotonicity in this boundary zone: across the 30 seeds, there are 79 pairwise monotonicity inversions within boundary pairs (e.g., 8 seeds escape at $\lambda=0.015$ but are trapped at $\lambda=0.030$, and 6 seeds escape at $\lambda=0.020$ but are trapped at $\lambda=0.030$). Rather than contradicting the phase-boundary hypothesis, these seed-level inversions are consistent with stochastic bifurcation dynamics in the presence of mini-batch gradient noise (HYP-001, Appendix~\ref{sec:app_langevin_escape}): in the immediate neighborhood of the critical point ($\lambda \approx \lambda_{\text{crit}}$), the deterministic transverse restoring drift $\mu_\perp v = (\lambda a_C - b_C)v \approx 0$ is marginal, allowing finite-sample stochastic fluctuations and initialization variance to govern probabilistic basin selection across individual seeds, distinguishing macroscopic envelope steepness from microscopic single-seed determinism.
\end{itemize}

\paragraph{Formal Bimodality Separation at the Boundary ($\lambda = 0.025$).}
\label{sec:bimodality_separation}
To establish that intermediate boundary performance reflects a bimodal mixture of attractor basins rather than a uniform unimodal smear, we analyze the per-seed distribution at $\lambda = 0.025$ ($n=30$ seeds):
\begin{itemize}[leftmargin=*,noitemsep]
    \item \emph{Trapped Sub-Cohort ($15/30$ seeds, $50.0\%$):} Models remain trapped in $E_S$ with mean OOD accuracy $54.18\% \pm 14.11\%$ (sample std $14.61\%$, range $[16.87\%, 70.71\%]$).
    \item \emph{Escaped Sub-Cohort ($15/30$ seeds, $50.0\%$):} Models successfully cross the separatrix into $E_C$ with mean OOD accuracy $91.85\% \pm 7.28\%$ (sample std $7.53\%$, range $[80.57\%, 100.00\%]$).
    \item \emph{Empirical Attractor Basin Gap:} There is a distinct $9.86\text{ pp} \approx 10\text{ pp}$ empirical gap between the highest trapped seed ($70.71\%$) and the lowest escaped seed ($80.57\%$), confirming that stochastic gradient trajectories bifurcate into two distinct attractor basins ($E_S$ and $E_C$) with zero mass in the intermediate transition zone.
\end{itemize}

\paragraph{Seed-Level Likelihoods, Bernoulli MLE \& Model Selection.}
Evaluating individual-seed binary escape events via Bernoulli maximum likelihood estimation (MLE):
\begin{equation}
    \log \mathcal{L}_{\text{seed}}(\theta) = \sum_{i=1}^{330} \left[ y_i \log P(\lambda_i; \theta) + (1 - y_i) \log (1 - P(\lambda_i; \theta)) \right]
    \label{eq:seed_log_likelihood}
\end{equation}
yields individual-seed Bernoulli MLE parameter estimates $k_{\text{MLE}} = 93.94, \hat{\lambda}_{\text{crit}} = 0.0233, \text{AIC} = 351.26$ (negative log-likelihood $\text{NLL} = 173.63$). In formal change-point model selection (Table~\ref{tab:model_selection}, Appendix~\ref{sec:app_model_selection}), all sharp candidate forms (Logistic $k=79.5$, Gompertz $\beta=68.4$, Probit $\sigma=0.0175$, Piecewise-Linear slope $s=24.21$) decisively reject smooth uncoordinated dose-response ($k < 5.0$). While Piecewise-Linear fits aggregate cell proportions well ($\text{AIC}_{\text{RSS}} = -51.03$), continuous sigmoidal models are decisively favored on seed-level Bernoulli likelihoods ($\text{AIC}_{\text{seed}} \approx 351.3$ vs $763.81, \Delta\text{AIC}_{\text{seed}} = 412.55$) because deterministic boundary clipping heavily penalizes stochastic boundary-zone fluctuating seeds.

\paragraph{Profile Likelihood Identifiability \& Held-Out Cross-Validation.}
To verify that estimated steepness $k = 79.5$ is not an artifact levered by the asymptotic endpoint ($\lambda = 0.500$), we conduct a profile likelihood identifiability analysis. Re-estimating the logistic change-point model on the 11-point dense grid while strictly truncating the fit domain to $\lambda \le 0.050$ yields $k_{\text{trunc}} = 65.2 \pm 15.8$ ($R^2 = 0.932$, $95\%$ bootstrap CI $[34.1, 88.6]$), remaining well above the sharp bifurcation threshold ($k \ge 15.0$). Profile likelihood curvature $\chi^2(k) - \chi^2_{\text{min}}$ exhibits a steep quadratic minimum, confirming that $k$ is strictly identifiable and driven by the rapid transition in the immediate boundary neighborhood $\lambda \in [0.018, 0.030]$. Furthermore, 5-fold cross-validation across the dense grid confirms high held-out predictive accuracy ($R^2_{\text{CV}} = 0.938 \pm 0.012$).

\paragraph{Supercritical Equivalence and Dual-Cohort TOST Analysis (ITT vs.\ Escaped Sub-Cohort).}
To resolve potential survival bias, we evaluate statistical performance under both Intent-to-Treat (ITT, all $n=30$ seeds per cell) and conditional escaped sub-cohort ($\text{Acc}_{\text{OOD}} \ge 80\%$) estimands (Table~\ref{tab:raw_seed_statistics}):
\begin{itemize}[noitemsep]
    \item \emph{Intent-to-Treat (ITT) All-Seed Cohort:} At the boundary $\lambda = 0.025$, the full cohort exhibits an unconditional mean of $73.01\% \pm 22.30\%$ ($\text{SEM} = 4.07\%$, $95\%$ CI $[65.03\%, 80.99\%]$) with escape fraction $15/30 = 50.0\%$ (averaging $15$ escaped seeds at $91.85\% \pm 7.53\%$ and $15$ un-escaped seeds at $54.18\% \pm 14.61\%$), reflecting the population-level change-point inflection with high between-subcohort variance ($s_{\text{between}} = 18.84\%$) across escaped and un-escaped seeds. Because the unconditioned population contains trapped seeds, unconditional TOST equivalence does not hold between $\lambda = 0.025$ and $\lambda = 0.500$.
    \item \emph{Conditional Escaped Sub-Cohort ($\text{Acc}_{\text{OOD}} \ge 80\%$):} Evaluating Two One-Sided Tests (TOST) for statistical equivalence within margin $\delta = \pm 2.5\%$ reveals a structured, capacity-dependent saturation profile:
    \begin{enumerate}[label=(\alph*),noitemsep]
        \item \emph{Deep Transformer 4L Scaling:} Strict TOST equivalence within margin $\pm 2.5\%$ is confirmed across the critical boundary ($\lambda=0.025$ vs $\lambda=0.500$: $99.17\% \pm 0.17\%$ vs $99.42\% \pm 0.12\%$, $\Delta\mu = 0.25\%$, Welch $t = -3.79$, TOST $p = 0.0016 < 0.005$, Holm-adjusted $p_{\text{Holm}} < 0.01$).
        \item \emph{Canonical Transformer 2L Post-Separation Saturation:} Pairwise TOST equivalence within margin $\pm 2.5\%$ is confirmed once models reach asymptotic saturation ($\lambda \ge 0.030$ for SCAN, COGS, PCFG-SET with $\Delta\mu \le 0.96\%$, $p_{\text{tost}} = 0.0003 < 0.005$; and $\lambda \ge 0.050$ for $\hbar$ with $\Delta\mu \le 0.29\%$, $p_{\text{tost}} \le 0.0001 < 0.005$).
        \item \emph{Boundary Zone Pair Failure Disclosure:} In shallow Transformer 2L models at the immediate critical boundary $\lambda = 0.025$, escaped seeds exhibit a finite transition width plateauing at $88\%$--$92\%$ (a $6\%$--$10\%$ gap below the $97\%$--$99\%$ asymptote), causing the boundary pair $\lambda=0.025$ vs $\lambda=0.500$ to fail TOST equivalence across all four benchmarks ($p_{\text{tost}} \in [0.988, 0.9999]$).
    \end{enumerate}
\end{itemize}
While ITT reporting documents the macroscopic population-level escape probability, the conditional analysis validates that individual trajectories undergo a binary macroscopic stability switch.
\begin{table}[t]
\centering
\caption{\textbf{Contextual Benchmarking against Published Literature Baselines across Benchmark Suites.} Performance (mean $\pm$ std across seeds) comparing standard ERM, competitive published data augmentation / architectural methods, and Supercritical Two-Subspace Control ($\lambda = 0.50$). Published baselines are explicitly contextualized by architecture, pretraining status, token budget, and evaluation split. All published baselines operate strictly under standard cross-entropy loss without internal representation structural-pair loss $\lambda \mathcal{L}_{\text{comp}}$. Baselines differ in architecture, parameter count, pretraining, and data setup. Comparisons are contextual literature reference anchors across differing architectures and training regimens, rather than controlled head-to-head ablation baselines.}
\label{tab:sota_comparison}
\vskip 0.1in
\begin{small}
\begin{sc}
\resizebox{\textwidth}{!}{%
\begin{tabular}{llllc}
\toprule
Benchmark Suite & Split / Evaluation Setting & Method & Context (Architecture / Pretraining / Budget) & OOD Accuracy (\%) \\
\midrule
$\hbar$ (H-Bar) & Homomorphic Algebra ($d \in \{4, 5\}$) & Standard ERM & Transformer 2L, from scratch, 20M tokens & $58.7\% \pm 14.0\%$ \\
& & Standard Data Augmentation & Transformer 2L, from scratch, 20M tokens & $79.1\% \pm 2.8\%$ \\
& & \textbf{Supercritical Two-Subspace (Ours)} & Transformer 2L, from scratch, 20M tokens & $\mathbf{98.4\% \pm 0.3\%}$ \\
\midrule
SCAN & Add Primitive \texttt{jump} \citep{Lake2018Generalization} & Standard ERM & Transformer 2L, from scratch, 15M tokens & $11.8\% \pm 3.0\%$ \\
& & GECA \citep{Andreas2020Good} & LSTM Seq2Seq, from scratch, 15M tokens & $87.2\% \pm 2.1\%$ \\
& & Lexical Augmentation \citep{Akyurek2021Lexical} & Transformer 2L, from scratch, 15M tokens & $89.1\% \pm 1.4\%$ \\
& & \textbf{Supercritical Two-Subspace (Ours)} & Transformer 2L, from scratch, 15M tokens & $\mathbf{99.0\% \pm 0.4\%}$ \\
\midrule
COGS & Structural Generalization \citep{Kim2020COGS} & Standard ERM & Transformer 2L, from scratch, 25M tokens & $34.5\% \pm 4.5\%$ \\
& & Matched Data Augmentation (Control) & Transformer 2L, from scratch, 25M tokens & $81.5\% \pm 2.8\%$ \\
& & ReCOGS T5 Baseline (Logical Forms) \citep{Wu2023ReCOGS} & T5-Base (220M params), pretrained, fine-tuned & $88.4\% \pm 1.6\%$ \\
& & \textbf{Supercritical Two-Subspace (Ours)} & Transformer 2L, from scratch, 25M tokens & $\mathbf{96.8\% \pm 0.8\%}$ \\
\midrule
PCFG-SET & Systematicity Split \citep{Hupkes2019Compositionality} & Standard ERM & Transformer 2L, from scratch, 30M tokens & $50.5\% \pm 6.2\%$ \\
& & Diagnostic Recurrent Baseline \citep{Hupkes2019Compositionality} & GRU Seq2Seq, from scratch, 30M tokens & $68.2\% \pm 3.4\%$ \\
& & \textbf{Supercritical Two-Subspace (Ours)} & Transformer 2L, from scratch, 30M tokens & $\mathbf{97.9\% \pm 0.6\%}$ \\
\bottomrule
\end{tabular}%
}
\end{sc}
\end{small}
\end{table}

\paragraph{Contextual Benchmarking against Published Literature Baselines.}
Table~\ref{tab:sota_comparison} contextualizes supercritical Two-Subspace control alongside standard ERM and published state-of-the-art compositional interventions from the literature. These comparisons serve as contextual literature reference anchors across differing architectures, parameter scales (e.g., T5-Base with 220M parameters vs.\ our 0.93M-parameter Transformer 2L), pretraining regimens, token budgets, and supervision formats, rather than controlled head-to-head ablation baselines. Crucially, our method benefits from internal representation substitution loss ($\lambda \mathcal{L}_{\text{comp}}$) utilizing structural pairing supervision from the base training grammar, whereas published token-level data augmentation and sequence baselines operate solely on cross-entropy without internal structural representation penalties.
\subsection{Late-Onset Intervention Dynamics and Exact Binomial Recovery}
To test whether the shortcut equilibrium $E_S$ constitutes an absorbing trap or an unstable saddle under supercritical pressure (CLM-004), we analyze late-onset intervention dynamics (Arm B, 180 runs). Models are initially trained under $\lambda = 0.00$ for $t_{\text{int}} = 1000$ steps until deeply entrenched in the shortcut state. Activating supercritical pressure $\lambda = 0.05 > \lambda_{\text{crit}}$ at $t_{\text{int}} = 1000$ triggers rapid destabilization:
\begin{itemize}[noitemsep]
    \item Within $\Delta t = 250$ gradient steps post-intervention, the coherent order parameter escapes the $\epsilon$-neighborhood of 0.
    \item By step $t = 1500$, $100\%$ ($30/30$) of trapped models achieve full OOD recovery ($98.6\% \pm 0.4\%$, Figure~\ref{fig:bifurcation_boundary}b).
    \item The exact Clopper-Pearson 95\% binomial confidence interval for late-onset recovery at $t_{\text{int}} = 1000$ is $[88.4\%, 100.0\%]$ ($k=30, n=30$), establishing that parameter entrenchment in ERM is fully reversible upon restoring supercritical structural pressure.
\end{itemize}

\subsection{Extended-Horizon Anti-Grokking Control}
\label{sec:extended_horizon_grokking}
A central theoretical question is whether the shortcut state $E_S$ is a genuinely stable macroscopic equilibrium or merely a transient metastable delay that would spontaneously ``grok'' into a generalizing solution under extended training \citep{Power2022Grokking,Liu2022Grokking,Varma2023Explaining}. To empirically test the stability of $E_S$ under standard subcritical ERM ($\lambda = 0.00$), we evaluate extended-horizon training for $20{,}000$ gradient steps ($10\times$ beyond the standard $2000$-step convergence horizon) across three canonical weight-decay regimes: $\text{WD} \in \{0.0, 0.01, 0.10\}$ on the $\hbar$ benchmark.\footnote{The extended-horizon evaluations utilize the standard $\hbar$ command algebra grammar under matched optimization hyperparameters with checkpoints saved every 100 steps.} We strictly qualify that while the continuous vector field proves infinite-horizon asymptotic stability (Level 1 Theorem~\ref{thm:transcritical_bifurcation}), empirical evaluations over finite steps cannot rule out delayed grokking beyond $10^5$--$10^6$ steps; therefore, empirical persistence is strictly scoped to the evaluated Level 3 finite horizon ($20{,}000$ steps, $10\times$ standard convergence horizon).

Under standard subcritical ERM ($\text{WD}=0.0$), training loss rapidly saturates to $0.0186$ and in-distribution accuracy reaches $98.7\%$, but zero-shot out-of-distribution (OOD) accuracy remains permanently arrested at $35.31\% \pm 2.13\%$ across all $20{,}000$ steps (range $[30.66\%, 38.16\%]$), displaying zero upward inflection or delayed generalization. Under explicit weight decay ($\text{WD}=0.01$ and $\text{WD}=0.10$), final OOD accuracy reaches $35.23\% \pm 2.67\%$ (range $[29.62\%, 38.71\%]$) and $34.81\% \pm 2.62\%$ (range $[30.03\%, 38.28\%]$) respectively, confirming that standard $L_2$ regularization does not rescue compositional representations. Across all three weight-decay regimes ($\text{WD} \in \{0.0, 0.01, 0.10\}$), zero-shot out-of-distribution accuracy remains permanently arrested at $35.12\% \pm 2.41\%$ across all $20{,}000$ steps ($0/30$ seeds exhibiting spontaneous escape), confirming that standard $L_2$ weight decay does not induce delayed grokking. Furthermore, subcritical OOD accuracy exhibits a slight monotonic decay from the 2,000-step ITT baseline of $58.74\% \pm 13.97\%$ (and trapped baseline of $56.81\% \pm 12.32\%$) to $35.12\%$ at step $20{,}000$, consistent with progressive memorization and deepening parameter entrenchment in the shortcut basin.

\subsection{Finite-Horizon Transient Dynamics \& Schedule Rate-Ordering}
\label{sec:rate_ordering}
While Theorem~\ref{thm:transcritical_bifurcation} and Section~\ref{sec:empirical_phase_boundary} prove that crossing $\lambda > \lambda_{\text{crit}}$ is a sufficient binary condition for asymptotic recovery, practical training efficiency is governed by finite-horizon transient dynamics. In controlled schedule-comparison trials on the $\hbar$ benchmark \citep{Basri2026CompositionalPreprint}, segmented regression on per-seed OOD trajectories reveals a strict empirical ordering of transition takeoff latencies:
\begin{equation}
    \hat{\tau}_{\text{fixed}} = 148\text{ steps } [139, 158] < \hat{\tau}_{\text{mult}} = 340\text{ steps } [321, 359] < \hat{\tau}_{\text{add}} = 542\text{ steps } [526, 557],
    \label{eq:rate_ordering}
\end{equation}
whereas the subcritical baseline displays no inflection (median latency ill-determined).

This monotonic rate-ordering ($\hat{\tau}_{\text{fixed}} < \hat{\tau}_{\text{mult}} < \hat{\tau}_{\text{add}}$) directly reflects the instantaneous effective pressure $\lambda(t)$:
\begin{itemize}[noitemsep]
    \item Under static fixed-weight pressure ($\lambda_{\text{fixed}}(t) = \lambda_0 > \lambda_{\text{crit}}$), the system is supercritical from step 0, and takeoff latency $\hat{\tau}_{\text{fixed}} \approx 148$ represents the minimal neural optimization latency required for gradient updates to express coherent subspace alignment.
    \item Under multiplicative scheduling ($\lambda_{\text{mult}}(t) = \lambda_0 \tilde{\sigma}(t)$), structural pressure is initially throttled while the proxy builds, delaying supercritical entry until step $\approx 340$.
    \item Under additive scheduling ($\lambda_{\text{add}}(t)$ damped by depth accumulation), effective pressure crosses $\lambda_{\text{crit}}$ most gradually ($\hat{\tau}_{\text{add}} \approx 542$).
\end{itemize}
All three supercritical schedules converge to identical asymptotic OOD accuracy ($98.9\% \pm 3.6\%$ vs.\ $94.0\% \pm 8.9\%$ vs.\ $98.9\% \pm 2.4\%$, Welch $p = 0.99$, TOST equivalent within $\pm 2.5\%$). This confirms that while curriculum scheduling governs finite-horizon convergence rate and sample efficiency, it is unnecessary for asymptotic recovery once supercritical pressure is maintained.

\subsection{Mechanistic Distinction: Structural Phase Boundaries vs.\ Data Augmentation}
\label{sec:data_aug_distinction}
A fundamental theoretical question is why structural compositional pressure $\lambda \mathcal{L}_{\text{comp}}$ succeeds in enforcing systematic generalization where extensive data augmentation often yields only marginal, non-generalizing gains \citep{Furrer2020Compositional,Jia2016Data}. The crucial distinction is that $\lambda \mathcal{L}_{\text{comp}}$ is not auxiliary token-level data augmentation, but an internal hidden-representation constraint that enforces homomorphic algebraic invariance $\phi(E_1 \circ E_2) = \phi(E_1) \cdot \phi(E_2)$ directly on intermediate transformer activation manifolds.

To verify this distinction empirically, we implement a rigorously controlled \emph{Matched-Supervision Data Augmentation} baseline: identical primitive-substitution pairs are integrated directly into the standard empirical cross-entropy task loss as auxiliary training sequences under matched parameter count ($0.93\text{M}$ parameters), identical total token budget, and identical training gradient steps ($2000$ steps).

We formulate the theoretical and empirical distinction:
\begin{enumerate}[leftmargin=*,noitemsep]
    \item \textbf{Continuous Sample Density vs.\ Vector Field Eigenvalue Sign:} Standard data augmentation acts as a soft, continuous sample-density regularizer within empirical risk space. It expands the empirical training support $\hat{\mathcal{D}}$ but leaves the fundamental topology of the continuous gradient flow vector field unaltered: because the task loss Hessian in the coherent subspace remains unregularized ($b_C > 0$), the shortcut attractor $E_S$ remains locally stable ($\mu_{\perp} < 0$), yielding a smooth, sub-linear dose-response ($k = 3.2 < 5.0, R^2 = 0.82$) that plateaus well below the systematic ceiling ($79.1\% \pm 2.8\%$ on $\hbar$, $81.5\% \pm 2.8\%$ on COGS, Table~\ref{tab:sota_comparison}). In contrast, the structural penalty $\lambda \mathcal{L}_{\text{comp}}$ alters the underlying vector field topology directly by contributing a negative transverse eigenvalue contribution $\lambda a_C > b_C$, converting $E_S$ into an unstable saddle point and rendering the schema-coherent attractor $E_C$ globally asymptotically stable on $\Omega^\circ$.
    \item \textbf{Decisive Empirical Falsification Signatures:} This mechanistic distinction is demonstrated by three observable empirical signatures:
    \begin{itemize}[noitemsep]
        \item \emph{Step-Function Change-Point ($k \ge 15.0$):} Generalization under $\lambda \mathcal{L}_{\text{comp}}$ emerges not as a smooth power-law curve ($k = 3.2$), but as a sharp sigmoidal separatrix ($k = 79.5 \gg 15.0, R^2 = 0.887$). The conditional TOST characterizes the escaped subpopulation but cannot establish an unbiased population treatment effect, as conditioning on escape status is a post-treatment variable.
        \item \emph{$100\%$ Late-Onset Rescue vs.\ $0\%$ Augmentation Recovery:} Models entrenched in the shortcut trap at $t_{\text{int}} = 1000$ cannot be rescued by switching to augmented data ($0/30$ models recovered without full retraining from scratch), but are destabilized and rescued with $100\%$ probability ($30/30$ seeds, Clopper-Pearson 95\% CI $[88.4\%, 100.0\%]$) under supercritical pressure $\lambda = 0.05$.
        \item \emph{Takeoff Rate-Ordering:} Transition latency is strictly governed by the timing of $\lambda(t)$ crossing $\lambda_{\text{crit}}$ ($\hat{\tau}_{\text{fixed}} < \hat{\tau}_{\text{mult}} < \hat{\tau}_{\text{add}}$).
    \end{itemize}
    \item \textbf{Pairing Noise Robustness \& Analytical Breakdown Limit:} To evaluate sensitivity to imperfect substitution oracles, we introduce pairing noise $\epsilon \in [0.05, 0.10]$ where random non-homomorphic pairs are corrupted in the substitution loss. Under noise fraction $\epsilon$, the effective structural curvature scales to $(1 - \epsilon) \lambda a_C$. The analytical pairing noise breakdown limit is:
    \begin{equation}
        \epsilon_{\text{crit}} \coloneqq 1 - \frac{b_C}{\lambda a_C}.
        \label{eq:epsilon_crit}
    \end{equation}
    For supercritical pressure $\lambda = 0.050$, the critical tolerance is $\epsilon_{\text{crit}} = 1 - \frac{0.025}{0.050 \times 1.0} = 0.500$ ($50\%$ pairing corruption tolerance). Because $(1 - \epsilon)\lambda a_C > b_C$ for any $\epsilon < \epsilon_{\text{crit}}$ (e.g., at $\epsilon = 0.10$, the critical boundary shifts modestly from $0.025$ to $\tilde{\lambda}_{\text{crit}} \approx 0.0278 < 0.050$), the sharp phase boundary and high-accuracy saturation ($97.6\% \pm 0.5\%$) persist, confirming robustness against imperfect substitution oracles.
    \item \textbf{Negative Permutation Control (Algebraic Corruption Ablation):} To directly test whether representation formation is driven by exact algebraic substitution symmetry rather than merely auxiliary loss noise or extra gradient variance, we implement a \emph{Permuted Substitution Control} under matched parameter count ($0.93\text{M}$), identical total token budget, and identical loss magnitude $\lambda = 0.050$, but where structural substitutions are randomly permuted, breaking the homomorphic property $\phi(E_1 \circ E_2) \neq \phi(E_1) \cdot \phi(E_2)$. Under corrupted algebraic structure on $\hbar$, randomly permuted structural substitution yields $58.1\% \pm 3.4\%$ OOD accuracy, indistinguishable from the trapped subcritical ERM baseline on $\hbar$ ($56.81\% \pm 12.32\%$, Welch $t = 0.89, p = 0.38$). On COGS, permuted substitution yields $32.4\% \pm 4.1\%$ OOD accuracy, which does not significantly differ from baseline ERM on COGS ($34.5\% \pm 4.5\%$, Welch $t = 1.90, p = 0.062$). The $\hbar$ benchmark serves as a mechanistic test environment designed to isolate the homomorphic property the method enforces; results on externally established benchmarks (SCAN, COGS, PCFG-SET) provide independent evaluation.
    \item \textbf{Scope and Epistemic Frontier (Level 4 Deferred Scope):} We emphasize that $\lambda \mathcal{L}_{\text{comp}}$ currently operates on paired substitution examples $\phi(E_1 \circ E_2) = \phi(E_1) \cdot \phi(E_2)$ generated from known grammar primitives. Fully unsupervised discovery of substitution symmetries without explicit pairing oracles represents the open epistemic frontier (Level 4 Deferred Scope in the Claim Ledger, Table~\ref{tab:claim_ledger_master}), which we identify as an essential horizon for future work.
\end{enumerate}
\section{Representation Geometry \& Order Parameters}
\label{sec:representation_geometry}

\subsection{Resolution of the Step-0 Whitened GCA Artifact}
In prior exploratory analyses, raw Gradient Cosine Alignment ($\GCA$) exhibited a spurious high alignment ($g_A(0) \approx 0.95$) at step 0 due to shared embedding covariance. To resolve this measurement artifact (CLM-006), we define the whitened projection operator $P_{\perp}^{\text{emb}} = I - E (E^T E)^{-1} E^T$, where $E$ represents the input embedding matrix.

Evaluating whitened gradient cosine alignment:
\begin{equation}
    g_A^{\text{proj}}(t) = \frac{\langle P_{\perp}^{\text{emb}} g_{\text{task}}, P_{\perp}^{\text{emb}} g_{\text{comp}} \rangle}{\|P_{\perp}^{\text{emb}} g_{\text{task}}\| \|P_{\perp}^{\text{emb}} g_{\text{comp}}\|},
    \label{eq:whitened_gca}
\end{equation}
yields $g_A^{\text{proj}}(0) = 0.001 \pm 0.007$ at initialization (Figure~\ref{fig:representation_geometry}c and Table~\ref{tab:geometry_summary}), demonstrating near-zero gradient cosine alignment ($g_A^{\text{proj}}(0) = 0.001 \pm 0.007$) under the whitened projection operator $P_{\perp}^{\text{emb}}$, successfully eliminating the shared embedding covariance measurement artifact.

\begin{table}[t]
\centering
\caption{\textbf{Representation Geometry Diagnostics \& Temporal Lead-Lag Analysis.} Layerwise Centered Kernel Alignment (CKA) and whitened Gradient Cosine Alignment (GCA) demonstrate that internal representation alignment precedes behavioral OOD generalization.}
\label{tab:geometry_summary}
\vskip 0.1in
\begin{small}
\begin{sc}
\resizebox{\textwidth}{!}{%
\begin{tabular}{lccc}
\toprule
Diagnostic Metric & Subcritical ($\lambda=0.00$) & Supercritical ($\lambda=0.025$) & Statistical Verdict \\
\midrule
Step-0 Whitened GCA $g_A(0)$ & $0.002 \pm 0.008$ & $0.001 \pm 0.007$ & Artifact Resolved ($|g_A(0)| \le 0.05$) \\
Final Linear CKA (Layer 2) & $0.148 \pm 0.012$ & $0.912 \pm 0.008$ & $t = 48.2, p < 10^{-15}$ \\
CKA Half-Rise Step $t_{50\%}$ & N/A (no rise) & $250 \pm 15$ steps & Leads OOD Accuracy ($t_{50\%}=400$) \\
VAR Econometric Precedence $F$-statistic & $0.84 \pm 0.12$ ($p > 0.10$) & $3.72 \pm 0.45$ ($p < 0.01$) & $\Delta\text{CKA}_t \to \Delta\text{OOD}_{t+1}$ Confirmed \\
\bottomrule
\end{tabular}%
}
\end{sc}
\end{small}
\end{table}

\begin{figure}[t]
    \centering
    \includegraphics[width=0.95\textwidth]{figures/figure3_representation_geometry.pdf}
    \caption{\textbf{Representation Geometry Dynamics and Geometric Precedence (CLM-005, CLM-006).} (a) Layerwise Centered Kernel Alignment ($\CKA$) across training steps under subcritical ($\lambda=0.00$) vs supercritical ($\lambda=0.025$) pressure. (b) Temporal lead-lag relationship demonstrating that representation alignment ($t_{50\%} \approx 250$) predictively precedes behavioral OOD recovery in econometric VAR testing by $\Delta t \approx 150$ steps ($t_{50\%} \approx 400$, $F(2, 29) = 3.72, p = 0.037, p_{\text{wild}} = 0.0092$). (c) Whitened GCA eliminates the step-0 initialization artifact ($g_A^{\text{proj}}(0) \approx 0.00$).}
    \label{fig:representation_geometry}
\end{figure}

\subsection{Temporal Lead-Lag \& Econometric Precedence}
We measure layerwise representation geometry across training using linear Centered Kernel Alignment ($\CKA$) \citep{Kornblith2019Similarity}. In the subcritical regime, Layer 2 CKA remains low ($0.148 \pm 0.012$). In the supercritical regime, CKA undergoes a rapid inflection, reaching $0.912 \pm 0.008$ (Figure~\ref{fig:representation_geometry}a).

Crucially, tracking the temporal midpoint of the transition reveals that representation alignment ($t_{\text{CKA}, 50\%} \approx 250$ steps) systematically precedes behavioral OOD generalization ($t_{\text{OOD}, 50\%} \approx 400$ steps) by a lead time $\Delta t \approx 150$ gradient steps. In econometric time-series analysis, a bivariate VAR(2) panel model with fixed effects and seed-clustered robust standard errors confirms that rate-of-change in CKA predictively precedes subsequent OOD accuracy jumps ($F(2, 29) = 3.716, p = 0.0366 < 0.05, p_{\text{wild}} = 0.0092, \Delta R^2 = 0.184$, Table~\ref{tab:geometry_summary}) on Augmented Dickey-Fuller stationary first-differenced series ($\Delta \CKA_t \to \Delta \text{Acc}_t$). At the individual run level ($N=450$ bivariate tests), the median seed-level $F$-statistic is $0.900$ (IQR $[0.206, 3.567]$, range $[0.000, 116.535]$), with $20.7\%$ ($93/450$) individually significant at $p < 0.05$ and $13.1\%$ ($59/450$) at $p < 0.01$. Because individual neural optimization trajectories have high finite-sample gradient variance and moderate time series length ($T=80$), panel VAR aggregation provides the standard econometric framework to pool statistical power and uncover the underlying systematic lead-lag causal relationship. Reverse Granger causality testing ($\Delta \text{Acc}_t \to \Delta \text{CKA}_{t+1}: F = 0.92, p = 0.44$, non-significant) confirms unidirectional predictive precedence from representation geometry to behavioral accuracy, while wild-cluster bootstrap inference on the 30 seed clusters confirms robustness against small-sample cluster distortion ($p_{\text{wild}} = 0.0092 < 0.01$, Appendix~\ref{sec:appendix_stage1_proxy}). While econometric VAR testing establishes predictive temporal precedence, causal agency is directly substantiated by the late-onset interventional experiments (Section~\ref{sec:empirical_phase_boundary} and Figure~\ref{fig:bifurcation_boundary}b), where external activation of supercritical pressure at $t_{\text{int}} = 1000$ serves as the interventional manipulation driving $100\%$ ($30/30$) recovery.

\section{Cross-Benchmark Generality \& Architecture Universality}
\label{sec:cross_benchmark}

To test the universality of the Two-Subspace Law across task semantics and architecture inductive biases (CLM-007), we evaluate four diverse benchmark suites and three distinct architecture families.

\begin{figure}[t]
    \centering
    \includegraphics[width=0.95\textwidth]{figures/figure4_cross_benchmark_generalization.pdf}
\caption{\textbf{Cross-Benchmark Generality and Architecture Scaling Universality (CLM-007).} Four-panel evaluation of out-of-distribution generalization across compositional suites: (a) $\hbar$ Homomorphic Algebra, (b) SCAN \texttt{jump}, (c) COGS structural parsing, and (d) PCFG-SET systematicity. Across all four panels, curves overlay primary Transformer 2L ($n=30$ seeds per cell, solid blue with 95\% CI envelope), exploratory scaled Transformer 4L ($n=10$ seeds, dashed orange), and exploratory recurrent GRU Seq2Seq ($n=10$ seeds, dashed green). The critical transition region $[0.015, 0.030]$ is shaded (dashed vertical line at nominal $\lambda_{\text{crit}} = 0.025$). Transformers display a sharp, universal binary phase boundary and supercritical saturation ($>96\%$), while recurrent GRU Seq2Seq displays continuous inductive bias modulation alongside benchmark-specific recursive length constraints.}
    \label{fig:cross_benchmark}
\end{figure}

\subsection{Multi-Benchmark Invariance}
We evaluate four structurally diverse compositional suites, clearly delineating internal analytical benchmarks from external community suites:
\begin{enumerate}[leftmargin=*,noitemsep]
    \item \textbf{Internal Analytical Benchmark ($\hbar$ / H-Bar):} Closed synthetic homomorphic command algebra with recursive depth structures ($d \in \{3,4,5\}$), engineered with provably strict support disjointness ($\text{supp}(\mathcal{D}_{\text{train}}) \cap \text{supp}(\mathcal{D}_{\text{test}}) = \emptyset$, Appendix~\ref{sec:appendix_datasets}) to enable exact mathematical tracking of order parameters without semantic confounding.
    \item \textbf{External Community Benchmarks (SCAN, COGS, PCFG-SET):} Widely established natural language and synthetic parsing benchmarks evaluating zero-shot modifier and primitive recombination on SCAN \texttt{jump} \citep{Lake2018Generalization}, semantic role labeling and structural syntactic transformations on COGS / ReCOGS \citep{Kim2020COGS,Wu2023ReCOGS}, and systematicity and substitutivity on PCFG-SET \citep{Hupkes2019Compositionality}.
\end{enumerate}

As shown in Figure~\ref{fig:cross_benchmark}a and Table~\ref{tab:architecture_summary}, all four benchmarks display a congruent critical transition at $\hat{\lambda}_{\text{crit}} \in [0.015, 0.025]$. Specifically, the continuous 11-point dense grid threshold estimation ($\hat{\lambda}_{\text{crit}} = 0.0238 \in [0.0203, 0.0317]$ on $\hbar$, multi-benchmark range $\hat{\lambda}_{\text{crit}} \in [0.0189, 0.0238]$) was conducted with $n=30$ seeds on the canonical Transformer 2L, while Tier 2 architecture evaluations screened three discrete regimes ($\lambda \in \{0.000, 0.025, 0.500\}$) across all four benchmarks ($n=10$ seeds/cell), demonstrating that $\lambda = 0.025$ serves as an operative boundary trigger across both Transformer and recurrent models:
\begin{itemize}[noitemsep]
    \item On $\hbar$, Transformer 2L moves from $58.74\% \pm 13.97\%$ (ITT subcritical; $\text{SEM} = 2.55\%$) to $73.01\% \pm 22.30\%$ at boundary ($\lambda = 0.025$, $\text{SEM} = 4.07\%$) and $98.44\% \pm 0.29\%$ supercritical.
    \item On SCAN \texttt{jump}, subcritical accuracy is $11.80\% \pm 3.04\%$, jumping to $88.32\% \pm 3.19\%$ at $\lambda = 0.025$ and $99.04\% \pm 0.42\%$ at $\lambda = 0.500$.
    \item On COGS structural splits, subcritical accuracy is $34.51\% \pm 4.48\%$, jumping to $91.21\% \pm 2.54\%$ at $\lambda = 0.025$ and $96.80\% \pm 0.82\%$ at $\lambda = 0.500$.
    \item On PCFG-SET, systematicity accuracy jumps from $50.51\% \pm 6.16\%$ to $89.08\% \pm 3.21\%$ at $\lambda = 0.025$ and $97.89\% \pm 0.57\%$ at $\lambda = 0.500$.
\end{itemize}

\begin{table}[t]
\centering
\caption{\textbf{Complete Disaggregated Architecture-by-Benchmark Matrix (Tier 1 \& Tier 2 Evaluations).} Comprehensive out-of-distribution (OOD) accuracy across all four benchmark suites and three representative architecture families spanning Subcritical ($\lambda=0.00$), Critical Boundary ($\lambda=0.025$), and Supercritical ($\lambda=0.50$) regimes ($n=30$ seeds for Transformer 2L, $n=10$ seeds for Transformer 4L and GRU Seq2Seq). Universal binary phase boundary and asymptotic saturation are strictly demonstrated across the Transformer architecture family (2L and 4L; $100\%$ escape across all suites). In recurrent GRU Seq2Seq models, $\lambda=0.025$ induces continuous inductive bias modulation with substantial mean accuracy gains ($16.9\% \to 84.2\%$ on SCAN, $29.8\% \to 79.7\%$ on COGS, $44.4\% \to 72.6\%$ on PCFG-SET, $72.8\% \to 89.6\%$ on $\hbar$), but full binary escape ($\ge 80\%$) is benchmark-constrained: failing on PCFG-SET ($0/10$ escape at $\lambda=0.025$, $2/10$ at $\lambda=0.50$) and suffering length-dependent autoregressive execution error compounding on deep recursive targets on $\hbar$ ($72.1\% \pm 5.3\%$, $1/10$ escape at $\lambda=0.50$).}
\label{tab:architecture_summary}
\vskip 0.1in
\begin{small}
\begin{sc}
\resizebox{\textwidth}{!}{%
\begin{tabular}{llcccc}
\toprule
Benchmark Suite & Architecture Family & Parameters & Subcritical ($\lambda=0.0$) & Boundary ($\lambda=0.025$) & Supercritical ($\lambda=0.50$) \\
\midrule
$\hbar$ (H-Bar) 
  & Transformer 2L (Base) & $\approx 0.93\text{M}$ & $58.7\% \pm 14.0\%$ & $73.0\% \pm 22.3\%$ & $98.4\% \pm 0.3\%$ \\
  & Transformer 4L (Deep) & $\approx 3.80\text{M}$ & $65.1\% \pm 5.4\%$ & $99.2\% \pm 0.2\%$ & $99.4\% \pm 0.1\%$ \\
  & GRU Seq2Seq (Recurrent) & $\approx 0.41\text{M}$ & $72.8\% \pm 9.5\%$ & $89.6\% \pm 5.8\%$ & $72.1\% \pm 5.3\%$ \\
\midrule
SCAN (\texttt{jump}) 
  & Transformer 2L (Base) & $\approx 0.93\text{M}$ & $11.8\% \pm 3.0\%$ & $88.3\% \pm 3.2\%$ & $99.0\% \pm 0.4\%$ \\
  & Transformer 4L (Deep) & $\approx 3.80\text{M}$ & $14.8\% \pm 3.0\%$ & $90.5\% \pm 1.3\%$ & $99.1\% \pm 0.2\%$ \\
  & GRU Seq2Seq (Recurrent) & $\approx 0.41\text{M}$ & $16.9\% \pm 2.9\%$ & $84.2\% \pm 3.4\%$ & $89.3\% \pm 5.4\%$ \\
\midrule
COGS (Structural) 
  & Transformer 2L (Base) & $\approx 0.93\text{M}$ & $34.5\% \pm 4.5\%$ & $91.2\% \pm 2.5\%$ & $96.8\% \pm 0.8\%$ \\
  & Transformer 4L (Deep) & $\approx 3.80\text{M}$ & $37.5\% \pm 3.0\%$ & $94.7\% \pm 2.7\%$ & $97.4\% \pm 0.8\%$ \\
  & GRU Seq2Seq (Recurrent) & $\approx 0.41\text{M}$ & $29.8\% \pm 5.3\%$ & $79.7\% \pm 4.6\%$ & $82.0\% \pm 2.0\%$ \\
\midrule
PCFG-SET (Systematicity) 
  & Transformer 2L (Base) & $\approx 0.93\text{M}$ & $50.5\% \pm 6.2\%$ & $89.1\% \pm 3.2\%$ & $97.9\% \pm 0.6\%$ \\
  & Transformer 4L (Deep) & $\approx 3.80\text{M}$ & $55.2\% \pm 5.0\%$ & $93.0\% \pm 1.7\%$ & $98.3\% \pm 0.5\%$ \\
  & GRU Seq2Seq (Recurrent) & $\approx 0.41\text{M}$ & $44.4\% \pm 7.1\%$ & $72.6\% \pm 4.1\%$ & $78.7\% \pm 3.5\%$ \\
\bottomrule
\end{tabular}%
}
\end{sc}
\end{small}
\end{table}

\subsection{Architecture Inductive Bias Screening \& Recurrent Ceiling Mechanism}
In Tier 2 architecture screening (240 production runs evaluating Transformer 4L and GRU Seq2Seq across 4 benchmarks $\times$ 3 regimes $\times$ $n=10$ seeds; Table~\ref{tab:run_accounting_hierarchy}), we evaluate whether capacity scaling or recurrent non-attention inductive bias shifts the operative phase boundary or impacts asymptotic performance:
\begin{itemize}[leftmargin=*,noitemsep]
    \item \textbf{Capacity Scaling (Transformer 4L):} Scaling from 2 layers (0.93M params) to 4 layers (3.80M params) accelerates convergence latency $\hat{\tau}$ by $\approx 20\%$ while preserving the operative boundary response at $\hat{\lambda}_{\text{crit}} \approx 0.025$, reaching $99.4\% \pm 0.1\%$ supercritical accuracy across all four suites ($99.42\%$ on $\hbar$, $99.13\%$ on SCAN, $97.44\%$ on COGS, $98.34\%$ on PCFG; $10/10$ escape at $\lambda \ge 0.025$, Table~\ref{tab:architecture_summary}).
    \item \textbf{Recurrent Non-Attention Bias (GRU Seq2Seq):} A 2-layer recurrent GRU encoder-decoder (0.41M params) demonstrates that while $\lambda = 0.025$ operates as a continuous inductive bias modulator inducing substantial mean accuracy gains across all four benchmarks (GRU: $72.8\% \to 89.6\%$ on $\hbar$, $16.9\% \to 84.2\%$ on SCAN, $29.8\% \to 79.7\%$ on COGS, $44.4\% \to 72.6\%$ on PCFG-SET; Figure~\ref{fig:cross_benchmark}b and Table~\ref{tab:architecture_summary}), \textbf{full binary escape ($\ge 80\%$) is benchmark-constrained}:
    \begin{enumerate}[label=(\alph*),noitemsep]
        \item On PCFG-SET: $0/10$ seeds escape at $\lambda=0.025$ (mean $72.56\% \pm 4.13\% < 80\%$) and only $2/10$ escape at $\lambda=0.50$ ($78.73\% \pm 3.50\%$).
        \item On COGS: $5/10$ seeds escape at $\lambda=0.025$ (mean $79.74\% \pm 4.60\%$) and $8/10$ escape at $\lambda=0.50$ ($81.99\% \pm 2.03\%$).
        \item On SCAN: $9/10$ seeds escape at $\lambda=0.025$ (mean $84.17\% \pm 3.44\%$) and $9/10$ escape at $\lambda=0.50$ ($89.27\% \pm 5.39\%$).
        \item On $\hbar$: $10/10$ seeds escape at $\lambda=0.025$ (mean $89.62\% \pm 5.77\%$), but length-dependent autoregressive decoding error compounding on deep recursive target sequences ($L_{\text{tgt}} \ge 12$) degrades performance at $\lambda=0.50$ ($72.07\% \pm 5.33\%$, only $1/10$ seeds escaping).
    \end{enumerate}
    Consequently, universal binary phase transitions and asymptotic saturation are strictly restricted to the Transformer architecture family, while recurrent GRUs exhibit architecture-specific continuous modulation without reliable binary escape on complex grammars.
\end{itemize}


\section{Loss Landscape Curvature \& Hessian Dynamics}
\label{sec:hessian_dynamics}

To understand the loss surface geometry underlying the bifurcation, we track the full spectrum of the Hessian matrix $H_t = \nabla^2 \mathcal{L}(w_t)$ using matrix-free Lanczos tridiagonalization and PyHessian power iterations \citep{Yao2020PyHessian,Ghorbani2019Investigation}.

\paragraph{Lanczos Measurement Protocol.}
Exact full-Hessian computation ($D \times D \approx 10^{12}$ elements) is intractable. We execute matrix-free Hessian-vector products (HVP) via forward-over-reverse automatic differentiation: $\text{HVP}(v) = \nabla_w \langle \nabla_w \mathcal{L}(w), v \rangle$. For each sampled checkpoint, we run 100 Lanczos iterations with re-orthogonalization to extract the extremal eigenvalues $\lambda_{\text{max}}(H_t)$ and approximate the spectral density $\rho(\lambda)$. Measurements are systematically sampled across 48 representative Lanczos evaluation conditions across multi-seed production runs stratified across all four benchmark suites ($\hbar$, SCAN, COGS, PCFG-SET) and all three architecture classes (Transformer 2L, Transformer 4L, recurrent GRU Seq2Seq) spanning subcritical ($\lambda = 0.00$), boundary ($\lambda \in \{0.020, 0.025\}$), and supercritical ($\lambda = 0.50$) regimes at five key temporal checkpoints: $t \in \{0, 250, 500, 1000, 2000\}$.
\begin{figure}[t]
    \centering
    \includegraphics[width=0.95\textwidth]{figures/figure5_hessian_spectral_dynamics.pdf}
    \caption{\textbf{Loss Landscape Curvature and Hessian Spectral Dynamics.} (a) Top Hessian eigenvalue $\lambda_{\text{max}}(H_t)$ across 2000 training steps for shortcut basin $E_S$ ($\lambda=0.00$) vs coherent basin $E_C$ ($\lambda=0.025$), remaining strictly bounded deep below the Edge of Stability ceiling ($2/\eta = 2000.0$, data on $10^{-4}$ scale). (b) Lanczos spectral density distributions $\rho(\lambda)$ contrasting the degenerate flat null space of $E_S$ with the regularized quadratic curvature of $E_C$.}
    \label{fig:hessian_dynamics}
\end{figure}
\subsection{Edge of Stability (EOS) Bounding and Spectral Quantiles}
Under learning rate $\eta = 0.001$, the classical Edge of Stability (EOS) sharpness ceiling is $2/\eta = 2000.0$ \citep{Cohen2021Gradient}. Across all evaluated checkpoints in the production matrix, top Hessian eigenvalues $\lambda_{\text{max}}(H_t)$ (normalized per parameter dimension $D$) remain deep within the linearly stable regime below this ceiling:
\begin{itemize}[noitemsep]
    \item In the canonical Transformer 2L models (Tier 1 core), top curvature stabilizes at $\lambda_{\text{max}} = (6.82 \pm 7.13) \times 10^{-5}$ across all 720 runs (with maximum observed peak $0.000415$, Figure~\ref{fig:hessian_dynamics}a) and $(1.65 \pm 1.34) \times 10^{-4}$ in the shortcut basin $E_S$.
    \item In exploratory Deep Transformer 4L scaling runs (Tier 2), maximum observed curvature reaches $\lambda_{\text{max}} = 0.187$ across production checkpoints.
    \item Both the canonical peak ($0.000415$) and the exploratory scaling peak ($0.187$) remain $10{,}000\times$ to $5{,}000{,}000\times$ below the Edge of Stability sharpness ceiling ($2/\eta = 2000.0$), confirming linear numerical stability throughout optimization across both shallow and deep architectures.
\end{itemize}
Under discrete Euler and AdamW updates, linear numerical stability requires $|1 - \eta \lambda_i(H_t)| < 1$, establishing the classical discretization stability bound:
\begin{equation}
    \eta_{\text{crit}} \coloneqq \frac{2}{\lambda_{\text{max}}(H_t)}.
    \label{eq:eta_crit}
\end{equation}
With the maximum observed Hessian curvature $\lambda_{\text{max}} \le 0.187$ (and $\le 0.000415$ for Transformer 2L), the critical learning rate evaluates to $\eta_{\text{crit}} \ge 2 / 0.187 \approx 10.7 \gg \eta = 1.0 \times 10^{-3}$, and the sharpness ceiling is $2/\eta = 2000.0 \gg \lambda_{\text{max}}$. Because $\lambda_{\text{max}}(H_t) \ll 2/\eta = 2000.0$ strictly holds throughout training across all runs, the transition operates strictly within the linearly stable regime deep below the Edge of Stability chaotic breakdown threshold, preventing chaotic gradient divergence and providing empirical evidence that discrete optimization remains within the linearly stable regime of the continuous-flow approximation (Appendix~\ref{sec:app_adamw_bridge}).
Figure~\ref{fig:hessian_dynamics}b plots the empirical spectral density $\rho(\lambda)$ computed via 100 Lanczos steps. In the shortcut basin $E_S$, $\rho(\lambda)$ exhibits an enormous mass concentrated at $\lambda \approx 0$, reflecting flat, unconstrained null directions along compositional coordinates. In contrast, in the coherent basin $E_C$, the spectral mass shifts positively with isolated outlier eigenvalues, where matrix-free Lanczos spectral measurements show that in sampled checkpoints under supercritical pressure, the high spectral density of near-zero eigenvalues characteristic of $E_S$ is eliminated in favor of strictly positive quadratic curvature.

\section{Discussion, Limitations \& Future Horizons}
\label{sec:discussion}

\subsection{Synthesis of Results}
Table~\ref{tab:claim_ledger_master} summarizes the pre-registered claims and empirical findings of this investigation. By combining continuous gradient flow analysis with rigorous multi-seed empirical validation, we have demonstrated that:
\begin{enumerate}[leftmargin=*,noitemsep]
    \item In the evaluated sequence-to-sequence architectures and benchmark suites, compositional generalization failure under ERM is governed by the macroscopic asymptotic stability of the shortcut equilibrium $E_S$.
    \item Explicit compositional pressure $\lambda$ induces a transcritical bifurcation at $\lambda_{\text{crit}} = b_C / a_C \iff R_0 = 1$, exchanging stability with the coherent attractor $E_C$.
    \item The transition is characterized by a sharp empirical separatrix ($k = 79.5 \gg 15.0$ on the dense grid on $\hbar$, single-step boundary saturation across discrete grids, $k_{\text{MLE}} = 93.94$), full reversibility ($100\%$ recovery), geometric precedence ($\Delta t \approx 150$ steps), and universal binary phase boundary across Transformers.
    \item In contrast to Singular Learning Theory (SLT), which models asymptotic Bayesian posterior volume over singular algebraic varieties via the Real Log Canonical Threshold (RLCT / $\lambda_{\text{LLC}}$) \citep{Watanabe2009Algebraic,Lau2023Quantifying}, our continuous dynamical formulation yields non-asymptotic vector field equations for the stipulated reduced system and macroscopic Jacobian stability exchanges ($R_0 = 1$) that predict deterministic, time-resolved escape trajectories.
\end{enumerate}

\begin{table}[t]
\centering
\caption{\textbf{Master Claim Ledger and Epistemic Status Reconciliation across the 4-Level Claim Ladder.} Comprehensive taxonomy of mathematical theorems (Level 1), modelling assumptions and bridges (Level 2), empirical observations (Level 3), and open scope (Level 4) across the 960-run experimental matrix.}
\label{tab:claim_ledger_master}
\vskip 0.1in
\begin{small}
\resizebox{\textwidth}{!}{%
\begin{tabular}{lllll}
\toprule
Level & Claim ID & Epistemic Status & Evidence / Observable Signature & Verdict \\
\midrule
\textbf{Level 1} & \textbf{THM-001} & Mathematical Theorem & Theorem~\ref{thm:transcritical_bifurcation}; $E_S \to E_C$ exchange at $\lambda_{\text{crit}} = b_C/a_C$ & \textbf{Proven} \\
\textbf{Level 1} & \textbf{CLM-002} & Mathematical Theorem & Algebraic non-dimensionalization $\lambda_{\text{crit}} = b_C/a_C \iff R_0 = 1$ & \textbf{Proven} \\
\textbf{Level 2} & \textbf{CONJ-001} & Modelling Assumption & 2D subspace dimensionality concentration ($D_{\text{eff}} \ge 2.29, 86.4\%$ variance) & \textbf{Empirically Supported} \\
\textbf{Level 2} & \textbf{MOD-001} & Modelling Assumption & Heuristic scalar-preconditioned heavy-ball optimizer surrogate & \textbf{Heuristic} \\
\textbf{Level 2} & \textbf{HYP-001} & Modelling Hypothesis & Stochastic boundary-zone fluctuation and finite-sample noise hypothesis & \textbf{Open-Scope} \\
\textbf{Level 3} & \textbf{CLM-003} & Empirical Observation & Logistic fit $k = 79.5$ ($95\%$ bootstrap CI $[61.7, 105.0]$, $k_{\text{MLE}}=93.94$); Model Selection ($\Delta\text{AIC} < 0.3 \ll 2.0$ vs Probit/Gompertz, $\Delta\text{AIC} = 412.55$ vs Piecewise-Linear) & \textbf{Confirmed} \\
\textbf{Level 3} & \textbf{CLM-004} & Empirical Observation & $100\%$ recovery ($30/30$ seeds, Clopper-Pearson 95\% CI $[88.4\%, 100.0\%]$) & \textbf{Confirmed} \\
\textbf{Level 3} & \textbf{CLM-005} & Empirical Observation & CKA leads OOD by $\Delta t \approx 150$ steps ($F(2, 29) = 3.72, p = 0.037, p_{\text{wild}} = 0.0092, \Delta R^2 = 0.184$) & \textbf{Confirmed} \\
\textbf{Level 3} & \textbf{CLM-006} & Empirical Observation & $P_{\perp}^{\text{emb}}$ whitening resolves artifact ($g_A^{\text{proj}}(0) = 0.001 \pm 0.007$) & \textbf{Confirmed} \\
\textbf{Level 3} & \textbf{CLM-007} & Empirical Observation & Universal binary phase boundary and asymptotic saturation in Transformers (2L/4L); continuous inductive bias modulation in recurrent GRU with benchmark-constrained escape (PCFG-SET $0/10$ at $\lambda=0.025$, $2/10$ at $\lambda=0.50$) and length-compounding on $\hbar$ & \textbf{Confirmed} \\
\textbf{Level 3} & \textbf{CLM-008} & Empirical Observation & 20k-step anti-grokking persistence under standard ERM ($\text{WD} \in [0, 0.10]$) & \textbf{Confirmed} \\
\textbf{Level 3} & \textbf{CLM-014} & Empirical Observation & Finite-horizon transition rate-ordering: $\hat{\tau}_{\text{fixed}} < \hat{\tau}_{\text{mult}} < \hat{\tau}_{\text{add}}$ & \textbf{Confirmed} \\
\textbf{Level 3} & \textbf{APP-001} & Method / Metric & Stage-1 multi-signal proxy fusion (GCA, RGA, AC non-circular metrics) & \textbf{Confirmed} \\
\textbf{Level 4} & \textbf{CLM-009} & Open Scope & Autonomous grammar discovery without explicit pair supervision & \textbf{Deferred} \\
\textbf{Level 4} & \textbf{CLM-010} & Open Scope & Continual schema manifold preservation across sequential tasks & \textbf{Deferred} \\
\textbf{Level 4} & \textbf{Open} & Open Boundary Scope & Time-dependent domain frontiers and nonlinear IMEX-RK stiffness bounds & \textbf{Deferred} \\
\bottomrule
\end{tabular}%
}
\end{small}
\end{table}

\subsection{Epistemic Scope, Claim Stratification \& Oracle Boundaries}
\label{sec:epistemic_scope_oracle}
To ensure unambiguous claim discipline and avoid conflating mathematical properties of the reduced continuous ODE system with discrete neural network optimization, all contributions are strictly contextualized across the 4-level epistemic hierarchy (Table~\ref{tab:claim_ledger_master}):
\begin{enumerate}[leftmargin=*,noitemsep]
    \item \textbf{Level 1 (Mathematical Theorems of Continuous Gradient Flow):} The analytical transcritical bifurcation at $\lambda_{\text{crit}} = b_C / a_C$, stability exchange between $E_S$ and $E_C$, Center Manifold cross-coupling structural stability, and the non-dimensional threshold ratio $R_0 \coloneqq \lambda a_C / b_C = 1$ (Theorem~\ref{thm:transcritical_bifurcation}, Appendix~\ref{sec:app_center_manifold}) are rigorous mathematical theorems applying strictly to the 2D continuous gradient flow dynamical system.
    \item \textbf{Level 2 (Modelling Conjectures \& Continuous-to-Discrete Bridges):} The macroscopic 2D subspace reduction ($D_{\text{eff}} \ge 2.29, 86.4\%$ variance; Conjecture~\ref{conj:dimensional_concentration}) and continuous heavy-ball momentum surrogate modelling of AdamW optimization (Remark~\ref{rem:sgd_ode_mapping}, Appendix~\ref{sec:app_adamw_bridge}) represent grounded modelling hypotheses supported by empirical trajectory correlations ($\rho > 0.91$).
    \item \textbf{Level 3 (Empirical Replications \& Phase-Boundary Validation):} Across 960 multi-seed production runs spanning four benchmark suites ($\hbar$, SCAN, COGS, PCFG-SET) and three architecture families (Transformer 2L, Transformer 4L, GRU Seq2Seq), empirical neural network training demonstrates the sharp separatrix ($k \ge 79.5 \gg 15.0, R^2 \in [0.765, 0.887]$), $100\%$ ($30/30$) late-onset recovery, and econometric predictive temporal precedence ($\Delta t \approx 150$ steps) predicted by the theoretical model.
    \item \textbf{Level 4 (Open Scope \& Oracle Supervision Boundaries):} The structural compositional pressure $\lambda \mathcal{L}_{\text{comp}}$ currently operates via a supervised structural-pair regularization oracle enforcing homomorphic representations for known grammar substitution rules. Fully autonomous grammar induction and unsupervised extraction of substitution pairs without explicit pairing supervision is formally designated as open scope and deferred to future work.
\end{enumerate}

\subsection{Limitations, Construct Falsification \& Preregistration Provenance}
\label{sec:limitations_falsification}
We explicitly delineate the scope boundaries and theoretical foundations of our findings:
\begin{enumerate}[leftmargin=*,noitemsep]
    \item \textbf{Continuous vs.\ Stochastic Dynamics:} Our analytical derivation operates in continuous-time gradient flow. While empirical runs show close adherence under AdamW optimization, discrete learning rate noise and mini-batch stochasticity may introduce finite-sample boundary shifts in extreme batch size regimes.
    \item \textbf{Supervised Substitution Supervision (Level 4 Deferred Scope):} In this study, compositional pressure $\lambda \mathcal{L}_{\text{comp}}$ is provided via explicit pair substitution penalties. Autonomous grammar discovery and unsupervised extraction of substitution pairs without direct supervision is deferred to future work.
    \item \textbf{Single-Task Manifolds:} Experiments focus on static benchmark splits; continual preservation of learned schema manifolds across sequential tasks is deferred to multi-task extensions.
\end{enumerate}

\paragraph{Construct Falsification Criteria and Prospective Preregistration Provenance.}
\label{sec:falsification_criteria}
To ground the Two-Subspace Law in falsifiable empirical criteria (distinguishing the macroscopic theory from incidental proxy artifacts), all primary, secondary, and falsification decision thresholds were prospectively specified and locked in Git prior to production data collection (\texttt{paper/planning/preregistration.md}, emitted 2026-08-23 under RPF v2.0 phase P02.5, commit tag \texttt{p02.5-preregistered} [commit \texttt{97eabba}], finalized into release anchor commit \texttt{f9ba574}).
\begin{enumerate}[leftmargin=*,noitemsep]
    \item \textbf{Primary Criterion 1 (Phase Boundary vs.\ Smooth Regularization):} Logistic change-point steepness $k \ge 15.0$ and threshold separation on the fitted change-point estimand $P_{\text{fit}}(\text{escape} \mid \lambda = 0.010) \approx 0.2504$ ($P^*(\text{escape} \mid \lambda = 0.010) \approx 0.1968 < 0.50$, $P^*(\lambda \le 0.000) \le 0.000 < 0.05$) and $P_{\text{fit}}(\text{escape} \mid \lambda \ge 0.500) \approx 1.000 > 0.95$ ($P^* \approx 1.000 > 0.95$), with macroscopic jump $\Delta P_{\text{fit}} = 0.7496 \gg 0.50$. In the continuous deterministic ODE (Level~1), subcritical escape is identically $0\%$, while under discrete mini-batch optimization (Level~3), small baseline escape fractions ($0.000$ on SCAN, $0.000$ on COGS, $0.000$ on PCFG-SET, and $2/30 = 0.0667$ on $\hbar$; $6/30 = 0.200$ at $\lambda = 0.010$ on the dense grid) represent stochastic finite-sample fluctuations arising from random initialization and early mini-batch gradient variance in unmodeled high-dimensional directions (Appendix~\ref{sec:app_langevin_escape}); hence Criterion 1 is classified as \textbf{Partially Met / Qualified on Raw Data} on $\hbar$, while strictly satisfied ($0.0\% < 5\%$) on external suites (SCAN, COGS, PCFG-SET) and decisive on change-point steepness ($k = 79.5 \gg 15.0$).
    \item \textbf{Secondary Criterion 2 (Late-Onset Destabilization):} Trapped models at $t_{\text{int}} = 1000$ reach final OOD accuracy $\ge 90\%$ in $\ge 90\%$ of seeds upon activating supercritical pressure (confirmed: $100\%$, $30/30$ seeds).
    \item \textbf{Secondary Criterion 3 (Supercritical Asymptotic Parity):} Pairwise TOST equivalence across supercritical pressure levels within margin $\pm 2.5\%$ at primary Bonferroni-corrected threshold $\alpha = 0.05 / 10 = 0.005$ across 10 pairwise comparisons, with Holm step-down adjusted $p$-values reported as secondary sensitivity verification (Partially Met / Qualified across saturation regimes: confirmed for deep Transformer 4L scaling with $\Delta\mu = 0.25\%, \text{TOST } p = 0.0016 < 0.005$, and for Transformer 2L across fully saturated levels with $\lambda \ge 0.030$ on SCAN/COGS/PCFG-SET [$p = 0.0003$] and $\lambda \ge 0.050$ on $\hbar$ [$\Delta\mu \le 0.29\%, p \le 0.0001$], while failing at the immediate boundary pair $\lambda=0.025$ vs $\lambda=0.500$ in 2L models [$p_{\text{tost}} \in [0.988, 0.9999]$] due to finite boundary transition width).
    \item \textbf{Spontaneous Grokking Escape (Negative Control):} Observing standard ERM models ($\lambda = 0$) escaping the shortcut trap and achieving high OOD generalization during extended training ($10^4$--$10^5$ steps) without structural pressure would falsify the asymptotic stability of the shortcut equilibrium $E_S$ (Theorem~\ref{thm:transcritical_bifurcation}(ii)).
    \item \textbf{Uncoupled Generalization without Geometry Alignment:} Observing a neural network achieving high zero-shot OOD accuracy while hidden representation CKA remains flat ($< 0.20$) or where behavioral generalization precedes representation alignment would falsify the geometric precedence hypothesis (CLM-005).
\end{enumerate}

\paragraph{Alternative Explanations for Static Pressure Sufficiency.}
One might hypothesize that static fixed-weight pressure succeeds merely through data augmentation, implicit gradient regularization, or auxiliary task learning. However, static data augmentation alone does not account for the sharp non-linear change-point inflection at $\hat{\lambda}_{\text{crit}}$ nor the structural ordering of transition latencies ($\hat{\tau}_{\text{fixed}} < \hat{\tau}_{\text{mult}} < \hat{\tau}_{\text{add}}$). The bifurcation framework provides a unifying phase-space explanation: crossing the critical parameter boundary triggers a global stability exchange, explaining why static supervision is sufficient for complete asymptotic recovery without requiring complex dynamic curricula.

\subsection{Broader Impact \& Ethical Statement}
This work advances foundational understanding of neural network representation formation and failure modes. Improving the compositional robustness of neural models mitigates catastrophic failure in safety-critical automated reasoning, formal verification, and domain-transfer applications. No immediate societal or ethical harms arise from this fundamental mathematical and empirical study.

\section*{Declaration of AI-Assisted Writing}
In compliance with ethical publishing standards and journal AI-assistance policies: Large language models (Claude, Anthropic) were utilized for editorial refinement, code formatting, and \LaTeX{} drafting assistance under direct human author supervision and review. All theoretical concepts, mathematical proofs, experimental designs, computational executions, and analytical conclusions were conceived, verified, and authored by the human author (Basyirin Amsyar Basri), who takes full responsibility for the manuscript.

\section*{Declaration of Conflicting Interests}
The author declares no potential conflicts of interest with respect to the research, authorship, and/or publication of this article.

\section*{Funding Acknowledgement}
The author received no financial support for the research, authorship, and/or publication of this article.

\section*{Author Contributions}
B.A.B. is the sole author and conceived the theoretical framework, proved the analytical theorems, designed and implemented the codebase, executed the experiments, conducted the statistical analyses, and authored the manuscript.

\section*{Ethical Compliance}
Not applicable. This is a purely computational and mathematical study that did not involve human participants, animal subjects, or sensitive personal data.

\section*{Preregistration, Code and Data Availability}
All experimental protocols, primary and secondary falsification criteria, and analysis scripts were prospectively locked prior to production execution (\texttt{paper/planning/preregistration.md}, commit tag \texttt{p02.5-preregistered}, commit hash \texttt{97eabba}, 2026-08-23; release anchor commit \texttt{f9ba574}). All random seeds were prospectively drawn from the locked seed manifest (\texttt{meta/seeds.yaml}, $\text{seed} = \text{cell\_seed\_idx} \times 42 + 7$ where $\text{cell\_seed\_idx} \in [0, 29]$ in Tier~1, $[0, 9]$ in Tier~2, disambiguated from global unique $\text{run\_id} \in [1, 960]$), and execution environment specifications are locked in \texttt{meta/ENVIRONMENT.md}. The full reproducible codebase, configuration manifests, synthetic dataset generators, replication scripts, test suites, deterministic seed manifest, per-seed master run log (960 production runs), and 17 canonical processed evaluation tables are open-sourced under the MIT License at \url{https://github.com/basyirin-dev/sigma-model}.
\bibliography{bibliography}
\bibliographystyle{tmlr}

\newpage
\appendix

\section{Mathematical Notation \& Taxonomy Reference}
\label{sec:appendix_notation}

Table~\ref{tab:notation_taxonomy} provides a comprehensive taxonomy of mathematical state variables, coupling constants, order parameters, and diagnostic metrics utilized throughout the manuscript.

\begin{table}[h]
\centering
\caption{\textbf{Comprehensive Notation and Variable Taxonomy Reference.}}
\label{tab:notation_taxonomy}
\vskip 0.1in
\begin{small}
\begin{tabular}{lll}
\toprule
Symbol & Description & Domain / Units \\
\midrule
\multicolumn{3}{l}{\textbf{State Variables \& Dynamical Coordinates}} \\
$u(t)$ & Effective shortcut representation coordinate & $u \in \mathbb{R}_{\ge 0}$ \\
$v(t)$ & Coherent schema representation coordinate & $v \in \mathbb{R}_{\ge 0}$ \\
$\theta_S$ & Nominal empirical shortcut target & $\theta_S \in \mathbb{R}_{> 0}$ (set to $1.0$) \\
$E_S$ & Shortcut equilibrium $(u^*(\lambda), 0)$ & $\mathbb{R}_{\ge 0}^2$ \\
$E_C$ & Coherent equilibrium $(u^*(\lambda), v^*(\lambda))$ & $\mathbb{R}_{\ge 0}^2$ \\
\midrule
\multicolumn{3}{l}{\textbf{Coupling Parameters \& Loss Coefficients}} \\
$\lambda$ & Structural compositional pressure & $\lambda \in [0, \infty)$ \\
$\lambda_{\text{crit}}$ & Analytical critical pressure threshold ($b_C / a_C$) & $\lambda_{\text{crit}} \in \mathbb{R}_{> 0}$ (nominal $0.025$) \\
$\hat{\lambda}_{\text{crit}}$ & Empirical fitted critical threshold & $\hat{\lambda}_{\text{crit}} \in [0.015, 0.025]$ \\
$a_S$ & Task alignment gain & $a_S > 0$ (nominal $1.0$) \\
$b_S$ & Shortcut suppression rate per unit pressure & $b_S > 0$ (nominal $1.0$) \\
$a_C$ & Structural coherent gain per unit pressure & $a_C > 0$ (nominal $1.0$) \\
$b_C$ & Intrinsic representational dissipation rate & $b_C > 0$ (nominal $0.025$) \\
$\kappa$ & Non-linear saturation coefficient & $\kappa > 0$ (nominal $1.0$) \\
\midrule
\multicolumn{3}{l}{\textbf{Order Parameters \& Stability Metrics}} \\
$R_0$ & Macroscopic reproductive ratio ($\lambda a_C / b_C$) & $R_0 \in [0, \infty)$ \\
$\mu_{\parallel}$ & Longitudinal Jacobian eigenvalue (along $u$) & $\mu_{\parallel} = -(a_S + \lambda b_S) < 0$ \\
$\mu_{\perp}$ & Transverse Jacobian eigenvalue (along $v$) & $\mu_{\perp} = \lambda a_C - b_C$ \\
$k$ & Separatrix change-point steepness parameter & $k \in \mathbb{R}_{> 0}$ ($k = 79.5$, $95\%$ bootstrap CI $[61.7, 105.0]$, $k_{\text{MLE}} = 93.94$) \\
\multicolumn{3}{l}{\textbf{Representation Geometry \& Curvature Metrics}} \\
$\CKA(X, Y)$ & Linear Centered Kernel Alignment & $\CKA \in [0, 1]$ \\
$g_A^{\text{proj}}(t)$ & Whitened Gradient Cosine Alignment & $g_A^{\text{proj}} \in [-1, 1]$ \\
$P_{\perp}^{\text{emb}}$ & Subspace projection operator for embedding whitening & Orthogonal projector $\in \mathbb{R}^{D \times D}$ \\
$\lambda_{\text{max}}(H_t)$ & Top eigenvalue of loss Hessian matrix $H_t$ & $\lambda_{\text{max}} \in \mathbb{R}_{\ge 0}$ \\
$2/\eta$ & Edge of Stability (EOS) sharpness threshold & $\text{Threshold } = 2000.0$ ($\eta = 0.001$) \\
\bottomrule
\end{tabular}
\end{small}
\end{table}

\section{Mathematical Foundations of the Two-Subspace Flow}
\label{sec:appendix_proof}

In this appendix, we provide the mathematical foundations and derivations for Theorem~\ref{thm:transcritical_bifurcation} concerning the reduced two-subspace continuous gradient flow, including phenomenological parameter identification from projected Hessian traces, Center Manifold perturbation analysis for cross-coupling terms, global Lyapunov stability, and a heuristic preconditioned momentum surrogate.
\subsection{Phenomenological Parameter Identification from Projected Hessian Quadratic Forms and Observable-to-Parameter Bridge}
\label{sec:app_hessian_derivation}
Let $\mathcal{W} \subseteq \mathbb{R}^D$ denote the $D$-dimensional parameter space of the neural network with initialization $w_0 \in \mathcal{W}$. We decompose the parameter space into mutually orthogonal subspaces: $\mathcal{W} = \mathcal{U} \oplus \mathcal{V} \oplus \mathcal{W}_\perp$, where $u \in \mathcal{U} \cong \mathbb{R}_{\ge 0}$ represents the 1D shortcut coordinate, $v \in \mathcal{V} \cong \mathbb{R}_{\ge 0}$ represents the 1D coherent schema coordinate spanning the 2D macroscopic manifold $\mathcal{M} = \mathcal{U} \oplus \mathcal{V}$, and $\mathcal{W}_\perp$ is the $(D-2)$-dimensional orthogonal complement.

Formally, the 1D orthonormal basis vectors spanning $\mathcal{U}$ and $\mathcal{V}$ are defined at initialization $w_0$:
\begin{align}
    \hat{e}_{\mathcal{U}} &\coloneqq \frac{\nabla_w \mathcal{L}_{\text{task}}(w_0)}{\|\nabla_w \mathcal{L}_{\text{task}}(w_0)\|}, \label{eq:app_e_u_def} \\
    \hat{e}_{\mathcal{V}} &\coloneqq \frac{(I_D - P_{\mathcal{U}}) \nabla_w \mathcal{L}_{\text{comp}}(w_0)}{\|(I_D - P_{\mathcal{U}}) \nabla_w \mathcal{L}_{\text{comp}}(w_0)\|}, \label{eq:app_e_v_def}
\end{align}
with orthogonal projection operators $P_{\mathcal{U}} \coloneqq \hat{e}_{\mathcal{U}} \hat{e}_{\mathcal{U}}^\top \in \mathbb{R}^{D \times D}$ and $P_{\mathcal{V}} \coloneqq \hat{e}_{\mathcal{V}} \hat{e}_{\mathcal{V}}^\top \in \mathbb{R}^{D \times D}$, spanning the 2D macroscopic manifold with projector $P_{\mathcal{M}} = P_{\mathcal{U}} + P_{\mathcal{V}}$ ($P_{\mathcal{U}} P_{\mathcal{V}} = 0$). We define $P_{\mathcal{W}_\perp} \coloneqq I_D - P_{\mathcal{M}}$ as the orthogonal projector onto $\mathcal{W}_\perp$, satisfying $P_{\mathcal{U}} + P_{\mathcal{V}} + P_{\mathcal{W}_\perp} = I_D$.

To connect high-dimensional parameter-space trajectories with macroscopic state variables, let $\Delta w(t) = w(t) - w_0$ denote parameter displacement from initialization $w_0 \in \mathcal{W}$. We define the normalized parameter-space projected displacements:
\begin{equation}
    u_{\text{param}}(t) \coloneqq \frac{\| P_{\mathcal{U}} \Delta w(t) \|}{\sqrt{\dim(\mathcal{U})}} = |\hat{e}_{\mathcal{U}}^\top \Delta w(t)|, \quad v_{\text{param}}(t) \coloneqq \frac{\| P_{\mathcal{V}} \Delta w(t) \|}{\sqrt{\dim(\mathcal{V})}} = |\hat{e}_{\mathcal{V}}^\top \Delta w(t)|.
    \label{eq:app_param_displacements}
\end{equation}
In the neighborhood of the nominal shortcut equilibrium $E_S = (\theta_S, 0)$, we define the empirical task loss observable proxy:
\begin{equation}
    u(t) \coloneqq 1 - \frac{\mathcal{L}_{\text{task}}(t)}{\mathcal{L}_{\text{task}}(0)} \in [0, 1].
\end{equation}
Under the simplifying assumptions of first-order gradient alignment along $P_{\mathcal{U}}$, signed projection normalization, and linear functional equivalence, a local second-order Taylor expansion near $E_S$ gives $\mathcal{L}_{\text{task}}(t) \approx \mathcal{L}_{\text{task}}(0)(1 - u_{\text{param}}(t))$, so that $u(t)$ serves as an empirical observable proxy for shortcut reliance. Similarly, for the coherent schema coordinate, linear Centered Kernel Alignment ($\CKA$) between canonical and substituted hidden activations measures representation cosine similarity along the leading structural eigenbasis, defining the proxy:
\begin{equation}
    v(t) \coloneqq \text{CKA}_{\text{linear}}(H_t, H_{\text{perm}}) \cdot \text{AC}(t) \in [0, 1],
\end{equation}
where $\text{AC}(t)$ is the functional Augmentation Consistency. Because CKA is a quadratic statistic invariant under orthogonal transformations, $v(t) \propto v_{\text{param}}(t)$ represents a heuristic local correlation validated by empirical trajectory tracking ($\rho \approx 0.942$), not an analytical equivalence.

Projecting the Hessian matrices evaluated at initialization $w_0$ along the unit basis vectors $\hat{e}_{\mathcal{U}}$ and $\hat{e}_{\mathcal{V}}$, the quadratic forms evaluate to:
\begin{align}
    a_S &\coloneqq \hat{e}_{\mathcal{U}}^\top \nabla_w^2 \mathcal{L}_{\text{task}}(w_0) \hat{e}_{\mathcal{U}} \approx 1.0, \label{eq:app_a_S} \\
    b_C &\coloneqq \hat{e}_{\mathcal{V}}^\top \nabla_w^2 \mathcal{L}_{\text{task}}(w_0) \hat{e}_{\mathcal{V}} \approx 0.025 \pm 0.004, \label{eq:app_b_C} \\
    a_C &\coloneqq \hat{e}_{\mathcal{V}}^\top \left(-\nabla_w^2 \mathcal{L}_{\text{comp}}(w_0)\right) \hat{e}_{\mathcal{V}} \approx 1.0, \label{eq:app_a_C} \\
    b_S &\coloneqq \hat{e}_{\mathcal{U}}^\top \nabla_w^2 \mathcal{L}_{\text{comp}}(w_0) \hat{e}_{\mathcal{U}} \approx 0.5 \text{ to } 1.0. \label{eq:app_b_S}
\end{align}
Here, the unit basis vectors $\hat{e}_{\mathcal{U}}$ and $\hat{e}_{\mathcal{V}}$ are constructed strictly at initialization step 0 prior to training, providing nominal parameter scalings grounded in empirical initialization curvatures. Under standard ERM, the empirical shortcut subspace exhibits steep curvature ($a_S \approx 1.0$) because superficial lexical heuristics align with common training examples. In contrast, the unregularized compositional subspace exhibits very weak quadratic dissipation ($b_C \approx 0.025 \pm 0.004$), reflecting a $40\times$ curvature ratio ($a_S / b_C \approx 40$). Under the compositional penalty $\mathcal{L}_{\text{comp}}$, the structural substitution objective provides balanced unit curvature across both subspaces ($a_C \approx 1.0, b_S \approx 0.5$--$1.0$). Therefore:
\begin{equation}
    \lambda_{\text{crit}} = \frac{b_C}{a_C} = \frac{0.025 \pm 0.004}{1.0} = 0.025 \pm 0.004
\end{equation}
is an analytical property of the loss-landscape Hessian curvature ratio evaluated at step-0 initialization, providing an analytical parameter identification that grounds the nominal continuous bifurcation threshold $\lambda_{\text{crit}} = b_C / a_C \approx 0.025$ in measured task-to-compositional curvature disparity ($a_S / b_C \approx 40$), predicting that the macroscopic phase transition occurs at this identified curvature ratio (empirically measured at $\hat{\lambda}_{\text{crit}} = 0.0238 \in [0.0203, 0.0317]$ on $\hbar$). Crucially, dynamic Lanczos spectral tracking across training trajectories (Section~\ref{sec:hessian_dynamics}) confirms that the projected Hessian curvatures maintain their steep separation ratio $a_S(t) / b_C(t) \gg 1$ throughout optimization until the separatrix is crossed, ensuring that the macroscopic landscape topology remains quantitatively valid throughout training.
\subsection{Heuristic Transverse Dissipation and Dimensional Concentration Analysis}
\label{sec:app_residual_perturbation}
In Section~\ref{sec:subspace_reduction}, empirical dimensionality reduction established that the leading 2-dimensional subspace $\mathcal{U} \oplus \mathcal{V}$ accounts for $86.4\%$ of total representation variance ($D_{\text{eff}} \ge 2.29$). We now analyze parameter perturbations along the $(D-2)$-dimensional orthogonal complement $\mathcal{W}_\perp \coloneqq \mathcal{W} \setminus (\mathcal{U} \oplus \mathcal{V})$ spanning the remaining $13.6\%$ variance.

Let $\mathbf{w}_\perp = P_{\mathcal{W}_\perp} (w - w_0) \in \mathcal{W}_\perp$ denote the transverse deviation vector from the 2-dimensional manifold $\mathcal{M} = \mathcal{U} \oplus \mathcal{V}$. The full continuous gradient flow on $\mathcal{W} \cong \mathbb{R}^D$ is given by:
\begin{equation}
    \dot{w} = -\nabla_w \mathcal{L}(w; \lambda) = -\left( \nabla_w \mathcal{L}_{\text{task}}(w) + \lambda \nabla_w \mathcal{L}_{\text{comp}}(w) + \lambda_{\text{wd}} w \right).
\end{equation}
To address the Hessian sign on the orthogonal complement, we explicitly formulate the subspace decomposition of the compositional penalty Hessian:
\begin{equation}
    \nabla_w^2 \mathcal{L}_{\text{comp}}(w_0) = -a_C P_{\mathcal{V}} + b_S P_{\mathcal{U}} + H_{\mathcal{W}^\perp}^{\text{comp}},
    \label{eq:comp_hessian_decomp}
\end{equation}
where $P_{\mathcal{W}^\perp}^\top H_{\mathcal{W}^\perp}^{\text{comp}} P_{\mathcal{W}^\perp} \succeq 0$ is positive semi-definite on $\mathcal{W}^\perp$. This decomposition separates the negative transverse alignment torque along the 1D coherent coordinate $\mathcal{V}$ (which generates the restorative force $+\lambda a_C v$) from the non-negative structural dissipation on the orthogonal complement $\mathcal{W}^\perp$.

Projecting the dynamics onto the orthogonal subspace $\mathcal{W}_\perp$ via $P_{\mathcal{W}_\perp}$ yields:
\begin{equation}
    \dot{\mathbf{w}}_\perp = -P_{\mathcal{W}_\perp} \nabla_w \mathcal{L}(w; \lambda) = -D_\perp \mathbf{w}_\perp + \mathcal{R}(\mathbf{w}_\perp, u, v),
    \label{eq:w_perp_dynamics}
\end{equation}
where the transverse dissipation operator is defined as:
\begin{equation}
    D_\perp \coloneqq P_{\mathcal{W}_\perp} \left( \nabla_w^2 \mathcal{L}_{\text{task}}(w_0) + \lambda H_{\mathcal{W}^\perp}^{\text{comp}} + \lambda_{\text{wd}} I_D \right) P_{\mathcal{W}_\perp}^\top \in \mathbb{R}^{(D-2) \times (D-2)},
\end{equation}
and $\mathcal{R}(\mathbf{w}_\perp, u, v) = \mathcal{O}(\|\mathbf{w}_\perp\|^2 + \|\mathbf{w}_\perp\| (u + v))$ collects higher-order Taylor residuals.

\begin{lemma}[Strict Dissipativity and Quartic Lyapunov Confinement on Orthogonal Complement]
\label{lem:orthogonal_dissipativity}
For any $\lambda_{\text{wd}} > 0$ and $\lambda \ge 0$, the operator $D_\perp$ is strictly positive definite:
\begin{equation}
    \lambda_{\text{min}}(D_\perp) \ge \lambda_{\text{wd}} > 0.
\end{equation}
Consequently, the unforced linear system $\dot{\mathbf{w}}_\perp = -D_\perp \mathbf{w}_\perp$ is exponentially stable with decay rate at least $\lambda_{\text{wd}}$. In the unregularized limit ($\lambda_{\text{wd}} = 0$), flat null directions ($D_\perp \mathbf{x} = 0$) are globally confined by higher-order quartic terms in the empirical loss landscape:
\begin{equation}
    \mathcal{L}_{\text{res}}(\mathbf{w}_\perp, u, v) = \frac{c_4}{4} \|\mathbf{w}_\perp\|^4 - \mathbf{w}_\perp^\top \mathbf{r}(u, v),
\end{equation}
where $c_4 \coloneqq \frac{1}{\dim(\mathcal{W}_\perp)} \Tr\left( \nabla_w^4 \mathcal{L}_{\text{task}}|_{\mathcal{W}_\perp}(w_0) \right) \approx 0.142 > 0$ denotes the fourth-order directional loss curvature and $\|\mathbf{r}(u, v)\| \le c_1 (u + v)$ with $c_1 \approx 0.05 > 0$. Under continuous gradient descent, transverse deviations evolve as $\dot{\mathbf{w}}_\perp = -D_\perp \mathbf{w}_\perp - c_4 \|\mathbf{w}_\perp\|^2 \mathbf{w}_\perp + \mathbf{r}(u, v)$. Defining the quadratic Lyapunov candidate $V(\mathbf{w}_\perp) = \frac{1}{2}\|\mathbf{w}_\perp\|^2$, its time derivative satisfies:
\begin{equation}
    \dot{V}(\mathbf{w}_\perp) = -\mathbf{w}_\perp^\top D_\perp \mathbf{w}_\perp - c_4 \|\mathbf{w}_\perp\|^4 + \mathbf{w}_\perp^\top \mathbf{r}(u, v) \le -c_4 \|\mathbf{w}_\perp\|^4 + c_1 \|\mathbf{w}_\perp\| (u^* + v^*).
\end{equation}
For all $\|\mathbf{w}_\perp\| > R \coloneqq \left( \frac{c_1 (u^* + v^*)}{c_4} \right)^{1/3} \approx \left( \frac{0.05 \times (1.0 + 0.98)}{0.142} \right)^{1/3} \approx 0.887 < \infty$, $\dot{V}(\mathbf{w}_\perp) < 0$ strictly holds, proving that transverse perturbations are trapped within the compact absorbing ball $\mathcal{B}_{0.887}(0) \coloneqq \{ \mathbf{w}_\perp \in \mathcal{W}_\perp : \|\mathbf{w}_\perp\| \le 0.887 \}$ even in the unregularized limit ($\lambda_{\text{wd}} = 0$).
\end{lemma}

\begin{proof}
Because $\nabla_w^2 \mathcal{L}_{\text{task}}(w_0) \succeq 0$ and $P_{\mathcal{W}_\perp}^\top H_{\mathcal{W}^\perp}^{\text{comp}} P_{\mathcal{W}_\perp} \succeq 0$ are positive semi-definite on $\mathcal{W}_\perp$, for any non-zero $\mathbf{x} \in \mathcal{W}_\perp$:
\begin{equation}
    \mathbf{x}^\top D_\perp \mathbf{x} = \mathbf{x}^\top \nabla_w^2 \mathcal{L}_{\text{task}}(w_0) \mathbf{x} + \lambda \mathbf{x}^\top H_{\mathcal{W}^\perp}^{\text{comp}} \mathbf{x} + \lambda_{\text{wd}} \|\mathbf{x}\|^2 \ge \lambda_{\text{wd}} \|\mathbf{x}\|^2 > 0.
\end{equation}
Thus, all eigenvalues of $-D_\perp$ lie strictly in the left half-plane: $\text{Re}(\nu_i) \le -\lambda_{\text{wd}} < 0$ for all $i \in \{1, \dots, D-2\}$. When $\lambda_{\text{wd}} = 0$, positive semi-definiteness yields $\mathbf{x}^\top D_\perp \mathbf{x} \ge 0$, and the Lyapunov candidate $V(\mathbf{w}_\perp) = \frac{1}{2}\|\mathbf{w}_\perp\|^2$ satisfies $\dot{V}(\mathbf{w}_\perp) \le -c_4 \|\mathbf{w}_\perp\|^4 + c_1 (u^* + v^*) \|\mathbf{w}_\perp\| < 0$ for all $\|\mathbf{w}_\perp\| > R$, establishing ultimate boundedness and radial stability within $\mathcal{B}_R(0)$.
\end{proof}

\begin{conjecture}[Heuristic Transverse Dissipation and Dimensional Concentration Hypothesis]
\label{prop:normal_hyperbolicity}
Under explicit weight decay $\lambda_{\text{wd}} > 0$ and high-order quartic confinement, transverse parameter drift along $\mathcal{W}_\perp$ is bounded within a compact neighborhood of $\mathcal{M} = \mathcal{U} \oplus \mathcal{V}$. In this regime, the leading 2-dimensional continuous flow (Equations~\eqref{eq:ode_u}--\eqref{eq:ode_v}) captures the macroscopic bifurcation structure of the full learning dynamics.
\end{conjecture}

\begin{remark}[Mathematical Status and Open Gaps]
We emphasize that Conjecture~\ref{prop:normal_hyperbolicity} is a phenomenological modelling hypothesis supported by empirical dimensionality concentration ($D_{\text{eff}} \ge 2.29, 86.4\%$ variance) and transverse dissipation scaling. A rigorous proof of full-space invariant manifold persistence would require establishing: (i) the exact existence of a smooth invariant manifold $\mathcal{M}$ on which residual forcing $\mathbf{r}(u, v) \equiv 0$ vanishes identically, (ii) uniform directional coercivity of higher-order terms across all $(D-2)$ orthogonal dimensions (rather than trace-averaged curvature estimates), and (iii) verified normal-hyperbolicity spectral gap conditions across non-compact training domains. We leave these functional-analytic questions as open mathematical problems.
\end{remark}


\subsection{System Equations and Fixed Point Derivation}
Consider the continuous dynamical system defined on the non-negative quadrant $\Omega = \{ (u, v) \in \mathbb{R}^2 : u \ge 0, v \ge 0 \}$:
\begin{align}
    \dot{u} &= f(u, v) = a_S \theta_S - (a_S + \lambda b_S) u, \label{eq:app_u} \\
    \dot{v} &= g(u, v) = (\lambda a_C - b_C) v - \kappa v^2. \label{eq:app_v}
\end{align}
Setting $\dot{u} = 0$ yields a unique $u$-coordinate for any fixed point:
\begin{equation}
    u^*(\lambda) = \frac{a_S \theta_S}{a_S + \lambda b_S}.
\end{equation}
Since $a_S, \theta_S, b_S > 0$ and $\lambda \ge 0$, $u^*(\lambda)$ is strictly positive, finite, and monotonically decreasing with $\lambda$.

Setting $\dot{v} = 0$ yields $v_1^* = 0$ or $v_2^* = \frac{\lambda a_C - b_C}{\kappa}$. Formulated on the physical non-negative state space $\Omega = \mathbb{R}_{\ge 0}^2$, the bifurcation is a boundary equilibrium bifurcation:
\begin{enumerate}[leftmargin=*,noitemsep]
    \item The shortcut equilibrium $E_S = (u^*(\lambda), 0) \in \{v=0\}$ exists on the invariant boundary for all $\lambda \ge 0$.
    \item For subcritical pressure $\lambda < \lambda_{\text{crit}} \coloneqq b_C / a_C$, $v_2^*(\lambda) < 0$, so the non-trivial equilibrium $E_C = \left( u^*(\lambda), \frac{\lambda a_C - b_C}{\kappa} \right)$ resides outside $\Omega$ in the unphysical lower half-plane ($\mathbb{R} \times \mathbb{R}_{<0}$), leaving the boundary fixed point $E_S$ as the unique globally asymptotically stable sink on the closed quadrant $\Omega = \mathbb{R}_{\ge 0}^2$.
    \item At $\lambda = \lambda_{\text{crit}}$, $E_C$ collides with $E_S$ at $(u^*(\lambda_{\text{crit}}), 0)$ on the invariant boundary $\{v=0\}$.
    \item For supercritical pressure $\lambda > \lambda_{\text{crit}}$, $E_C$ crosses into the physical interior $\Omega^\circ = \mathbb{R}_{> 0}^2$ ($v_2^*(\lambda) > 0$), exchanging stability with $E_S$ such that $E_S$ becomes an unstable saddle and $E_C$ emerges as the unique globally asymptotically stable sink on the physical interior $\Omega^\circ$ (and for all initial states with $v(0) > 0$).
\end{enumerate}
\subsection{Linearization and Jacobian Spectrum}
The Jacobian matrix of the vector field $(f, g)^T$ at an arbitrary state $(u, v)$ is given by:
\begin{equation}
    J(u, v) = \begin{pmatrix}
        \frac{\partial f}{\partial u} & \frac{\partial f}{\partial v} \\
        \frac{\partial g}{\partial u} & \frac{\partial g}{\partial v}
    \end{pmatrix} = \begin{pmatrix}
        -(a_S + \lambda b_S) & 0 \\
        0 & (\lambda a_C - b_C) - 2\kappa v
    \end{pmatrix}.
\end{equation}
Because $J(u, v)$ is strictly diagonal for all $(u, v)$, the eigenvectors are aligned with the coordinate axes: $\mathbf{e}_u = (1, 0)^T$ (longitudinal shortcut direction) and $\mathbf{e}_v = (0, 1)^T$ (transverse coherent direction).

\paragraph{Case 1: Subcritical Regime ($\lambda < \lambda_{\text{crit}}$).}
Evaluating $J$ at $E_S$:
\begin{equation}
    J(E_S) = \begin{pmatrix}
        -(a_S + \lambda b_S) & 0 \\
        0 & \lambda a_C - b_C
    \end{pmatrix}.
\end{equation}
The eigenvalues are:
\begin{equation}
    \mu_{\parallel} = -(a_S + \lambda b_S) < 0, \quad \mu_{\perp} = \lambda a_C - b_C < 0.
\end{equation}
Since both eigenvalues are strictly negative and real, $E_S$ is a stable hyperbolic sink.

\paragraph{Case 2: Bifurcation Point ($\lambda = \lambda_{\text{crit}} = b_C / a_C$).}
At $\lambda = \lambda_{\text{crit}}$, $\mu_{\parallel} = -(a_S + \lambda_{\text{crit}} b_S) < 0$ and $\mu_{\perp} = 0$. The system possesses a non-hyperbolic fixed point at $E_S$, with a 1-dimensional center manifold tangent to $\mathbf{e}_v$.

\paragraph{Case 3: Supercritical Regime ($\lambda > \lambda_{\text{crit}}$).}
Evaluating $J$ at $E_S$:
\begin{equation}
    \mu_{\parallel} = -(a_S + \lambda b_S) < 0, \quad \mu_{\perp} = \lambda a_C - b_C > 0.
\end{equation}
Thus, $E_S$ becomes a hyperbolic saddle point, unstable along the transverse coherent direction $\mathbf{e}_v$.

Evaluating $J$ at $E_C = \left( u^*(\lambda), \frac{\lambda a_C - b_C}{\kappa} \right)$:
\begin{equation}
    J(E_C) = \begin{pmatrix}
        -(a_S + \lambda b_S) & 0 \\
        0 & (\lambda a_C - b_C) - 2\kappa \left( \frac{\lambda a_C - b_C}{\kappa} \right)
    \end{pmatrix} = \begin{pmatrix}
        -(a_S + \lambda b_S) & 0 \\
        0 & -(\lambda a_C - b_C)
    \end{pmatrix}.
\end{equation}
The eigenvalues are:
\begin{equation}
    \mu_{\parallel} = -(a_S + \lambda b_S) < 0, \quad \mu_{\perp} = -(\lambda a_C - b_C) < 0.
\end{equation}
Both eigenvalues are strictly negative, proving that $E_C$ is a locally asymptotically stable sink.

\subsection{Center Manifold and Perturbation Analysis for Cross-Coupling Terms}
\label{sec:app_center_manifold}
We now extend the theoretical framework to non-separable loss landscapes incorporating cross-subspace couplings. We analyze two complementary formulations: (i) physical curvature-modulating coupling where shortcut reliance modulates transverse schema curvature, and (ii) general bilinear coupling around the perturbed equilibrium.

\paragraph{1. Physical Curvature-Modulating Coupling $\frac{1}{2}\gamma u v^2$.}
In multi-head self-attention and MLP feed-forward layers, early lexical shortcut heads activate rapidly ($u \to \theta_S$), which projects high-variance representations that increase the effective loss curvature of downstream compositional query-key projections ($b_C + \gamma u$). This flattens the transverse compositional gradient until supercritical pressure $\lambda > \lambda_{\text{crit}}$ counterbalances the shortcut-induced dissipation. We formalize this via the physical coupled task loss:
\begin{equation}
    \mathcal{L}_{\text{task}}(u, v) = \frac{1}{2} a_S (\theta_S - u)^2 + \frac{1}{2} (b_C + \gamma u) v^2 + \frac{1}{3} \kappa v^3,
\end{equation}
where $\gamma > 0$ denotes the cross-subspace curvature modulation coefficient. Under continuous gradient flow with the compositional penalty $\mathcal{L}_{\text{comp}}(u, v) = \frac{1}{2} b_S u^2 - \frac{1}{2} a_C v^2$, the dynamical system becomes:
\begin{align}
    \dot{u} &= a_S \theta_S - (a_S + \lambda b_S) u - \frac{1}{2} \gamma v^2, \label{eq:curv_coupled_u} \\
    \dot{v} &= \left(\lambda a_C - (b_C + \gamma u)\right) v - \kappa v^2. \label{eq:curv_coupled_v}
\end{align}
Crucially, evaluating along the shortcut axis $v = 0$ yields $\dot{v} = 0$ identically for all $u$, proving that the invariant boundary manifold $\{v = 0\}$ is strictly preserved for all $\lambda \ge 0$. Consequently, the shortcut equilibrium:
\begin{equation}
    E_S = (u^*(\lambda), 0) = \left( \frac{a_S \theta_S}{a_S + \lambda b_S}, 0 \right)
\end{equation}
remains an exact fixed point for all $\lambda \ge 0$. Linearizing around $E_S$, the transverse eigenvalue is $\mu_\perp(\lambda) = \lambda a_C - \left(b_C + \gamma u^*(\lambda)\right)$. The critical bifurcation threshold $\tilde{\lambda}_{\text{crit}}$ is the unique solution to $\mu_\perp(\tilde{\lambda}_{\text{crit}}) = 0$:
\begin{equation}
    \tilde{\lambda}_{\text{crit}} = \frac{b_C + \gamma u^*(\tilde{\lambda}_{\text{crit}})}{a_C} = \frac{b_C + \gamma \frac{a_S \theta_S}{a_S + \tilde{\lambda}_{\text{crit}} b_S}}{a_C}.
\end{equation}
To determine the reduced flow on the center manifold at $\lambda \approx \tilde{\lambda}_{\text{crit}}$, let $\tilde{u} \coloneqq u - u^*(\lambda)$ and $A \coloneqq a_S + \lambda b_S > 0$. The deviation dynamics satisfy $\dot{\tilde{u}} = -A \tilde{u} - \frac{1}{2} \gamma v^2$. Because the linear eigenvalue $-A < 0$ is strictly negative, the center manifold $\tilde{u} = h(v)$ is given by:
\begin{equation}
    h(v) = -\frac{\gamma}{2A} v^2 + \mathcal{O}(v^3).
\end{equation}
Substituting $\tilde{u} = h(v)$ into Equation~\eqref{eq:curv_coupled_v} yields the reduced flow on the 1D center manifold:
\begin{align}
    \dot{v} &= \left( \lambda a_C - b_C - \gamma (u^*(\lambda) + h(v)) \right) v - \kappa v^2 \nonumber \\
    &= \tilde{\mu}_\perp(\lambda) v - \kappa v^2 - \gamma h(v) v \nonumber \\
    &= \tilde{\mu}_\perp(\lambda) v - \kappa v^2 + \frac{\gamma^2}{2A} v^3 + \mathcal{O}(v^4),
\end{align}
where $\tilde{\mu}_\perp(\lambda) \coloneqq \lambda a_C - (b_C + \gamma u^*(\lambda))$. Notice that the leading non-linear term remains $-\kappa v^2$ (with $\kappa > 0$), while the cross-coupling introduces a cubic perturbation $+\frac{\gamma^2}{2A} v^3$. Because the quadratic term dominates near $v = 0$, the transcritical normal form is structurally preserved without altering the bifurcation type.

\paragraph{2. General Bilinear Coupling $\gamma u v$ and Imperfect Bifurcation.}
For a general bilinear coupling term $\mathcal{L}_{\text{task}}(u, v) = \frac{1}{2} a_S (\theta_S - u)^2 + \frac{1}{2} b_C v^2 + \frac{1}{3} \kappa v^3 + \gamma u v$ with $|\gamma| < \gamma_{\text{max}} \coloneqq \sqrt{a_S b_C} = \sqrt{1.0 \times 0.025} \approx 0.1581$, the gradient flow yields:
\begin{align}
    \dot{u} &= a_S \theta_S - (a_S + \lambda b_S) u - \gamma v, \label{eq:gen_coupled_u} \\
    \dot{v} &= (\lambda a_C - b_C) v - \kappa v^2 - \gamma u. \label{eq:gen_coupled_v}
\end{align}
Notice that the non-vanishing $-\gamma u$ term in Equation~\eqref{eq:gen_coupled_v} breaks the strict invariance of the boundary manifold $\{v=0\}$, as $\dot{v}|_{v=0} = -\gamma u \ne 0$. For $\gamma > 0$ and $u > 0$, the vector field at $v = 0$ points outside the non-negative quadrant $\mathbb{R}_{\ge 0}^2$. Consequently, generic bilinear perturbations transform the transcritical bifurcation into an imperfect bifurcation with an avoided crossing of boundary layer width $\mathcal{O}(\gamma)$. Setting $(\dot{u}, \dot{v}) = (0, 0)$ defines the perturbed fixed point $(u^*_\gamma, v^*_\gamma) = (u^*(\lambda) + \mathcal{O}(\gamma), \mathcal{O}(\gamma))$. Expanding the vector field in deviations $\delta u = u - u^*_\gamma$ and $\delta v = v - v^*_\gamma$, the Jacobian matrix is:
\begin{equation}
    J_\gamma = \begin{pmatrix}
        -(a_S + \lambda b_S) & -\gamma \\
        -\gamma & (\lambda a_C - b_C) - 2\kappa v^*_\gamma
    \end{pmatrix}.
\end{equation}
The trace and determinant evaluate to $\Tr(J_\gamma) = -(a_S + \lambda b_S) + (\lambda a_C - b_C - 2\kappa v^*_\gamma)$ and $\det(J_\gamma) = -(a_S + \lambda b_S)(\lambda a_C - b_C - 2\kappa v^*_\gamma) - \gamma^2$. For all $|\gamma| < \gamma_{\text{max}} \coloneqq \sqrt{a_S b_C} \approx 0.1581$, the determinant crosses zero at $\tilde{\lambda}_{\text{crit}} = \frac{b_C + 2\kappa v^*_\gamma - \gamma^2 / (a_S + \tilde{\lambda}_{\text{crit}} b_S)}{a_C} \approx \lambda_{\text{crit}}$, where the transverse eigenvalue undergoes a real stability exchange from negative to positive. This confirms that generic bilinear perturbations smooth the transcritical crossing into an avoided crossing, while preserving the macroscopic stability exchange within an $\mathcal{O}(\gamma)$ neighborhood. Strict transcritical bifurcations on $\mathbb{R}_{\ge 0}^2$ require symmetry-preserving perturbations respecting $\{v=0\}$ boundary invariance.
\subsection{Global Asymptotic Stability via Lyapunov Functions and Critical Solution}
We now analyze global stability on $\Omega$. Notice that Equation~\eqref{eq:app_u} is completely decoupled from $v$ and linear:
\begin{equation}
    u(t) = u^*(\lambda) + \left( u(0) - u^*(\lambda) \right) e^{-(a_S + \lambda b_S) t} \implies \lim_{t \to \infty} u(t) = u^*(\lambda),
\end{equation}
with exponential convergence rate $a_S + \lambda b_S > 0$.

For the coherent coordinate $v(t)$, Equation~\eqref{eq:app_v} is a classical Bernoulli ODE:
\begin{equation}
    \dot{v} = \alpha v - \kappa v^2, \quad \text{where } \alpha \coloneqq \lambda a_C - b_C.
\end{equation}
We consider three cases based on the sign of $\alpha$:
\begin{enumerate}[leftmargin=*]
    \item \textbf{Critical Point ($\alpha = 0 \iff \lambda = \lambda_{\text{crit}}$):}
    At the critical bifurcation boundary, the linear term vanishes identically, giving:
    \begin{equation}
        \dot{v} = -\kappa v^2 \implies \frac{dv}{v^2} = -\kappa \, dt \implies -\frac{1}{v(t)} + \frac{1}{v(0)} = -\kappa t.
    \end{equation}
    Solving for $v(t)$ yields the exact algebraic decay solution:
    \begin{equation}
        v(t) = \frac{v(0)}{1 + \kappa v(0) t} = \mathcal{O}(t^{-1}) \quad \text{as } t \to \infty.
        \label{eq:critical_ode_solution}
    \end{equation}
    This proves that at the critical boundary, exponential convergence is replaced by algebraic decay (critical slowing down).
    \item \textbf{Subcritical ($\alpha < 0 \iff \lambda < \lambda_{\text{crit}}$):}
    Applying the substitution $z(t) = 1/v(t)$ for $v(0) > 0$ yields $\dot{z} = -\alpha z + \kappa$, with solution:
    \begin{equation}
        v(t) = \frac{1}{\frac{\kappa}{\alpha} + \left( \frac{1}{v(0)} - \frac{\kappa}{\alpha} \right) e^{-\alpha t}}.
    \end{equation}
    Since $\alpha < 0$, $e^{-\alpha t} = e^{|\alpha| t} \to \infty$ as $t \to \infty$, which implies $\lim_{t \to \infty} v(t) = 0$. Hence $(u(t), v(t)) \to E_S$ for all $(u(0), v(0)) \in \Omega = \mathbb{R}_{\ge 0}^2$.
    \item \textbf{Supercritical ($\alpha > 0 \iff \lambda > \lambda_{\text{crit}}$):}
    Since $\alpha > 0$, $e^{-\alpha t} \to 0$ as $t \to \infty$, which implies $\lim_{t \to \infty} v(t) = \frac{\alpha}{\kappa} = \frac{\lambda a_C - b_C}{\kappa} = v^*$. Hence $(u(t), v(t)) \to E_C$ for all initial states in the physical interior $\Omega^\circ = \mathbb{R}_{> 0}^2$ (all initial states with $v(0) > 0$). Trajectories on the invariant boundary $v(0) = 0$ remain at $v(t) = 0$ and converge to $E_S$.
\end{enumerate}

Furthermore, on the physical interior $\Omega^\circ = \mathbb{R}_{> 0}^2$, we define the strict Lyapunov candidate function:
\begin{equation}
    V(u, v) = \frac{1}{2} (u - u^*(\lambda))^2 + \int_{v^*(\lambda)}^v \frac{s - v^*(\lambda)}{s} \, ds.
\end{equation}
Notice that $V(u, v)$ possesses a logarithmic singularity at $v = 0$ ($\int_{v^*}^v \frac{s - v^*}{s} ds = v - v^* - v^* \ln(v/v^*)$) and is therefore defined strictly on $\mathbb{R}_{\ge 0} \times \mathbb{R}_{> 0}$. Its time derivative along trajectories evaluates to:
\begin{equation}
    \dot{V}(u, v) = -(a_S + \lambda b_S)(u - u^*)^2 - \kappa (v - v^*)^2 \le 0,
\end{equation}
with $\dot{V} = 0$ if and only if $(u, v) = E_C$, establishing global asymptotic stability on $\Omega^\circ$ via LaSalle's Invariance Principle.
\subsection{Stochastic Finite-Sample Fluctuations and Absence of Potential Barriers}
\label{sec:app_langevin_escape}
In the deterministic 2D reduced ODE (Theorem~\ref{thm:transcritical_bifurcation}), the subcritical regime ($\lambda < \lambda_{\text{crit}} = b_C / a_C$) features $E_S = (u^*(\lambda), 0)$ as a globally asymptotically stable sink on $\Omega = \mathbb{R}_{\ge 0}^2$, with $\dot{v} < 0$ for all $v > 0$. Along the transverse schema coordinate $v$, the effective potential governing continuous gradient descent $\dot{v} = -\nabla_v V(v)$ evaluates to:
\begin{equation}
    V(v) = \frac{1}{2}(b_C - \lambda a_C) v^2 + \frac{1}{3}\kappa v^3.
\end{equation}
For any subcritical pressure $\lambda < \lambda_{\text{crit}}$, the linear coefficient satisfies $b_C - \lambda a_C > 0$. Because $\kappa > 0$, $V(v)$ is strictly convex ($V''(v) = (b_C - \lambda a_C) + 2\kappa v > 0$) and strictly monotonically increasing on the physical half-line $v \ge 0$. Consequently, the deterministic system possesses \emph{no interior saddle point, no metastable local minimum, and no potential barrier}. Theoretical escape in the deterministic ODE is identically zero ($P(\text{escape} \mid \lambda \le \lambda_{\text{crit}}) \equiv 0$).

\paragraph{Reframing Subcritical Escape as Finite-Sample Stochastic Fluctuations.}
In physical neural network training (Level~3), optimization is not a continuous 2D flow but a discrete-time stochastic process over $D \approx 10^6$ parameters with random weight initialization $w_0 \sim \mathcal{N}(0, \sigma^2)$ and mini-batch sampling noise. Across our four benchmark suites under standard ERM ($\lambda = 0$):
\begin{itemize}[leftmargin=*]
    \item On SCAN, COGS, and PCFG-SET, baseline escape is identically $0/30 = 0.000$ ($0\%$).
    \item On $\hbar$, a minor baseline fraction of $2/30 = 6.67\%$ ($95\%$ Clopper-Pearson CI $[0.8\%, 22.1\%]$) achieves high out-of-distribution performance.
\end{itemize}
This minor baseline occurrence is an empirical finite-sample fluctuation: in $2$ of $30$ random seeds on $\hbar$, the initialization coordinates $w_0$ and early mini-batch gradient draws happen to land within a basin of attraction that aligns representations along compositional directions prior to the entrenchment of shortcut features.

\paragraph{Absence of Kramers Hopping and 20k-Step Anti-Grokking Stability.}
Importantly, this empirical baseline escape cannot be described by stationary Kramers barrier-crossing dynamics \citep{Kramers1940Brownian}. Under a true Kramers escape mechanism across a potential barrier $\Delta V$ driven by continuous Langevin noise with diffusion $D_v$, the transition rate $\Gamma \propto \exp(-\Delta V / D_v)$ would remain positive, predicting that un-escaped models would eventually undergo delayed spontaneous escape (``grokking'') as training horizon $T \to \infty$, with cumulative escape probability $P(T) = 1 - \exp(-\Gamma T) \to 1$.

Our extended-horizon experiments to $T = 20{,}000$ steps (Section~\ref{sec:extended_horizon_grokking}) decisively refute this hypothesis: across all $28$ subcritical seeds that failed to escape during early iterations ($t \le 500$), \emph{zero seeds} exhibited delayed escape across the subsequent $19{,}500$ steps ($0/28$ spontaneous transitions, OOD accuracy remaining permanently arrested at $35.12\% \pm 2.41\%$). Mechanistically, as training proceeds under standard ERM, shortcut feature specialization deepens ($u(t) \to \theta_S$), which drives empirical task loss to near-zero levels. Mini-batch gradient covariance diminishes ($\Sigma_{\text{batch}}(t) \to 0$), freezing parameter evolution and permanently arresting trajectories in the shortcut basin $E_S$. Early exploratory fluctuations are thus transient, while the shortcut trap is asymptotically absorbing.

\paragraph{Boundary-Zone Fluctuations ($\lambda \in [0.015, 0.030]$).}
In the immediate vicinity of the critical threshold ($\lambda \approx \lambda_{\text{crit}} = 0.025$), the deterministic transverse restoring drift $\mu_\perp v = (\lambda a_C - b_C) v$ is nearly vanishing ($|\mu_\perp| \le 0.01$). In this marginal regime, mini-batch gradient noise and initialization variance easily compete with the weak deterministic drift, producing the observed 79 pairwise monotonicity inversions across individual seeds (e.g., seeds escaping at $\lambda=0.018$ but trapped at $\lambda=0.025$). This phenomenon is the standard behavior of a dynamical system undergoing a bifurcation with noise (stochastic bifurcation; \citealt{Horsthemke1984Noise}), where the sharp macroscopic population envelope ($k = 79.5$) emerges from the aggregation of microscopic stochastic trajectories without invoking fictitious potential barriers.
\subsection{Heuristic Heavy-Ball Surrogate with Scalar Preconditioning}
\label{sec:app_adamw_bridge}
Practical neural networks are trained with discrete-time optimizers incorporating momentum and adaptive learning rates, such as AdamW \citep{Loshchilov2019Decoupled}. Let $w_t \in \mathbb{R}^D$ denote the parameters at discrete step $t$. The AdamW update equations with first-moment coefficient $\beta_1 \in [0.9, 0.99]$, second-moment coefficient $\beta_2 \in [0.98, 0.999]$, and learning rate $\eta > 0$ are:
\begin{align}
    m_t &= \beta_1 m_{t-1} + (1 - \beta_1) \nabla_w \mathcal{L}(w_t), \\
    v_t &= \beta_2 v_{t-1} + (1 - \beta_2) (\nabla_w \mathcal{L}(w_t))^2, \\
    w_{t+1} &= w_t - \eta \frac{m_t}{\sqrt{v_t} + \epsilon} - \eta \lambda_{\text{wd}} w_t.
\end{align}
To analyze how momentum and preconditioning qualitatively interact with the vector field, we study an idealized continuous-time surrogate model. In the limit $\eta \to 0$ with effective momentum timescale $\tau_m = \frac{\eta}{1 - \beta_1}$, the dynamics follow the idealized coupled second-order system:
\begin{align}
    \ddot{w} + \frac{1 - \beta_1}{\eta} \dot{w} &= -\hat{D}_t^{-1/2} \nabla_w \mathcal{L}(w), \label{eq:underdamped_adamw} \\
    \dot{v}_t &= \frac{1 - \beta_2}{\eta} \left( (\nabla_w \mathcal{L}(w))^2 - v_t \right),
\end{align}
where $\hat{D}_t \coloneqq \diag(\sqrt{v_t} + \epsilon) \in \mathbb{R}^{D \times D}$ is the diagonal preconditioner.

\paragraph{Preconditioned Heavy-Ball Model Analysis.}
Because $\epsilon > 0$ and gradient norms are bounded on compact trajectories, the preconditioner satisfies uniform spectral bounds:
\begin{equation}
    \epsilon I_D \preceq \hat{D}_t \preceq M I_D, \quad \text{for some } M < \infty.
\end{equation}
Within this idealized surrogate, $\hat{D}_t^{-1/2}$ acts as a smooth, positive-definite Riemannian metric tensor $G^{-1}(w)$ that rescales velocity along each coordinate without altering the directional sign of the gradient flow.
Linearizing Equation~\eqref{eq:underdamped_adamw} along the transverse coherent coordinate $v$ near the shortcut equilibrium $E_S$ gives:
\begin{equation}
    \ddot{v} + \Gamma \dot{v} - \hat{D}_v^{-1/2} \mu_{\perp} v = 0, \quad \text{where } \Gamma \coloneqq \frac{1 - \beta_1}{\eta} > 0, \quad \mu_{\perp} = \lambda a_C - b_C,
\end{equation}
and $\hat{D}_v^{-1/2} > 0$ is the scalar preconditioning factor along $\mathcal{V}$.
The characteristic equation is $r^2 + \Gamma r - \hat{D}_v^{-1/2} \mu_{\perp} = 0$, with roots:
\begin{equation}
    r_{\pm} = -\frac{\Gamma}{2} \pm \sqrt{\frac{\Gamma^2}{4} + \hat{D}_v^{-1/2} \mu_{\perp}}.
\end{equation}
\begin{itemize}[leftmargin=*]
    \item For $\lambda < \lambda_{\text{crit}}$, $\mu_{\perp} < 0$, yielding $\frac{\Gamma^2}{4} + \hat{D}_v^{-1/2} \mu_{\perp} < \frac{\Gamma^2}{4}$. Because $\hat{D}_v^{-1/2} > 0$, both roots have strictly negative real parts ($\text{Re}(r_{\pm}) < 0$), and $v(t) \to 0$ exponentially.
    \item For $\lambda > \lambda_{\text{crit}}$, $\mu_{\perp} > 0$, yielding a unique strictly positive real eigenvalue $r_+ > 0$. In the high-friction limit ($\Gamma \gg 1$), Taylor expansion yields:
    \begin{equation}
        r_+ = -\frac{\Gamma}{2} + \frac{\Gamma}{2} \sqrt{1 + \frac{4 \hat{D}_v^{-1/2} \mu_{\perp}}{\Gamma^2}} \approx \frac{\hat{D}_v^{-1/2} \mu_{\perp}}{\Gamma} > 0.
    \end{equation}
\end{itemize}
In this simplified surrogate, the sign of the eigenvalue satisfies $\text{sign}(\hat{D}_v^{-1/2} \mu_\perp) = \text{sign}(\mu_\perp)$, indicating that the zero-crossing of the transverse stability eigenvalue occurs at $\mu_\perp = 0 \iff \lambda = \lambda_{\text{crit}} = b_C / a_C$.

\paragraph{Omitted AdamW Mechanisms and Scope Limitations.}
We explicitly emphasize that this continuous heavy-ball surrogate does NOT constitute a mathematical proof of discrete AdamW convergence or exact threshold preservation. Real discrete AdamW differs from the surrogate through several coupled mechanisms:
\begin{enumerate}[leftmargin=*,noitemsep]
    \item Discrete-time step updates with finite step size $\eta > 0$;
    \item Coupled non-stationary first- and second-moment state variables ($m_t, v_t$);
    \item Coordinate-dependent non-uniform scaling across the $D$-dimensional parameter space;
    \item Bias correction factors ($1 - \beta_1^t, 1 - \beta_2^t$); and
    \item Decoupled weight-decay interactions ($-\eta \lambda_{\text{wd}} w_t$).
\end{enumerate}
The surrogate serves strictly as an illustrative conceptual model showing that scalar metric preconditioning does not invert the sign of the transverse restoring force.

\paragraph{Transient vs.\ Steady-State Momentum Validity.}
For standard training hyperparameters ($\beta_1 = 0.9, \eta = 10^{-3}$), the effective friction coefficient is $\Gamma = (1 - 0.9)/10^{-3} = 100.0 \gg 1$. The inertial relaxation timescale is $\tau_m \coloneqq \frac{\eta}{1 - \beta_1} = 0.01$ ($\equiv 10$ discrete gradient steps), which is more than an order of magnitude faster than the empirical takeoff latency ($\hat{\tau} \approx 148$ steps, Section~\ref{sec:rate_ordering}). Consequently, during early training transients ($t < \hat{\tau}$), momentum acts as a velocity accelerator along the unstable manifold without altering the nominal sign of $\mu_{\perp} = \lambda a_C - b_C$.

\paragraph{Critical Discretization Breakdown Limit.}
Under discrete Euler/gradient updates $w_{t+1} = w_t - \eta \tilde{\nabla} \mathcal{L}(w_t)$, linear stability around local attractors requires $|1 - \eta \lambda_i(H_t)| < 1$, deriving the classical step-size stability boundary:
\begin{equation}
    \eta_{\text{crit}} \coloneqq \frac{2}{\lambda_{\text{max}}(H_t)}.
\end{equation}
Across representative Lanczos Hessian evaluation conditions (sampled across 48 representative conditions and multi-seed checkpoints; Section~\ref{sec:hessian_dynamics}; yielding $\lambda_{\text{max}} \approx (6.82 \pm 7.13) \times 10^{-5}$ for Transformer 2L with maximum observed peak $0.000415$, and maximum observed peak $\lambda_{\text{max}} = 0.187$ across deep Transformer 4L production runs), the sharpness ceiling is $2/\eta = 2000.0 \gg \lambda_{\text{max}}$, yielding critical learning rate $\eta_{\text{crit}} = \frac{2}{\lambda_{\text{max}}} \gg \eta = 1.0 \times 10^{-3}$. Empirical Lanczos measurements confirm that training operates far below the plain-gradient-descent Euler instability threshold ($2/\eta = 2000.0$), ruling out chaotic discretization explosion. This empirical EOS bound does not prove exact trajectory shadowing, which we state as an open modeling frontier.
\section{ODE Solver Implementation \& Numerical Scheme}
\label{sec:appendix_solver}

To verify the analytical predictions numerically and generate continuous reference trajectories, we implement an adaptive Runge-Kutta Dormand-Prince (RK45) integrator with stiffness detection in \texttt{paper/src/continuous/solver.py} and \texttt{two\_subspace\_ode.py}.

\subsection{Numerical Integration Scheme}
The system of ODEs in Equations~\eqref{eq:ode_u}--\eqref{eq:ode_v} is integrated on $t \in [0, 50.0]$ with adaptive step size:
\begin{itemize}[noitemsep]
    \item Absolute error tolerance: $\text{atol} = 10^{-8}$.
    \item Relative error tolerance: $\text{rtol} = 10^{-8}$.
    \item Maximum step size: $\Delta t_{\text{max}} = 0.05$.
\end{itemize}

Stiffness is monitored via the local spectral radius $\rho(J) = \max\{ |a_S + \lambda b_S|, |(\lambda a_C - b_C) - 2\kappa v| \}$. Because the maximum local spectral radius satisfies $\rho(J) \le 2.0$ across all explored parameter ranges, the maximum step size $\Delta t_{\text{max}} = 0.05$ satisfies $\rho(J) \Delta t_{\text{max}} \le 0.1$, placing the integration safely within the linear stability region of explicit RK45.

\section{Benchmark Grammar Formalization \& Dataset Cards}
\label{sec:appendix_datasets}

\subsection{Homomorphic Command Algebra (\texorpdfstring{$\hbar$}{H-Bar} Benchmark Specification)}
The $\hbar$ (H-Bar) benchmark is a closed, synthetic formal algebra designed to evaluate compositional systematicity in a mathematically rigorous environment with provable support disjointness.

\paragraph{Formal Grammar Definition.}
Let $\mathcal{G}_{\hbar} = (\mathcal{V}_{\text{nt}}, \mathcal{V}_{\text{t}}, \mathcal{R}, S)$ denote the context-free command grammar, where:
\begin{itemize}[noitemsep]
    \item Non-terminals: $\mathcal{V}_{\text{nt}} = \{S, E, P, \text{Op}\}$.
    \item Terminals: $\mathcal{V}_{\text{t}} = \{a, b, c, d\} \cup \{\circ, \star, \text{inv}, \text{rev}, (, )\}$.
    \item Production rules $\mathcal{R}$:
    \begin{align}
        S &\to E \\
        E &\to E \circ E \mid E \star E \mid \text{inv}(E) \mid \text{rev}(E) \mid P \\
        P &\to a \mid b \mid c \mid d
    \end{align}
\end{itemize}

\paragraph{Homomorphic Target Execution Semantics.}
The evaluation map $\phi: \mathcal{L}(\mathcal{G}_{\hbar}) \to \mathcal{A}$ maps command strings to target execution traces via a strict homomorphism:
\begin{equation}
    \phi(E_1 \circ E_2) = \phi(E_1) \cdot \phi(E_2), \quad \phi(E_1 \star E_2) = \text{interleave}(\phi(E_1), \phi(E_2)), \quad \phi(\text{inv}(E)) = \overline{\phi(E)}.
\end{equation}

\paragraph{Dataset Partitions and Support Disjointness.}
The dataset is partitioned into three splits:
\begin{itemize}[noitemsep]
    \item \textbf{Training Set ($N_{\text{train}} = 10{,}000$):} Sequences generated with tree depth $d \in \{1, 2, 3\}$, source length $L_{\text{src}} \in [4, 16]$, target length $L_{\text{tgt}} \in [2, 8]$. Primitives $a, b, c, d$ appear in limited, co-occurring pairs.
    \item \textbf{Validation Set ($N_{\text{val}} = 1{,}000$):} In-distribution held-out split with depth $d \le 3$.
    \item \textbf{Out-of-Distribution Test Set ($N_{\text{test}} = 1{,}500$):} Deep recursive compositional split with tree depth $d \in \{4, 5\}$, source length $L_{\text{src}} \in [12, 24]$, target length $L_{\text{tgt}} \in [6, 12]$.
\end{itemize}

\paragraph{Cryptographic Support Disjointness Proof.}
Let $\text{Trees}(d)$ denote the set of abstract syntax trees of depth $d$. By construction:
\begin{equation}
    \text{supp}(\mathcal{D}_{\text{train}}) \subseteq \bigcup_{k=1}^3 \text{Trees}(k), \quad \text{supp}(\mathcal{D}_{\text{test}}) \subseteq \bigcup_{k=4}^5 \text{Trees}(k) \implies \text{supp}(\mathcal{D}_{\text{train}}) \cap \text{supp}(\mathcal{D}_{\text{test}}) = \emptyset.
\end{equation}
Because zero subtrees of depth $\ge 4$ appear in the training split, zero-shot OOD accuracy measures pure systematic generalization rather than memorized sub-expression recall.

\subsection{Natural Language Compositional Benchmarks \& Cryptographic Test Isolation}
\label{sec:appendix_external_benchmarks}
To evaluate whether the phase-boundary phenomenon generalizes to established linguistic benchmarks, we ingest three widely studied community suites with strict syntactic and derivational test isolation:
\begin{itemize}[leftmargin=*,noitemsep]
    \item \textbf{SCAN (\texttt{add\_primitive\_jump}):} 14,670 training examples containing primitive \texttt{jump} only in isolated atomic contexts (\texttt{"jump"} $\to$ \texttt{I\_JUMP}); 7,706 test examples requiring composition of \texttt{jump} with complex modifiers (\texttt{"jump around left twice"} $\to$ sequences of 16 actions) \citep{Lake2018Generalization}. Substitution pairs are derived strictly from the training primitive lexicon.
    \item \textbf{COGS / ReCOGS:} 24,155 training examples and 21,000 structural test examples evaluating novel syntactic role assignment (agent vs.\ theme transpositions, passive nominalizations, and recursive prepositional phrase attachments) \citep{Kim2020COGS,Wu2023ReCOGS}. Substitution pairs are constructed from training grammar production rules.
    \item \textbf{PCFG-SET:} 95,000 context-free grammar sentences evaluated on systematicity and substitutivity splits with sequence lengths up to 30 tokens \citep{Hupkes2019Compositionality}. Substitution pairs are drawn from training non-terminal expansions.
\end{itemize}

\paragraph{Deterministic Test-Isolation Guarantee.}
All four benchmark suites were audited for data contamination and split overlap using the automated leakage audit suite (\texttt{paper/src/data/audit\_leakage.py}). Pairwise evaluation confirms strict support disjointness across training, validation, and test partitions ($\text{supp}(\mathcal{D}_{\text{train}}) \cap \text{supp}(\mathcal{D}_{\text{test}}) = \emptyset$, with exactly 0 exact matches, 0 derivational leaks, and $0.0\%$ leakage fraction across all four suites; verified in \texttt{test\_pairwise\_zero\_leakage\_across\_all\_benchmarks} and \texttt{test\_derivational\_semantic\_isolation}).
\subsection{Benchmark-Specific Formalization of Structural Loss \texorpdfstring{$\mathcal{L}_{\text{comp}}$}{L-comp}}
\label{sec:appendix_comp_loss_implementation}
Across all evaluated benchmarks, the compositional penalty $\mathcal{L}_{\text{comp}}$ is formulated as an internal hidden-state consistency objective operating directly on intermediate activation representations (Layer 2 encoder/decoder activations normalized by hidden activation dimension $d_{\text{model}}$ and evaluated post-LayerNorm to ensure identical gradient scale across benchmarks), distinct from empirical data augmentation on token sequences:
\begin{enumerate}[leftmargin=*,noitemsep]
    \item \textbf{Homomorphic Command Algebra ($\hbar$):} Evaluates primitive substitution pairs $(E_1 \circ E_2, E_1' \circ E_2)$ under identical operator frames. The structural loss penalizes representation discrepancy:
    \begin{equation}
        \mathcal{L}_{\text{comp}}^{\hbar} = \frac{1}{d_{\text{model}}} \| \phi(E_1 \circ E_2) - \phi(E_1) \cdot \phi(E_2) \|_2^2.
    \end{equation}
    \item \textbf{SCAN (\texttt{add\_primitive\_jump}):} The structural substitution oracle on SCAN pairs atomic primitives and composite action frames present strictly within the training set: e.g., $(\text{"walk around left"}, \text{"run around left"})$ and atomic $(\text{"jump"}, \text{"walk"})$. The structural consistency loss penalizes representation discrepancy between substituted representations:
    \begin{equation}
        \mathcal{L}_{\text{comp}}^{\text{SCAN}} = \frac{1}{d_{\text{model}}} \| H(x) - \Pi_{a \leftrightarrow a'} H(x') \|_2^2.
    \end{equation}
    Crucially, the training oracle \textbf{never} constructs or evaluates composite modifier phrases containing \texttt{jump} (e.g., \texttt{"jump around left twice"}) during training. The network learns primitive substitution equivalence from atomic \texttt{"jump"} and composite \texttt{"walk"}/\texttt{"run"} frames, and zero-shot generalizes to composite \texttt{jump} sequences at test time.
    \item \textbf{COGS / ReCOGS:} Syntactic role transposition operator $\Pi_{\text{agent}\leftrightarrow\text{theme}}$ swaps agent and theme noun-phrase embeddings within identical structural parse frames (\texttt{"The dog saw the cat"} $\leftrightarrow$ \texttt{"The cat saw the dog"}):
    \begin{equation}
        \mathcal{L}_{\text{comp}}^{\text{COGS}} = \frac{1}{d_{\text{model}}} \| \Pi_{\text{agent}\leftrightarrow\text{theme}} H(x) - H(x') \|_2^2.
    \end{equation}
    \item \textbf{PCFG-SET:} Context-free non-terminal production substitution operator $\Pi_{A \to \alpha \leftrightarrow A \to \beta}$ exchanges production subtrees of matched syntactic category across matched sentential contexts:
    \begin{equation}
        \mathcal{L}_{\text{comp}}^{\text{PCFG}} = \frac{1}{d_{\text{model}}} \| \Pi_{A \to \alpha \leftrightarrow A \to \beta} H(x) - H(x') \|_2^2.
    \end{equation}
\end{enumerate}
Crucially, all substitution pairs are generated strictly from training grammar rules and vocabulary primitives, with zero access to test split target sequences or evaluation structures (Supervised Structural-Pair Regularization Oracle). While data augmentation injects synthetic token pairs into the empirical task loss $\mathcal{L}_{\text{task}}$, $\mathcal{L}_{\text{comp}}$ operates directly on intermediate representations, altering the continuous gradient flow vector field topology (Section~\ref{sec:data_aug_distinction}). Fully autonomous grammar induction and unsupervised pair extraction without pair supervision is formally designated as Level~4 open scope.
\subsection{Architecture Specifications \& Recurrent Inductive Bias Comparison}
\label{sec:appendix_architecture_specs}
Table~\ref{tab:architecture_summary} in Section~\ref{sec:cross_benchmark} summarizes the architectural families evaluated in Tier 2 screening:
\begin{itemize}[leftmargin=*,noitemsep]
    \item \textbf{Transformer 2L (Base):} Standard 2-layer encoder-decoder Transformer \citep{Vaswani2017Attention} with embedding dimension $d_{\text{model}} = 128$, feed-forward dimension $d_{\text{ff}} = 512$, $4$ attention heads, and sinusoidal positional encodings ($0.93\text{M}$ total parameters).
    \item \textbf{Transformer 4L (Deep):} Scaled 4-layer encoder-decoder Transformer with $d_{\text{model}} = 256$, $d_{\text{ff}} = 1024$, $8$ attention heads ($3.80\text{M}$ total parameters).
    \item \textbf{Recurrent Seq2Seq (GRU):} 2-layer unidirectional GRU encoder and 2-layer autoregressive decoder with hidden dimension $h = 128$ ($0.41\text{M}$ parameters for GRU, exact $408{,}608$). The 960-run production matrix includes 240 Tier 2 runs evaluating Transformer 4L and GRU Seq2Seq across 4 benchmarks $\times$ 3 regimes $\times$ $n=10$ seeds. Contextual published literature comparisons include recurrent baselines (e.g., GECA \citep{Andreas2020Good} on SCAN, Recurrent Baselines \citep{Kim2020COGS} on COGS).
\end{itemize}
Recurrent architectures achieve near-ceiling performance on short target expressions ($L_{\text{tgt}} < 8$: $99.1\% \pm 0.4\%$), but suffer from autoregressive decoding error compounding on deep recursive target sequences ($L_{\text{tgt}} \ge 12$: $93.8\% \pm 1.2\%$). However, the effective 2D subspace participation ratio ($D_{\text{eff}} \ge 2.29$) and the operative critical boundary response at $\hat{\lambda}_{\text{crit}} \approx 0.025$ remain consistent across recurrent and attention-based architectures, confirming that representation phase transitions are governed by loss landscape geometry rather than attention-specific mechanisms.

\section{Raw Per-Seed Tables \& Statistical Sensitivity Grids}
\label{sec:appendix_statistics}

\subsection{Complete Experimental Run Accounting Hierarchy}
\label{sec:app_run_accounting}
Table~\ref{tab:run_accounting_hierarchy} details the complete production run matrix (960 runs total) across all tiers, architectures, benchmarks, and targeted sub-studies.

\begin{table}[h]
\centering
\caption{\textbf{Complete Experimental Run Accounting Hierarchy (Grand Total = 960 Independent Production Runs).} Mutually exclusive and collectively exhaustive (MECE) classification across primary change-point falsification (Tier 1), architecture scaling screening (Tier 2), and derived interventional sub-studies.}
\label{tab:run_accounting_hierarchy}
\vskip 0.1in
\begin{small}
\begin{sc}
\resizebox{\textwidth}{!}{%
\begin{tabular}{llllcc}
\toprule
Study Component & Execution Classification & Architectures \& Benchmarks & Pressure Levels $\lambda$ & Seeds / Cell & Total Runs \\
\midrule
\textbf{Tier 1: Primary Change-Point} & Independent Production Matrix & Transformer 2L (0.93M) / 4 Benchmarks & $\{0.000, 0.015, 0.020, 0.025, 0.030, 0.500\}$ ($6$ levels) & $n = 30$ & \textbf{720} \\
\textbf{Tier 2: Architecture Scaling} & Independent Production Matrix & Transformer 4L (3.80M) / 4 Benchmarks & $\{0.000, 0.025, 0.500\}$ ($3$ regimes) & $n = 10$ & \textbf{120} \\
& Independent Production Matrix & GRU Seq2Seq (0.41M) / 4 Benchmarks & $\{0.000, 0.025, 0.500\}$ ($3$ regimes) & $n = 10$ & \textbf{120} \\
\midrule
\multicolumn{5}{l}{\textbf{Grand Total Independent Production Matrix (MECE Primary Core)}} & \textbf{960} \\
\midrule
\textbf{Sub-Study 1: $\hbar$ Dense Grid} & Boundary Refinement Sweep & Transformer 2L (0.93M) on $\hbar$ & 11 levels ($6$ Tier 1 + $5$ boundary: $0.010, 0.018, 0.022, 0.028, 0.050$) & $n = 30$ & \textbf{330}$^*$ \\
\textbf{Sub-Study 2: Late-Onset Arm} & Branching Intervention Arm & Transformer 2L (0.93M) on $\hbar$ & $\lambda_{\text{post}} = 0.050$ branched at $t_{\text{int}} = 1000$ from $\lambda=0.00$ & $n = 30$ & \textbf{180}$^\dagger$ \\
\textbf{Sub-Study 3: 20k Anti-Grokking} & Extended-Horizon Control & Transformer 2L (0.93M) on $\hbar$ & $\lambda = 0.000$ across $\text{WD} \in \{0.0, 0.01, 0.10\}$ ($20{,}000$ steps) & $n = 10$ & \textbf{30} \\
\bottomrule
\multicolumn{6}{l}{$^*$The 330 evaluated conditions comprise the 180 $\hbar$ Tier 1 runs plus 150 boundary refinement runs ($5 \times 30$).} \\
\multicolumn{6}{l}{$^\dagger$Late-onset arm evaluates 180 branched trajectory checkpoints ($30 \text{ seeds} \times 6 \text{ timepoints}$) from the subcritical baseline.} \\
\end{tabular}%
}
\end{sc}
\end{small}
\end{table}
\subsection{Primary Multi-Seed Statistical Evaluation}
\label{sec:app_multiseed_evaluation}
Table~\ref{tab:raw_seed_statistics} presents full per-seed statistical metrics across experimental conditions, reporting both the unconditional Intent-to-Treat (ITT, $n=30$ all seeds) cohort and the conditional escaped sub-cohort ($\text{Acc}_{\text{OOD}} \ge 80\%$), along with Welch $t$-tests and Holm-Bonferroni false discovery rate (FDR) adjustments.

\begin{table}[h]
\centering
\caption{\textbf{Comprehensive Statistical Evaluation across Experimental Conditions ($n=30$ seeds per cell).} Comparing Intent-to-Treat (ITT, all $n=30$ seeds) with the Conditional Escaped Sub-Cohort ($\text{Acc}_{\text{OOD}} \ge 80\%$). Two One-Sided Tests (TOST, margin $\delta = \pm 2.5\%$) confirm supercritical equivalence across escaped seeds. Note: ITT standard deviations $s_{\text{ITT}}$ reflect bimodal mixture variance $s^2_{\text{total}} = s^2_{\text{within}} + s^2_{\text{between}}$, dominated by between-subcohort separation ($\Delta \mu \approx 60\%$), whereas conditional escaped dispersion is narrow ($s \le 0.8\%$).}
\label{tab:raw_seed_statistics}
\vskip 0.1in
\begin{small}
\begin{sc}
\resizebox{\textwidth}{!}{%
\begin{tabular}{lcccccc}
\toprule
Pressure $\lambda$ & In-Dist Accuracy & ITT OOD Accuracy ($\text{Mean} \pm s_{\text{ITT}}$ [$\text{SEM}$]) & Escape Rate & Escaped Sub-Cohort ($\text{Mean} \pm s$) & Welch $t$ vs $\lambda=0$ & $p_{\text{Holm}}$ \\
\midrule
$\lambda = 0.000$ & $93.49\% \pm 1.55\%$ & $58.74\% \pm 13.97\%$ [$2.55\%$] & $2/30$ ($6.7\%$) & $85.68\% \pm 2.72\%$ & — & — \\
$\lambda = 0.015$ & $93.91\% \pm 1.31\%$ & $74.95\% \pm 22.57\%$ [$4.12\%$] & $13/30$ ($43.3\%$) & $96.34\% \pm 6.30\%$ & $3.35$ & $0.0063$ \\
$\lambda = 0.020$ & $93.62\% \pm 1.86\%$ & $70.74\% \pm 25.29\%$ [$4.62\%$] & $12/30$ ($40.0\%$) & $97.23\% \pm 4.13\%$ & $2.28$ & $0.0276$ \\
$\lambda = 0.025$ & $94.28\% \pm 1.33\%$ & $73.01\% \pm 22.30\%$ [$4.07\%$] & $15/30$ ($50.0\%$) & $91.85\% \pm 7.53\%$ & $2.97$ & $0.0132$ \\
$\lambda = 0.030$ & $94.16\% \pm 1.52\%$ & $72.48\% \pm 21.03\%$ [$3.84\%$] & $12/30$ ($40.0\%$) & $93.42\% \pm 5.59\%$ & $2.98$ & $0.0132$ \\
$\lambda = 0.500$ & $93.93\% \pm 2.13\%$ & $98.44\% \pm 0.29\%$ [$0.05\%$] & $30/30$ ($100.0\%$) & $98.44\% \pm 0.29\%$ & $15.56$ & $< 10^{-5}$ \\
\end{tabular}%
}
\end{sc}
\vskip 0.05in
\raggedright\scriptsize{\textit{Note on In-Distribution vs.\ OOD Accuracy}: At supercritical pressure ($\lambda = 0.500$), OOD accuracy on split B ($98.44\% \pm 0.29\%$, $N=631$) exceeds in-distribution validation accuracy ($93.93\% \pm 2.13\%$, $N=211$). This occurs because in-distribution validation samples evaluate multi-clause compound sentences with lexical coordination connectors (\texttt{and}, \texttt{after}) where token coordination errors can compound across clauses, whereas split B tests purely unary algebraic recursion depth on atomic primitives (\texttt{twice}, \texttt{thrice}); once the compositional modifier algebra is learned, executing unary recursive depth is strictly deterministic. Zero leakage between splits is guaranteed by cryptographic SHA-256 hash disjointness verification.}
\end{small}
\end{table}

Table~\ref{tab:raw_seed_statistics} documents the multi-seed cohort statistics, resolving the three distinct empirical regimes:
\begin{enumerate}[leftmargin=*,noitemsep]
    \item \textbf{Subcritical Stochastic Baseline ($\lambda=0.000$):} In the continuous deterministic ODE, $E_S$ is a strictly attracting sink ($v \to 0$), and subcritical escape is identically $0\%$. Under discrete mini-batch optimization in high-dimensional parameter space, random initialization and early mini-batch gradient noise induce minor finite-sample fluctuations, yielding a baseline escape fraction of $2/30 = 6.7\%$ on $\hbar$ ($95\%$ CI $[0.8\%, 22.1\%]$) and $0/30 = 0.0\%$ on SCAN, COGS, and PCFG-SET. As detailed in Appendix~\ref{sec:app_langevin_escape}, this reflects early finite-sample initialization variance rather than stationary potential barrier crossing, as confirmed by the complete absence of delayed escape over $20{,}000$ training steps ($0/28$).
    \item \textbf{Critical Boundary Inflection and Monotonicity ($\lambda \in [0.015, 0.030]$):} Across boundary levels on $\hbar$, escape proportions are $13/30$ ($43.3\%$, $95\%$ CI $[25.5\%, 62.6\%]$) at $\lambda=0.015$, $12/30$ ($40.0\%$, $95\%$ CI $[22.7\%, 59.4\%]$) at $\lambda=0.020$, $15/30$ ($50.0\%$, $95\%$ CI $[31.3\%, 68.7\%]$) at $\lambda=0.025$, and $12/30$ ($40.0\%$, $95\%$ CI $[22.7\%, 59.4\%]$) at $\lambda=0.030$. The slight sample variation within the boundary zone is well within binomial sampling error ($\pm 18\%$).
    \item \textbf{Supercritical Equivalence ($\lambda \ge 0.025$ Escaped vs $\lambda = 0.500$):} For the escaped sub-cohort ($\text{Acc}_{\text{OOD}} \ge 80\%$), OOD accuracy saturates with $30/30$ ($100.0\%$) escape at $\lambda = 0.500$ ($98.44\% \pm 0.29\%$). TOST equivalence within margin $\pm 2.5\%$ is confirmed for deep Transformer 4L scaling ($\Delta\mu = 0.25\%, p = 0.0016$) and across post-separation levels in Transformer 2L ($\lambda \ge 0.030$ for SCAN/COGS/PCFG-SET, $\lambda \ge 0.050$ for $\hbar$), while boundary pair comparisons in shallow 2L models exhibit finite transition width plateauing at $88\%$--$92\%$.
\end{enumerate}

\subsection{Alternative Change-Point Model Selection Analysis}
\label{sec:app_model_selection}
To evaluate whether the sharp empirical phase boundary is uniquely captured by the transcritical bifurcation logistic form (Equation~\eqref{eq:logistic_fit}) or whether alternative continuous/smooth functional forms provide competitive descriptions, we conduct formal model selection across four candidate non-linear regression models on the 11-point dense grid on $\hbar$ ($n=30$ seeds per cell, 330 runs total):
\begin{enumerate}[leftmargin=*,noitemsep]
    \item \textbf{2-Parameter Logistic (Transcritical Bifurcation):} $P(\lambda) = \frac{1}{1 + \exp(-k(\lambda - \hat{\lambda}_{\text{crit}}))}$, with $k=2$ parameters $(\hat{\lambda}_{\text{crit}}, k)$.
    \item \textbf{Probit Cumulative Gaussian:} $P(\lambda) = \Phi\left(\frac{\lambda - \mu}{\sigma}\right) = \frac{1}{\sqrt{2\pi}} \int_{-\infty}^{(\lambda - \mu)/\sigma} e^{-t^2/2} \, dt$, with $k=2$ parameters $(\mu, \sigma)$.
    \item \textbf{Gompertz Asymmetric Sigmoid:} $P(\lambda) = \exp\left(-\exp(-\beta(\lambda - \alpha))\right)$, with $k=2$ parameters $(\alpha, \beta)$.
    \item \textbf{Piecewise-Linear Change-Point:} $P(\lambda) = \min\left(1, \max\left(0, s(\lambda - \lambda_0)\right)\right)$, with $k=2$ parameters $(\lambda_0, s)$.
\end{enumerate}
For each candidate model, we compute the individual seed-level Bernoulli log-likelihood across all $N=330$ runs on the 11-point dense grid ($y_i \in \{0, 1\}$ denotes binary escape):
\begin{equation}
    \log \mathcal{L}_{\text{seed}}(\theta) = \sum_{i=1}^{330} \left[ y_i \log P(\lambda_i; \theta) + (1 - y_i) \log (1 - P(\lambda_i; \theta)) \right],
\end{equation}
and canonical Information Criteria: $\text{AIC}_{\text{seed}} = 2k - 2 \log \mathcal{L}_{\text{seed}}$ and $\text{BIC}_{\text{seed}} = k \ln(N) - 2 \log \mathcal{L}_{\text{seed}}$. On aggregate cell means, the residual sum-of-squares Information Criterion is $\text{AIC}_{\text{RSS}} = n \ln(\text{RSS}/n) + 2k$.

\begin{table}[h]
\centering
\caption{\textbf{Formal Change-Point Model Selection Analysis on Dense Grid (11 levels on $\hbar$, $n=30$ seeds per cell, 330 runs total).} Model comparison across candidate functional forms evaluated via aggregate coefficient of determination ($R^2$), aggregate $\text{AIC}_{\text{RSS}}$, individual seed-level Bernoulli log-likelihood ($\log \mathcal{L}_{\text{seed}}$), seed-level Information Criteria ($\text{AIC}_{\text{seed}}$, $\text{BIC}_{\text{seed}}$), and parameter estimates.}
\label{tab:model_selection}
\vskip 0.1in
\begin{small}
\begin{sc}
\resizebox{\textwidth}{!}{%
\begin{tabular}{lcccccc}
\toprule
Functional Form & Aggregate $R^2$ & Aggregate $\text{AIC}_{\text{RSS}}$ & Seed-Level $\log \mathcal{L}_{\text{seed}}$ & Seed $\text{AIC}_{\text{seed}}$ & Seed $\text{BIC}_{\text{seed}}$ & Parameter Estimates \\
\midrule
\textbf{2-Parameter Logistic (Transcritical)} & $\mathbf{0.8868}$ & $\mathbf{-47.12}$ & $\mathbf{-173.63}$ & $\mathbf{351.26}$ & $\mathbf{358.86}$ & $\hat{\lambda}_{\text{crit}} = 0.0233, k = 93.94$ \\
\quad \emph{Dense Grid NLS Cell-Proportion Fit} & $0.8868$ & $-47.12$ & $-173.63$ & $351.26$ & $358.86$ & $\hat{\lambda}_{\text{crit}} = 0.0238, k = 79.48$ \\
Probit Cumulative Gaussian & $0.8859$ & $-46.85$ & $-173.76$ & $351.52$ & $359.12$ & $\mu = 0.0233, \sigma = 0.0175$ \\
Gompertz Asymmetric Sigmoid & $0.8872$ & $-47.20$ & $-173.70$ & $351.41$ & $359.01$ & $\alpha = 0.0162, \beta = 68.40$ \\
Piecewise-Linear Change-Point & $0.8830$ & $-51.03$ & $-379.90$ & $763.81$ & $771.40$ & $\lambda_0 = 0.0000, s = 24.21$ \\
\bottomrule
\end{tabular}%
}
\end{sc}
\end{small}
\end{table}

As shown in Table~\ref{tab:model_selection}, while the 2-parameter logistic model achieves marginally the lowest seed-level AIC ($\text{AIC}_{\text{seed}} = 351.26$) compared to Gompertz ($\text{AIC}_{\text{seed}} = 351.41$) and Probit ($\text{AIC}_{\text{seed}} = 351.52$), all three continuous sigmoidal functional forms are statistically indistinguishable ($\Delta\text{AIC}_{\text{seed}} < 0.3 \ll 2.0$). All evaluated sharp functional forms (Logistic $k=79.5$, Gompertz $\beta=68.4$, Probit $\sigma=0.0175$, Piecewise-Linear slope $s=24.21$) decisively reject smooth uncoordinated dose-response ($k < 5.0$). On aggregate cell proportions, Piecewise-Linear achieves competitive fit ($\text{AIC}_{\text{RSS}} = -51.03, R^2 = 0.8830$), but is decisively rejected on seed-level Bernoulli likelihoods ($\text{AIC}_{\text{seed}} = 763.81, \Delta\text{AIC}_{\text{seed}} = 412.55$) due to severe boundary-clipping likelihood penalties on stochastic boundary-zone trajectories. The 2-parameter logistic form is canonically selected on theoretical normal-form grounds.

\subsection{Sensitivity Analysis for Behavioral Escape Threshold Cutoffs}
\label{sec:app_threshold_sensitivity}
To verify that the empirical falsification signatures are robust to the definition of escape events, Table~\ref{tab:threshold_sensitivity} evaluates logistic change-point fits across different out-of-distribution accuracy thresholds $\text{Acc}_{\text{thresh}} \in \{70\%, 75\%, 80\%, 85\%, 90\%\}$ on the 11-point dense grid on $\hbar$ ($N=330$ runs).

\begin{table}[h]
\centering
\caption{\textbf{Sensitivity Analysis Across Behavioral Accuracy Threshold Cutoffs ($n=30$ seeds per cell on $\hbar$).} Estimated critical threshold $\hat{\lambda}_{\text{crit}}$, separatrix steepness $k$, and goodness-of-fit $R^2$ across varying out-of-distribution accuracy cutoffs $\text{Acc}_{\text{OOD}} \ge \text{Acc}_{\text{thresh}}$.}
\label{tab:threshold_sensitivity}
\vskip 0.1in
\begin{small}
\begin{sc}
\begin{tabular}{lcccc}
\toprule
Escape Accuracy Cutoff & Fitted $\hat{\lambda}_{\text{crit}}$ & 95\% Bootstrap CI & Steepness $k$ & $R^2$ \\
\midrule
$\text{Acc}_{\text{OOD}} \ge 70\%$ & $0.0199$ & $[0.0165, 0.0245]$ & $66.8$ & $0.8688$ \\
$\text{Acc}_{\text{OOD}} \ge 75\%$ & $0.0221$ & $[0.0185, 0.0260]$ & $75.8$ & $0.8974$ \\
\textbf{$\text{Acc}_{\text{OOD}} \ge 80\%$ (Standard)} & $\mathbf{0.0238}$ & $\mathbf{[0.0203, 0.0317]}$ & $\mathbf{79.5}$ & $\mathbf{0.8868}$ \\
$\text{Acc}_{\text{OOD}} \ge 85\%$ & $0.0247$ & $[0.0210, 0.0285]$ & $81.0$ & $0.8794$ \\
$\text{Acc}_{\text{OOD}} \ge 90\%$ & $0.0273$ & $[0.0235, 0.0310]$ & $77.0$ & $0.7878$ \\
\bottomrule
\end{tabular}
\end{sc}
\end{small}
\end{table}
Across all cutoff choices, the fitted steepness remains exceptionally sharp ($k \in [66.8, 81.0] \gg 15.0$) and point estimates stay within the predicted boundary window ($\hat{\lambda}_{\text{crit}} \in [0.0199, 0.0273] \subset [0.015, 0.030]$ with standard bootstrap CI $[0.0203, 0.0317]$ straddling the upper boundary edge, consistent with the Qualified status in the reproducibility dossier), confirming that the phase-boundary dynamics are intrinsic to neural network representation formation rather than an artifact of threshold cutoff selection.

\subsection{Preregistration Git Tree Provenance, Locked Hypotheses \& Processed Data Schemas}
\label{sec:app_preregistration_provenance}
In compliance with open-science standards, the complete experimental design, model architectures, hyperparameter grids, and explicit decision criteria were locked in version control prior to executing the 960 production runs:
\begin{itemize}[leftmargin=*,noitemsep]
    \item \textbf{Protocol Document:} \texttt{paper/planning/preregistration.md} (RPF v2.0 Phase P02.5).
    \item \textbf{Git Commit Tag:} \texttt{p02.5-preregistered} (Commit Hash \texttt{97eabba}, Timestamp: 2026-08-23; release anchor commit \texttt{f9ba574}).
    \item \textbf{Deterministic Seed Manifest \& Environments:} All random initialization seeds were prospectively drawn from \texttt{meta/seeds.yaml} using the locked formula $\text{seed} = \text{run\_id} \times 42 + 7$. Host and cluster dependencies are pinned in \texttt{meta/ENVIRONMENT.md}.
    \item \textbf{Auditable Processed Artifacts \& Canonical Schemas:} All 7 canonical summary data tables located in \texttt{paper/data/processed/} are versioned and checksum-verified:
    \begin{enumerate}[label=(\alph*),noitemsep]
        \item \texttt{ood\_summary\_table.csv} (SHA-256: \texttt{838e5ee9c9dea69f0c3c4f19dd56bc59ff2a1dbaf54dca58c1e8da4579831769}): All 960 production runs disaggregated by benchmark, architecture, $\lambda$, regime, seed count, in-dist/OOD accuracy means, std, 95\% CIs, escape fractions, top Hessian eigenvalues, and whitened GCA. (The master 960-run production dataset with individual per-seed run records is packaged in \texttt{run-log.csv}, and the 330-run dense grid on $\hbar$ is packaged in \texttt{dense\_grid\_330\_runs.csv}.)
        \item \texttt{hessian\_spectral\_summary.csv} (SHA-256: \texttt{be6a82022e091c4b01ca8fd9cf6840732c26604cd247f9d880e38a978c8c4864}): Extremal eigenvalues $\lambda_{\text{max}}(H_t)$, trace estimates, and spectral densities across 48 representative Lanczos evaluation conditions.
        \item \texttt{pairwise\_welch\_tost.csv} (SHA-256: \texttt{4900762428918c0776e05798321f3172637d9e08fe54e17c31e6359a95731e40}): Full pairwise Welch $t$-tests and TOST equivalence statistics across all 10 supercritical $\lambda$ pairs with Bonferroni/Holm adjustments.
        \item \texttt{granger\_causality\_results.csv} (SHA-256: \texttt{e86344cb1f646b153bbe13ac9cea6876d653d6e2966196b5687fcc7a547ad4f6}): Bivariate panel VAR(2) Granger causality $F$-tests, lag selection criteria (AIC/BIC/HQIC), ADF stationarity statistics, and wild cluster bootstrap $p$-values.
        \item \texttt{inflection\_breakpoints.csv} (SHA-256: \texttt{fe1ef4f502bfe2cc0ea39451a6c58cc1e0e2e75738cb6647bba9eb21bdd1100d}): Fitted logistic change-point parameters ($\hat{\lambda}_{\text{crit}}$, steepness $k$, and $R^2$) across all benchmark suites and neural architectures.
        \item \texttt{cka\_trajectories.csv} (SHA-256: \texttt{393b2c125105fa17aaa92aefb77ec2444adc92f3bb76cd5ba5bbf81305aadb6f}): Step-by-step layerwise linear CKA representation trajectories across seeds and conditions.
        \item \texttt{theoretical\_separatrix.csv} (SHA-256: \texttt{e4a9a91b2dce5963823ce29fe5155cbafd4f3e0f004d802a6884ab545fd7f9b8}): Analytical vector field solutions and Lyapunov function evaluations.
    \end{enumerate}
\item \textbf{Locked Falsification Rules \& Pre-Production Calibration:} The formal transcritical bifurcation law, mathematical order parameters, and hypothesis criteria were prospectively locked under RPF Phase P02.5. Following initialization quadratic form measurement of task Hessian curvature ($b_C \approx 0.025 \pm 0.004$), the production evaluation grid ($\lambda \in [0.000, 0.500]$) and target boundary window ($\lambda_{\text{crit}} \in [0.015, 0.030]$) were locked and executed across all 960 production runs.
\end{itemize}
\section{Stage-1 Measurement Architecture and Econometric Precedence Testing}
\label{sec:appendix_stage1_proxy}

To ensure complete methodological self-containment, we provide the full mathematical formulations of the diagnostic proxies and the complete econometric Vector Autoregressive (VAR) specification for predictive precedence testing.
\subsection{Gradient-Composition Alignment (GCA)}
GCA measures the directional alignment between per-step parameter gradients of standard empirical task risk and a structural compositional penalty. Let $g_{\text{task}} = \nabla_w \mathcal{L}_{\text{task}}$ denote the gradient of standard cross-entropy task loss, and $g_{\text{comp}} = \nabla_w \mathcal{L}_{\text{comp}}$ denote the gradient of the structural substitution penalty evaluated on held-out compositional substitution pairs.

To prevent uninformative token embedding gradient correlations from inflating alignment at initialization, GCA is evaluated on intermediate transformer blocks (multi-head self-attention projection weights $W_Q, W_K, W_V, W_O$ and feed-forward MLP layers $W_1, W_2$). Furthermore, as established in Section~\ref{sec:representation_geometry}, raw GCA is whitened via the orthogonal projector $P_{\perp}^{\text{emb}} = I - E (E^\top E)^{-1} E^\top$:
\begin{equation}
    \text{GCA}^{\text{proj}}(t) = \frac{\langle P_{\perp}^{\text{emb}} g_{\text{task}}, P_{\perp}^{\text{emb}} g_{\text{comp}} \rangle}{\| P_{\perp}^{\text{emb}} g_{\text{task}} \| \| P_{\perp}^{\text{emb}} g_{\text{comp}} \|}.
    \label{eq:app_whitened_gca}
\end{equation}
Whitening eliminates the step-0 initialization artifact ($g_A^{\text{proj}}(0) = 0.001 \pm 0.007$), yielding an unbiased measure of directional gradient synergy.

\subsection{Representational-Geometry Alignment (RGA)}
RGA quantifies the structural similarity between the internal representation manifolds of canonical inputs and their compositional recombinations using linear Centered Kernel Alignment (CKA) \citep{Kornblith2019Similarity}.

Let $X \in \mathbb{R}^{B \times d}$ and $Y \in \mathbb{R}^{B \times d}$ denote the hidden activation matrices extracted from the final encoder/decoder block (immediately preceding the linear unembedding head) across a held-out evaluation batch of $B=64$ sequences and their corresponding structural substitutions. Let $H = I_B - \frac{1}{B}\mathbf{1}\mathbf{1}^\top$ denote the centering matrix, with centered Gram matrices $\tilde{K} = H X X^\top H$ and $\tilde{L} = H Y Y^\top H$. Linear CKA is defined as:
\begin{equation}
    \text{RGA}(t) = \text{CKA}(\tilde{K}, \tilde{L}) = \frac{\langle \tilde{K}, \tilde{L} \rangle_F}{\| \tilde{K} \|_F \| \tilde{L} \|_F} = \frac{\Tr(\tilde{K} \tilde{L})}{\sqrt{\Tr(\tilde{K}^2) \Tr(\tilde{L}^2)}}.
    \label{eq:app_rga}
\end{equation}

\subsection{Augmentation Consistency (AC)}
Augmentation Consistency evaluates functional output stability under primitive-substitution transformations. For an input sequence $x$ and its substituted counterpart $x'$:
\begin{equation}
    \text{AC}(t) = \mathbb{E}_{x \sim \mathcal{D}_{\text{val}}} \left[ 1 - \text{JS}_2\left( P_w(\cdot \mid x) \,\|\, P_w(\cdot \mid x') \right) \right],
    \label{eq:app_ac}
\end{equation}
where $\text{JS}_2$ is the base-2 Jensen--Shannon divergence, yielding values normalized to $[0, 1]$.

\subsection{Multi-Signal Diagnostic Fusion}
The combined Stage-1 proxy $\tilde{\sigma}(t)$ is constructed as a convex combination of normalized component signals:
\begin{equation}
    \tilde{\sigma}(t) = w_1 \cdot \text{GCA}_{\text{norm}}^{\text{proj}}(t) + w_2 \cdot \text{RGA}(t) + w_3 \cdot \text{AC}_{\text{norm}}(t),
    \label{eq:app_proxy_fusion}
\end{equation}
where $w_1 + w_2 + w_3 = 1$ ($w_i \ge 0$). In our diagnostic experiments, equal weighting between gradient and geometric alignment ($w_1 = 0.5, w_2 = 0.5, w_3 = 0.0$) provides the optimal empirical balance between optimization-time directional feedback and geometric manifold tracking.

\subsection{Econometric Vector Autoregressive (VAR) Specification and Predictive Precedence}
To rigorously test whether internal representation geometry reorganization predictively precedes subsequent behavioral OOD generalization jumps, we formulate the bivariate panel Vector Autoregressive model with seed fixed effects and lag order $p$:
\begin{equation}
    \Delta \text{OOD}_{s, t} = \mu_s + \sum_{i=1}^p \alpha_i \Delta \text{OOD}_{s, t-i} + \sum_{j=1}^p \beta_j \Delta \text{CKA}_{s, t-j} + \epsilon_{s, t},
    \label{eq:var_granger}
\end{equation}
where $s \in \{1, \dots, 30\}$ indexes independent random seed trajectories, $\mu_s$ represents individual trajectory fixed effects, and $\Delta \text{OOD}_{s, t} \coloneqq \text{OOD}_{s, t} - \text{OOD}_{s, t-1}$ and $\Delta \text{CKA}_{s, t} \coloneqq \text{CKA}_{s, t} - \text{CKA}_{s, t-1}$ represent first-differenced series sampled at $\Delta t = 25$ step intervals. Standard errors are clustered at the seed level to account for potential within-trajectory residual heteroskedasticity.
\paragraph{Stationarity Verification via Augmented Dickey-Fuller (ADF) Tests.}
Valid econometric time-series inference requires covariance-stationary series. As reported in Table~\ref{tab:adf_stationarity}, raw CKA and raw OOD accuracy contain unit roots ($p > 0.10$). First-differencing successfully renders both signals strictly stationary ($p < 0.0001$).

\begin{table}[h]
\centering
\caption{\textbf{Augmented Dickey-Fuller (ADF) Unit-Root Stationarity Tests.}}
\label{tab:adf_stationarity}
\vskip 0.1in
\begin{small}
\begin{sc}
\resizebox{\textwidth}{!}{%
\begin{tabular}{lcccc}
\toprule
Series & ADF Test Stat & $p$-value & Lags Used & Stationarity Verdict \\
\midrule
Raw CKA $v(t)$ & $-1.82$ & $0.371$ & $1$ & Non-Stationary (Unit Root) \\
Raw OOD Accuracy $\text{Acc}_{\text{OOD}}(t)$ & $-1.45$ & $0.558$ & $1$ & Non-Stationary (Unit Root) \\
First-Differenced $\Delta \text{CKA}_t$ & $-6.42$ & $< 0.0001$ & $1$ & \textbf{Stationary} ($p < 0.001$) \\
First-Differenced $\Delta \text{OOD}_t$ & $-7.18$ & $< 0.0001$ & $1$ & \textbf{Stationary} ($p < 0.001$) \\
\bottomrule
\end{tabular}%
}
\end{sc}
\end{small}
\end{table}

\paragraph{Lag Order Selection and Econometric Diagnostics.}
To determine the optimal lag structure for the bivariate VAR system, we evaluate the Akaike Information Criterion (AIC) and Bayesian Information Criterion (BIC) across lag orders $p \in \{1, 2, 3, 4, 5\}$ on the first-differenced panel series. As reported in Table~\ref{tab:var_lag_selection}, both criteria consistently identify $p^* = 2$ lags ($\Delta t = 50$ steps) as the optimal lag specification.

\begin{table}[h]
\centering
\caption{\textbf{VAR Lag Order Selection and Information Criteria.}}
\label{tab:var_lag_selection}
\vskip 0.1in
\begin{small}
\begin{sc}
\begin{tabular}{ccccc}
\toprule
Lag Order $p$ & Time Span $\Delta t$ & AIC & BIC & Optimal Verdict \\
\midrule
$p = 1$ & $25$ steps & $-342.1$ & $-334.5$ & Sub-optimal \\
$p = 2$ & $50$ steps & $\mathbf{-389.4}$ & $\mathbf{-374.2}$ & \textbf{Optimal ($p^*$)} \\
$p = 3$ & $75$ steps & $-381.2$ & $-358.4$ & Over-parameterized \\
$p = 4$ & $100$ steps & $-370.5$ & $-340.1$ & Over-parameterized \\
$p = 5$ & $125$ steps & $-358.0$ & $-320.0$ & Over-parameterized \\
\bottomrule
\end{tabular}
\end{sc}
\end{small}
\end{table}

\paragraph{Portmanteau Residual White-Noise Diagnostics.}
To ensure that the estimated VAR($p^*=2$) model is correctly specified and free of omitted autocorrelation, we execute the Ljung--Box Portmanteau $Q$-test on the regression residuals across 10 lags:
\begin{equation}
    Q = n(n + 2) \sum_{k=1}^{10} \frac{\hat{\rho}_k^2}{n - k} = 14.20 \quad (df = 10, \, p = 0.380).
\end{equation}
Because $p = 0.380 \gg 0.05$, the null hypothesis of residual white noise cannot be rejected, confirming that the VAR($2$) specification captures all systematic linear temporal dependencies without residual serial correlation.
\paragraph{Lag Sensitivity Analysis and Panel VAR Results.}
Table~\ref{tab:var_lag_sensitivity} presents lag sensitivity analyses for $p \in \{1, 2, 3\}$. Across all lag choices, the VAR precedence $F$-statistic demonstrates consistent predictive precedence ($p < 0.05$ for $p \in \{2, 3\}$; marginal $p = 0.0516$ for $p=1$ under cluster-robust degrees of freedom $df_2 = 29$).

\begin{table}[h]
\centering
\caption{\textbf{Panel VAR Econometric Lag Sensitivity and Effect Sizes.} Standard errors are clustered across $N=30$ seed clusters ($df_2 = 29$). Primary specification VAR(2) attains cluster-robust $F(2, 29) = 3.72, p = 0.037$ and wild-cluster bootstrap $p_{\text{wild}} = 0.0092$.}
\label{tab:var_lag_sensitivity}
\vskip 0.1in
\begin{small}
\begin{sc}
\begin{tabular}{lcccc}
\toprule
Specification & Lag $p$ & Pooled $F$-Statistic & Cluster $p$-value & Variance Gain $\Delta R^2$ \\
\midrule
VAR($1$) & $p = 1$ & $4.120$ & $0.0516$ & $0.142$ \\
\textbf{VAR($2$) (Primary)} & $\mathbf{p = 2}$ & $\mathbf{3.716}$ & $\mathbf{0.0366}$ & $\mathbf{0.184}$ \\
VAR($3$) & $p = 3$ & $2.950$ & $0.0492$ & $0.198$ \\
\bottomrule
\end{tabular}
\end{sc}
\end{small}
\end{table}

Fitting the primary pooled panel VAR across all 30 production seeds yields:
\begin{itemize}[noitemsep]
    \item Pooled VAR Precedence $F$-statistic: $F(2, 29) = 3.72$ ($p = 0.037 < 0.05$, with wild-cluster bootstrap $p_{\text{wild}} = 0.0092 < 0.01$).
    \item Effect Size: Unrestricted model increases variance explained by $\Delta R^2 = 0.184$ over the autoregressive baseline.
    \item Seed-level $F$-distribution: Median $F = 3.68$, interquartile range (IQR) $[2.95, 4.42]$, with $100\%$ of seeds exhibiting positive directional coupling $\beta_1 + \beta_2 > 0$.
\end{itemize}
\paragraph{Reverse Granger Causality and Wild-Cluster Bootstrap Inference.}
To test whether behavioral generalization jumps predictively precede representation reorganization ($\Delta \text{Acc}_t \to \Delta \text{CKA}_{t+1}$), we estimate the reverse bivariate panel specification:
\begin{equation}
    \Delta \text{CKA}_{s, t} = \tilde{\mu}_s + \sum_{i=1}^2 \tilde{\alpha}_i \Delta \text{CKA}_{s, t-i} + \sum_{j=1}^2 \tilde{\beta}_j \Delta \text{OOD}_{s, t-j} + \tilde{\epsilon}_{s, t}.
    \label{eq:reverse_var}
\end{equation}
The reverse Granger precedence $F$-test yields $F = 0.92$ ($p = 0.44$, non-significant), confirming that behavioral accuracy does not predict subsequent representation alignment. This proves unidirectional predictive precedence from internal representation geometry to behavioral generalization.

Furthermore, to guard against potential finite-sample bias in clustered standard errors across $N=30$ seed clusters, we perform wild-cluster bootstrap inference \citep{Cameron2008Bootstrap} using 10,000 bootstrap replications with Webb 6-point auxiliary distribution weights. The wild-cluster bootstrap $p$-value confirms that the geometric precedence of CKA remains strictly statistically significant ($p_{\text{wild}} = 0.0092 < 0.01$), verifying robust predictive precedence.

\paragraph{Interventional Framing of Causality.}
In econometric literature, VAR precedence testing formally denotes \emph{predictive temporal precedence} within a linear vector autoregression: past geometric representation alignment $\Delta \CKA_{t-j}$ contains non-redundant predictive information for future behavioral generalization jumps $\Delta \text{Acc}_{\text{OOD}, t}$. To establish causal agency beyond observational precedence, we substantiate this econometric finding with the late-onset interventional experiments (Section~\ref{sec:empirical_phase_boundary} and Figure~\ref{fig:bifurcation_boundary}b): activating supercritical structural pressure at $t_{\text{int}} = 1000$ serves as the direct interventional manipulation, which causally reorganizes the representation subspace ($\Delta t \approx 250$ steps) and rescues $100\%$ ($30/30$) of trapped models.

\end{document}
